% Catégories de variétés (Préliminaires au livre des variétés, n° 262)
% Source scans: U46.pdf (CSG, csg.igrothendieck.org)
% Model: gemini-3.1-pro-preview (thinking=medium)
% Markers written by the pipeline from the PDF page index
% Pages transcribed: 100/100
% Rebuilt: 2026-08-14

%% ===== Page 1 =====

% n° 262
% [marge: Archives Grothendieck. Avril 57.]

PRELIMINAIRES AU LIVRE DES VARIETES

CATEGORIES DE VARIETES

\section*{Paragraphe : 1. CATEGORIES TOPOLOGIQUES.}
1. Catégories prétopologiques.
2. Un procédé de recollement.
3. Produits, catégories topologiques.
4. Produits de $\mathcal{C}$-espaces locaux, Catégories topologiques recollables
6. Structures induites.
7. Structures quotient.
8. Espaces réduits à un point.
9. Espaces fibrés de type $(\mathcal{B}, \underline{\Phi}, \underline{\Phi}_0)$.
10. $\mathcal{C}$-espaces fibrés à $\mathcal{C}$-groupe structural.
11. $\mathcal{C}$-fibrés associés
11. bis - Extension et restriction du groupes structural
12. Sections de $\mathcal{C}$-fibrés.
13. Images réciproques de $\mathcal{C}$-fibrés.
14. Catégories de modèles.
15. Catégorie topologique définie par une catégorie de modèles.
16. Foncteurs topologiques.

\section*{Paragraphe : 2 CATEGORIES DE VARIETES.}
1. Espaces $K$-annelés.
2; Catégories de variétés sur $K$.
3. Structures induites.
4. Structures quotient.
5. Un axiome supplémentaire.
6. $\mathcal{V}$-Structures sur les espaces vectoriels.
7. $\mathcal{V}$-fibrés vectoriels, Relations avec les faisceaux de modules.
8. Opérations sur les fibrés vectoriels.
9. Sous-fibrés vectoriels, fibrés vectoriels quotients, $\mathcal{V}$-fibrés en drapeaux.
10. Extensions de fibrés vectoriels.

\newpage

%% ===== Page 2 =====

% n° 262

\emph{Commentaires du rédacteur.}

Ce rapport est destiné à esquisser les fondements fonctoriels pour
le Livre des variétés, mais il semble hors de doute que sa place véri-
table est dans le livre de Top. Elém..

Une grande partie de ce rapport sera bouffée par le rapport ultérieur
de Cartan-Sammy, néanmoins il était indispensable au rédacteur, pour
pouvoir commencer la rédaction des variétés,. On notera que le paragra-
phe 1, comprend en réalité trois parties assez distinctes, les n$^{\text{os}}$ 1 à
8 donnent les généralités sur les catégories topologiques, : les n$^{\text{os}}$ 9
à 13 les notions essentielles sur les fibrés, enfin, les n$^{\text{os}}$ 14,15 16
introduisent les catégories de modèles, qui ne sont que des intermédi-
aires techniques pour construire, par recollement, des vraies catégories
topologiques. (Faire attention qu'il n'y a pas de n$^\circ$ 5, mais en revan-
che un n$^\circ$ 11 bis).

Les notions introduites dans le paragraphe 1 sont utilisées aussi-
tôt dans le paragraphe 2, qui en est une bonne illustration,. Un para-
graphe 3 devait donner la définition des catégories usuelles de varié-
tés (en supposant connu seulement quelques résultats élémentaires de
dépotoirs divers) et les relations entre ces catégories de variétés,
mais le rédacteur a fini par se dégonfler. Bien que très probablement
un tel paragraphe serait finalement démembré, il semble indiqué qu'il
soit rédigé prochainement (Chevalley ?), pour qu'on se rende compte
si la présente rédaction donne tout ce qu'il faut.

Il va sans dire que le rédacteur n'a pas fait le moindre effort en
direction d'un style canonique.

\newpage

%% ===== Page 3 =====

% n° 262

\section*{\S 1. - \emph{Catégories topologiques.}}

1. \emph{Catégories prétopologiques.}

Soit $\mathcal{C}$ une catégorie, muni d'un foncteur $i_{\mathcal{C}}$ à valeurs dans la
catégorie $\mathcal{T}$ des espaces topologiques. On dit que $(\mathcal{C}, i_{\mathcal{C}})$ est une
\emph{catégorie prétopologique}, et que $i_{\mathcal{C}}$ est son \emph{foncteur fondamental}, si
les axiomes CT 1 à CT 4 qui suivent sont satisfaits.

CT 1) $i$ \emph{est fidèle} :

Rappelons (Ens. 2ème Ed., Chap.4) que cela signifie que, si $f, g :$
$A \to B$ sont deux $\mathcal{C}$-morphismes tels que $i(f) = i(g)$, alors $f = g$,
et que si $f$ est un isomorphisme dans $\mathcal{C}$, tel que $i(f)$ soit l'applica-
tion identique de $i(A)$ sur $i(B) = i(A)$, alors $A = B$ et $f$ est le morphis-
me identique. Pour tout $X \in \mathcal{T}$, soit $\mathcal{S}(X)$ la classe des objets $A$ de $\mathcal{C}$
tels que $i(A) = X$, on identifie alors, conformément aux conventions
générales relatives aux foncteurs fidèles, un $A \in \mathcal{S}(X)$ à l'espace $X$,
muni de la structure supplémentaire $A$, et on mettra le plus souvent
$A$ dans la notation, suivant un abus de notation couramment employé dans
tous les livres de notre traité. Les morphismes d'un $X \in \mathcal{C}$ dans un $Y \in \mathcal{C}$
sont donc certaines applications continues de l'espace $X$ dans l'espace
$Y$, appelées $\mathcal{C}$-\emph{applications}.

CT 2) $i$ \emph{est transportable}.

Rappelons (Ens. 2ème Ed. chap. 4) que cela signifie que si $A \in \mathcal{C}$ ,
$Y \in \mathcal{T}$ et si $f'$ est un isomorphisme (i.e. un homéomorphisme) de $i(A) = X$
sur $Y$, alors il existe un morphisme $f$ dans $\mathcal{C}$ tel que $i(f) = f'$. Il résul-
te alors de CT 1, que cet $f$ est unique, et appliquant l'axiome CT 2, à
$f^{\prime -1}$, on voit que $f$ est un isomorphisme. On voit donc qu'un homéomor-
phisme $f$ d'un espace $X$ sur un espace $Y$ définit une bijection

\newpage

%% ===== Page 4 =====

% typescript page 2 — n° 262

$\bar{f}$ de la classe $\mathcal{C}(X)$ des $\mathcal{C}$-structures sur $X$, sur la classe $\mathcal{C}(Y)$ des
$\mathcal{C}$-structures sur $Y$. Pour $f$ variable, $\bar{f}$ a les propriétés fonctorielles
usuelles. Dans tous les cas de foncteurs transportables que nous aurons
à considérer, les $\mathcal{C}(X)$ sont des ensembles (il n'y aurait donc guère
d'inconvénient à mettre ceci comme condition dans la définition d'un
foncteur transportable), et $X \to \mathcal{C}(X)$ est donc un foncteur sur la caté-
gorie restreinte de $\mathcal{C}$ (i.e. où on restreint les morphismes à être des
isomorphismes) à valeurs dans la catégorie $\mathcal{E}$ des espaces topologiques,
c'est donc une "espèce de structure" (Ens. 2è ed. chap.IV) - Les objets
de $\mathcal{C}$ seront aussi appelés des \emph{$\mathcal{C}$-espaces}.

Rappelons enfin, dans le contexte des foncteurs transportables la
notion de structure initiale et de structure finale ; choisissons
comme un foncteur transportable à valeurs dans la catégorie des ensembles
Soit $X$ un ensemble, soit $(f_i)$ une famille de morphismes, i.e. d'applica-
tions, de $X$ dans des $X_i \in \mathcal{C}$, on dit qu'une \emph{$\mathcal{C}$-structure $S$ sur $X$ est
structure initiale} pour la famille $(f_i)$, si, quel que soit $Y \in \mathcal{C}$, et
l'application $f$ de $Y$ dans $X$, $f$ est un $\mathcal{C}$-morphisme si et seulement si les
$f_i f$ le sont. Cette structure $S$ est alors manifestement unique. En par-
ticulier, on a la notion de \emph{$\mathcal{C}$-structure induite} sur une partie $U$ d'un
$X \in \mathcal{C}$, qui est uniquement déterminée si elle existe. La définition des
\emph{structures finales}, et en particulier des \emph{$\mathcal{C}$-structures quotient}, est
duale.

CT 3) \emph{Soit $X \in \mathcal{C}$, $U$ une partie ouverte de $X$. Alors $U$ admet une $\mathcal{C}$-struc-
ture induite, et la topologie correspondante de $U$ est celle induite
par $X$}

CT 4) \emph{Soient $X, Y \in \mathcal{C}$, $(U_i)$ un recouvrement ouvert de $X$, $f$ une application}

\newpage

%% ===== Page 5 =====

% typescript page 3 — n° 262

continue de $X$ dans $Y$, dont la restriction à tout $U_i$ est un $\mathcal{C}$-\emph{mor}-
phisme ; alors $f$ est un $\mathcal{C}$-\emph{morphisme}.

On déduit immédiatement de CT 4 :

PROPOSITION 1. - Soient $X$ un espace topologique, $(U_i)$ un recouvrement
ouvert de $X$, $S$ et $S'$ deux $\mathcal{C}$-structures sur $X$ qui induisent sur chaque
$U_i$ des $\mathcal{C}$-structures identiques. Alors $S = S'$.

Corollaire. - Soient $f$ une application d'un $\mathcal{C}$-espace $X$ dans un autre
$Y$, $(U_i)$ un recouvrement ouvert de $X$. Pour que $f$ soit un $\mathcal{C}$-isomorphisme de
$X$ sur un ouvert de $Y$, il faut et il suffit que $f$ soit un monomorphisme
et que sa restriction à tout $U_i$ soit un $\mathcal{C}$-isomorphisme de $U_i$ sur une
partie ouverte de $Y$.

\emph{Sous-catégories prétopologiques}. - Soient $\mathcal{C}$, $\mathcal{C}'$ deux catégories pré-
topologiques. On dit que $\mathcal{C}'$ est une sous-catégorie prétopologique de $\mathcal{C}$
si $\mathcal{C}' \subset \mathcal{C}$ , si le foncteur $i_{\mathcal{C}'}$ est induit par $i_{\mathcal{C}}$ et si pour
$X, Y \in \mathcal{C}'$, les $\mathcal{C}'$-morphismes de $X$ dans $Y$ coincident avec les $\mathcal{C}$-morphismes
de $X$ dans $Y$, et enfin si pour tout ouvert $U$ d'un $X \in \mathcal{C}'$, la $\mathcal{C}$-structure
induite est dans $\mathcal{C}'$ (donc identique à la $\mathcal{C}'$-structure induite). Il
s'ensuit que $\mathcal{C}'$ est complètement connue quand on connait la \emph{classe}
$|\mathcal{C}'|$ de ses éléments. Une sous-classe $\mathcal{C}'$ de $\mathcal{C}$ correspond à une sous-
catégorie prétopologique si et seulement si elle satisfait les conditions
suivantes :

(i) Un $X \in \mathcal{C}$ qui est $\mathcal{C}$-isomorphe à un $Y \in \mathcal{C}'$, est dans $\mathcal{C}'$.

(ii) Si $X \in \mathcal{C}'$, alors toute partie ouverte $U$ de $X$ munie de la $\mathcal{C}$-
structure induite, est dans $\mathcal{C}'$.

2. - \emph{Un procédé de recollement}.

Nous dirons qu'un espace $X$ est muni d'une $\mathcal{C}$-structure locale, ou
est un $\mathcal{C}$-\emph{espace local}, si on s'est donné un recouvrement ouvert $(U_i)$

\newpage

%% ===== Page 6 =====

% typescript page 4 — n° 262

de $X$ et des $\mathcal{C}$-structures $S_i$ sur les $U_i$, telle façon que pour tout couple
$(i,j)$ d'indices tels que $U_i \cap U_j \neq \emptyset$, on ait $S_i | U_i \cap U_j = S_j | U_i \cap U_j$ ; étant
entendu que deux telles données $(U_i,S_i)$ et $(U'_j,S'_j)$ sont considérées comme
définissant la même structure de $\mathcal{C}$-espace local, si et seulement si
$U_i \cap U'_j \neq \emptyset$ implique que $S_i$ et $S'_j$ induisent la même $\mathcal{C}$-structure sur
$U_i \cap U'_j$. (C'est bien là une relation d'équivalence dans l'ensemble des sys-
tèmes $(U_i,S_i)$ grâce à CT 4). On désigne par $\hat{\mathcal{C}}$ la classe des $\mathcal{C}$-espaces
locaux,. Tout $\mathcal{C}$-espace est un $\hat{\mathcal{C}}$-espace local (prendre le recouvrement
de $X$ réduit à un seul élément $X$), et deux $\mathcal{C}$-structures sur un espace $X$
sont identiques si et seulement si les $\hat{\mathcal{C}}$-structures locales correspon-
dantes sont identiques, donc on a $\mathcal{C} \subset \hat{\mathcal{C}}$. Soit $X$ un $\hat{\mathcal{C}}$-espace local
défini par un système $(U_i,S_i)$, $U$ une partie ouverte de $X$, alors les $(U \cap U_i,$
$S_i | U \cap U_i)$ définissent sur $U$ une $\hat{\mathcal{C}}$-structure locale qui ne dépend que de
celle de $X$, et appelée $\hat{\mathcal{C}}$-\emph{structure locale induite}. Si $X \in \mathcal{C}$, on retrouve
bien la $\mathcal{C}$-structure induite. D'après la transitivité (évidente) des struc-
tures induites, on voit que tout point d'un $\hat{\mathcal{C}}$-espace local admet un sys-
tème fondamental de voisinages ouverts qui sont dans $\mathcal{C}$ (pour la structure
induite). On dira qu'une application $f$ d'un $\hat{\mathcal{C}}$-espace local $X$ dans un
autre est une $\hat{\mathcal{C}}$-\emph{application}, si pour tout $x \in X$ on peut trouver un
voisinage ouvert $U$ de $x$ et une partie ouverte $V$ de $Y$, telles que $U$ et $V$
soient des $\mathcal{C}$-espaces, $f(U) \subset V$, et que $f$ soit un $\mathcal{C}$-morphisme de $U$ dans
$V$. D'après CT 4), si $X,Y \in \mathcal{C}$ , on retrouve la notion donnée de $\mathcal{C}$-appli-
cation. Ainsi $\hat{\mathcal{C}}$ devient une catégorie, de plus le foncteur naturel $i$
de $\mathcal{C}$ dans $\hat{\mathcal{C}}$ est trivialement représentable et fidèle, et la notion de
structure induite $S_i$ sur un ouvert $U \subset X$ introduite plus haut est bien
celle qui se déduit de la donnée de $i$ et des morphismes dans $\hat{\mathcal{C}}$,

\newpage

%% ===== Page 7 =====

% typescript page 5 — n° 262

Cela signifie en effet que si Y est un $\mathcal{C}$-espace local et f une applica-
tion continue de Y dans U, alors f est un $\mathcal{C}$-morphisme de Y dans U si et
seulement si c'est un $\mathcal{C}$-morphisme de Y dans X, ce qui résulte facilement
des définitions. Enfin, la vérification de l'axiome \underline{CT 4)} pour $\hat{\mathcal{C}}$ est
aussi triviale, ainsi $\hat{\mathcal{C}}$ devient une \underline{catégorie prétopologique}. D'ailleurs
$\mathcal{C}$ en est une sous-catégorie prélocale. On dit que $\mathcal{C}$ est \underline{sursaturée}
si $\mathcal{C} = \hat{\mathcal{C}}$ . Ainsi $\hat{\mathcal{C}}$ est une catégorie prélocale sursaturée, dite
associée à $\mathcal{C}$ .

PROPOSITION 2. - Soit X un ensemble, $(U_i)$ un recouvrement de X par des
ensembles $U_i$ munis de $\mathcal{C}$-structures $S_i$. On suppose que si $U_i \cap U_j \neq \emptyset$,
$U_i \cap U_j$ est la réunion d'ensembles W ayant les propriétés suivantes : W
est ouvert dans $U_i$ et $U_j$, et $S_i | W = S_j | W$. Alors il existe sur X une
$\hat{\mathcal{C}}$-structure locale S et une seule telle que les $U_i$ soient des parties
ouvertes de X, et $S_i = S | U_i$ pour tout i. La topologie sous-jacente est la
topologie pour laquelle les $U_i$ sont ouverts et induisent sur chaque $U_i$
la topologie de $U_i$.

Considérons sur X la topologie T la moins fine induisant sur chaque
$U_i$ une topologie moins fine que la topologie $T_i$ de $U_i$. Montrons que les
$U_i$ sont ouverts pour T, et que $T | U_i = T_i$. Pour ceci, il suffit de montrer
que si V est une partie de $U_i$ ouverte pour $T_i$, alors V est ouverte pour
T, i.e. $V \cap U_j$ est ouvert pour $T_j$ quel que soit j. Or $U_i \cap U_j$ est réunion
d'ensembles W, ouverts pour $T_i$ et $T_j$, sur lesquels $S_i$ et $S_j$, donc aussi
$T_i$ et $T_j$ coïncident, il suffit de montrer que pour un tel W, $W \cap V$ est
ouvert pour $T_j$ ($V \cap U_j$ étant la réunion de tels $W \cap V$). Or $W \cap V$ est ou-
vert dans W pour $T_i$, donc pour $T_j$, donc aussi ouvert dans $U_j$ pour $T_j$.
Il reste à prouver seulement que si $U_i \cap U_j \neq \emptyset$, alors $S_i | U_i \cap U_j = S_j | U_i \cap U_j$

\newpage

%% ===== Page 8 =====

% typescript page 6 - n° 262

Or, pour les $W$ de l'énoncé, ces deux $\mathcal{C}$-structures sur $U_i \cap U_j$ induisent
la même $\mathcal{C}$-structure sur $W$, elles sont donc identiques en vertu de prop.1.
Le résultat d'unicité de la prop.2 résulte aussitôt de la prop.1, car on
voit que la topologie associée à $S$ ne peut être que $T$.

3. - \emph{Produits - Catégories topologiques}.

Soient $X$ et $Y$ deux $\mathcal{C}$-espaces. Conformément aux définitions générales
(Ens, 2ème edt. chap.IV) on dit que $X \times Y$ admet une \emph{$\mathcal{C}$-structure produit}
si elle admet une $\mathcal{C}$-structure initiale pour les applications de projec-
tion $X \times Y \longrightarrow X$ et $X \times Y \longrightarrow Y$ ; alors $X \times Y$, muni de la $\mathcal{C}$-structure pro-
duit, est aussi un produit de $X$ et $Y$ dans la catégorie $\mathcal{C}$ , i.e. les $\mathcal{C}$-
morphismes d'un $Z \in \mathcal{C}$ dans $X \times Y$ correspondent biunivoquement aux couples
de morphismes de $Z$ dans $X$ et dans $Y$. On dit qu'une catégorie prétopologi-
que $\mathcal{C}$ est une \emph{catégorie topologique} si elle satisfait à l'axiome suivant

CT 5) - \emph{Si $X,Y$ sont deux $\mathcal{C}$-espaces, $X \times Y$ admet une $\mathcal{C}$-structure produit}

D'après les propriétés générales sur les produits dans une catégorie,
(loc. cité) le produit de $\mathcal{C}$-espaces est commutatif, et associatif ;
d'après la transitivité des structures initiales, si $U$ est un ouvert du
$\mathcal{C}$-espace $X$ et $V$ un ouvert du $\mathcal{C}$-espace $Y$, alors $U \times V$ est un ouvert du
$\mathcal{C}$-espace $X \times Y$, et la $\mathcal{C}$-structure induite est identique à la $\mathcal{C}$-
structure produit de $U$ et $V$. Le fait que $U \times V$ soit ouvert résulte du fait
que les projections de $X \times Y$ sur $X$ et $Y$, étant des $\mathcal{C}$-applications, sont
continues, en d'autres termes, :

PROPOSITION 3. - Soient $\mathcal{C}$ une catégorie topologique, $X$ et $Y \in \mathcal{C}$, la topo-
logie de $X \times Y$ associée à sa $\mathcal{C}$-structure est plus fine que la topologie
produit des topologies de $X$ et $Y$.

On fera attention qu'elle peut être strictement plus fine. Sauf men-
tion expresse du contraire, quand nous considérerons une topologie sur le

\newpage

%% ===== Page 9 =====

% typescript page 7 — n° 262
produit de deux $\mathcal{C}$-espaces, ce sera celle associée à la $\mathcal{C}$-structure
produit.

Définition 1. - Soit $\mathcal{C}$ une catégorie topologique. Un $\mathcal{C}$-espace $X$ est
dit \emph{séparé} si la diagonale de $X \times X$ est fermée.

Il suffit pour ceci que la diagonale soit fermée pour la topologie
produit, i.e. que la topologie de l'espace $X$ soit séparée, (et cette con-
dition est aussi nécessaire dans le cas, peu intéressant du reste, où la
topologie de $X \times X$ est identique à la topologie produit). Si $X$ est un $\mathcal{C}$-
espace séparé, toute partie de $X$ qui admet une $\mathcal{C}$-structure induite,
en particulier toute partie ouverte de $X$, est un $\mathcal{C}$-espace séparé. Le
produit de deux $\mathcal{C}$-espaces séparés est séparé, car la diagonale de
$(X \times Y) \times (X \times Y)$ est l'intersection des images réciproques des diagonales
de $X \times X$ et $Y \times Y$ par les applications de projection, qui sont continues.

Une catégorie locale est dite \emph{séparée} si elle satisfait à :

\emph{CT 6) Tout $X \in \mathcal{C}$ est séparé.}

\emph{Sous-catégories topologiques.}

Soient $\mathcal{C}$ et $\mathcal{C}'$ deux catégories topologiques, on dit que $\mathcal{C}'$ est une
sous-catégorie topologique de $\mathcal{C}$ si c'est une sous-catégorie prétopologi-
que (cf. n°1) et si pour $X, Y \in \mathcal{C}'$, la $\mathcal{C}'$-structure produit sur $X \times Y$ est
identique à la $\mathcal{C}$-structure produit. Pour qu'une sous-classe $\mathcal{C}'$ de $\mathcal{C}$
provienne d'une sous-catégorie topologique, il faut et il suffit qu'elle
satisfasse aux conditions (i) et (ii) du n°1 et à la condition sui-
vante :

(iii) Si $X$ et $Y$ sont dans $\mathcal{C}'$, alors $X \times Y$ (qui est dans $\mathcal{C}$ à priori) est
aussi dans $\mathcal{C}'$.

(C'est trivialement nécessaire : et c'est suffisant, car on vérifie
aussitôt que la $\mathcal{C}$-structure produit est aussi une $\mathcal{C}'$-structure produit)

\newpage

%% ===== Page 10 =====

% typescript page 8 — n° 262

Si $X \in \mathcal{C}'$, il revient alors au même de dire que X est séparé dans $\mathcal{C}'$ ,
ou dans $\mathcal{C}$ .

Soit par exemple $\mathcal{C}$ une catégorie locale, soit $\mathcal{C}^s$ la classe des $\mathcal{C}$-
espaces séparés, alors les conditions (i), (ii) et (iii) sont satisfaites,
donc :

PROPOSITION 4. - La classe $\mathcal{C}^s$ des $\mathcal{C}$-espaces séparés est une sous-catégo-
rie topologique séparée de la catégorie topologique $\mathcal{C}$ .

4. - \emph{Produits de $\mathcal{C}$-espaces locaux. Catégories topologiques saturées, par-
faites.}

PROPOSITION 5. - Si $\mathcal{C}$ est une catégorie topologique, alors $\hat{\mathcal{C}}$ est une
catégorie topologique, et $\mathcal{C}$ en est une sous-catégorie topologique .

Soient X et Y deux $\hat{\mathcal{C}}$-espaces locaux, $(U_i)$ (resp. $(V_j)$) un recouvrement
de X (resp. Y) par des ouverts qui sont des $\mathcal{C}$-espaces. Alors $X \times Y$ est un
ensemble recouvert par les ensembles $U_i \times V_j$, qui d'ailleurs sont munis de
$\mathcal{C}$-structures produits. La proposition 2 s'applique, et on trouve sur
$X \times Y$ une unique $\hat{\mathcal{C}}$-structure locale telle que les $U_i \times V_j$ soient des ou-
verts, et leur $\hat{\mathcal{C}}$-structure la structure induite. Je dis que muni de cette
$\hat{\mathcal{C}}$-structure locale, $X \times Y$ est un produit de X et Y au sens de $\hat{\mathcal{C}}$ . Soit
en effet, f une application d'un $\hat{\mathcal{C}}$-espace local Z dans $X \times Y$, il faut mon-
trer que f est un morphisme si et seulement si ses composantes le sont.
C'est évidemment nécessaire, car les projections de $X \times Y$ sont manifestement
des $\hat{\mathcal{C}}$-morphismes. C'est suffisant, car alors les images réciproques des
$U_i \times V_j$ sont ouvertes, on peut alors se ramener aussitôt au cas où Z est
un $\mathcal{C}$-espace et $f(Z)$ contenu dans un des $U_i \times V_j$, auquel cas c'est immédiat.
Reste à vérifier pour $\mathcal{C}$ plongé dans $\hat{\mathcal{C}}$ la condition (iii) du numéro précé-
dent, ce qui est immédiat.

On dira que la catégorie topologique $\mathcal{C}$ est \emph{saturée} si elle satisfait

\newpage

%% ===== Page 11 =====

% typescript page 9 — n° 262

à la condition suivante :

CT 7) \emph{Tout $\hat{\mathcal{C}}$-espace local séparé est un $\mathcal{C}$-espace}, i.e. $\hat{\mathcal{C}}^s \subset \mathcal{C}$.

Une catégorie topologique est dite \emph{parfaite}, si elle est séparée et
saturée (i.e. satisfait aux axiomes CT 1) \& CT 7)). Une variante de cette
notion , rencontrée en Géométrie Algébrique, est celle de catégorie loca-
\emph{le semi-parfaite} , et obtenue en remplaçant CT 7) par l'axiome le plus
faible :

CT 7 bis) \emph{Tout $\hat{\mathcal{C}}$-espace local séparé, réunion d'un nombre fini d'ouverts
qui sont des $\mathcal{C}$-espaces, est un $\mathcal{C}$-espace, et tout $\mathcal{C}$-espace
est un espace de Zariski.}

Rappelons qu'un espace topologique est dit un espace de Zariski si
toute suite croissante de parties ouvertes est stationnaire).

PROPOSITION 6. - Soit $\mathcal{C}$ une catégorie topologique. Alors la catégorie
$\hat{\mathcal{C}}^s$ des $\mathcal{C}$-espaces locaux séparés est une catégorie locale parfaite.

C'est en effet, une catégorie lacale, séparée (prop.4) et manifeste-
ment saturée.

Nous dirons qu'une catégorie topologique est \emph{recollable}, si elle est
soit sursaturée, soit parfaite, soit semi-parfaite.

6. - \emph{Structures induites.}

Dans la suite de ce \S , $\mathcal{C}$ désigne une catégorie topologique fixée
Soit $X$ un $\mathcal{C}$-espace, $Y$ une partie de $X$. Conformément aux définitions gé-
nérales (Ens. 2è ed, chap.IV) on dit qu'une $\mathcal{C}$-structure sur $Y$ est \emph{in-
duite} par celle de $X$, ou simplement induite par $X$, si c'est une structure
initiale pour l'application $f$ de l'ensemble $Y$ dans le $\mathcal{C}$-espace $X$.
(N.B. - le foncteur fondamental $i_{\mathcal{C}}$ de $\mathcal{C}$ est donc considéré ici comme
foncteur à valeurs dans la catégorie des ensembles). On fera attention
que la topologie sous-jacente à la structure induite n'est pas

\newpage

%% ===== Page 12 =====

% typescript page 10 — n° 262

nécessairement la topologie induite (géodésiques du tore !). Rappelons
(Ens., loc cité) les faits suivants : (i) Soit $X$ un $\mathcal{C}$-espace, $Y$ une
partie de $X$ admettant une $\mathcal{C}$-structure induite, et $Z$ une partie de $Y$ ;
Pour que $Z$ admette une structure induite dans $Y$, il faut et il suffit qu'
il admette une structure induite dans $X$, et ses deux structures induites
sont alors identiques ; (ii) Soient $X_i$ ($i=1,2$) des $\mathcal{C}$-espaces, soit $Y_i$
une partie de $X_i$ admettant une $\mathcal{C}$-structure induite, alors $Y_1 \times Y_2$
est une partie de $X_1 \times X_2$ admettant une $\mathcal{C}$-structure induite, identique
à la $\mathcal{C}$-structure produit de $Y_1$ et $Y_2$. - Voici un fait spécial au contex-
te des catégories locales :

PROPOSITION 7. - Soit $X$ un $\mathcal{C}$-espace, $Y$ une partie de $X$. Si $Y$ admet une
$\mathcal{C}$-structure induite, alors pour tout ouvert $U \subset X$, $U \cap Y$ est une partie
de $U$ admettant une structure induite, identique à la $\mathcal{C}$-structure induite
sur $U \cap Y$ considéré comme partie ouverte du $\mathcal{C}$-espace $Y$. Réciproquement,
supposons que $\mathcal{C}$ soit recollable, et que $X$ admette un recouvrement ouvert
$(U_i)$ tel que pour tout $i$, $U_i \cap Y$ soit une partie de $U_i$ admettant une $\mathcal{C}$-
structure induite. Alors $Y$ admet une $\mathcal{C}$-structure induite.

Appliquons (i) à $X \supset Y \supset U \cap Y$ : comme $U \cap Y$ est une partie ouverte de $Y$,
comme image inverse de l'ouvert $U \subset X$ par l'application d'injection $Y \to X$
qui est continue, $U \cap Y$ admet une $\mathcal{C}$-structure induite dans $Y$, donc aussi
dans $X$, et ses deux $\mathcal{C}$-structures sont égales. Appliquons mainteant (i)
à $X \supset U \supset U \cap Y$, il s'ensuit que $U \cap Y$ a une structure induite dans $U$, identi-
que à sa structure induite dans $X$. Soit maintenant $(U_i)$ un recouvrement
ouvert de $X$, tel que pour tout $i$, $U_i \cap Y$ ait une structure induite dans $U_i$
(ou dans $X$, cela revient au même d'après (i) ci-dessus). Montrons d'abord
que $Y$ admet alors une $\hat{\mathcal{C}}$-structure induite, où $\hat{\mathcal{C}}$ est la catégorie topolo-
gique sursaturée associée à $\mathcal{C}$. Soit $Y_i = U_i \cap Y$, montrons que les $Y_i$

\newpage

%% ===== Page 13 =====

% typescript page 11 - n° 262
satisfont à la condition de cohérence de la proposition 2; c'est à dire
que $Y_i \cap Y_j$ est ouvert dans $Y_i$ pour tout couple $(i,j)$, et que les $\mathcal{C}$-
structures induites sur $Y_i \cap Y_j$ par $Y_i$ et $Y_j$ sont les mêmes. Or $Y_i \cap Y_j$
est l'image réciproque de la partie ouverte $U_j$ de $X$ par l'application
d'injection $Y_i \to X$ qui est continue, elle est donc bien ouverte, de plus
la structure induite sur $Y_i \cap Y_j$ par $Y_i$ est identique à celle induite
par $X$ d'après la propriété de transitivité (i), donc aussi identique à
la $\mathcal{C}$-structure induite par $Y_j$. Il existe donc une $\mathcal{C}$-structure lo-
cale sur $Y$, pour laquelle les $Y_i$ soient des parties ouvertes, et les
structures induites sur les $Y_i$ soient les structures induites par $X$.
Montrons que c'est là une \emph{$\hat{\mathcal{C}}$-structure induite par $X$ sur $Y$}. Tout
d'abord, en vertu de CT 4), l'application d'injection $Y \to X$ est un mor-
phisme, il reste donc à montrer que si $f$ est une application d'un $\hat{\mathcal{C}}$-
espace $Z$ dans $Y$, qui soit un morphisme de $Z$ dans $X$, c'est un morphisme
de $Z$ dans $Y$. Par localisation, on peut supposer que $Z \in \mathcal{C}$, et de plus
que $f$ applique $Z$ dans un des $Y_i$. Mais comme la $\mathcal{C}$-structure de $Y_i$ est
induite par celle de $X$, $f$ est un morphisme de $Z$ dans $Y_i$, donc aussi de
$Z$ dans $Y$. Notons maintenant que si la $\mathcal{C}$-structure locale mise sur $Y$,
induite par $X$ d'après ce qui précède, est une $\mathcal{C}$-structure, ce sera
aussi une $\mathcal{C}$-structure induite. Or cette condition est remplie trivia-
lement si $\mathcal{C} = \hat{\mathcal{C}}$ i.e. si $\mathcal{C}$ est sursaturée. Elle l'est aussi si $\mathcal{C}$ est
parfaite, car il suffit de prouver que la $\mathcal{C}$-structure locale de $Y$
est séparée, ce qui résulte du fait que celle de $X$ l'est et de la

PROPOSITION 8. - Soit $f$ un morphisme injectif d'un $\mathcal{C}$-espace local $Y$
dans un autre $X$. Si $X$ est séparé, $Y$ est séparé.

En effet, la diagonale de $Y \times Y$ est l'image réciproque de celle de

\newpage

%% ===== Page 14 =====

% typescript page 12 - n° 262

$X \times X$ par l'application $f \times f$. Comme cette dernière est continue puisque
c'est un morphisme, et que la diagonale de $X \times X$ est fermée par hypothèse
il en est de même de celle de $Y \times Y$. Cela prouve la proposition 8, et
achève en même temps la démonstration de la proposition 7.

\section*{7. - $\mathcal{C}$-Structures quotient.}

Soient $X$ un $\mathcal{C}$-espace, $R$ une relation d'équivalence dans $X$, $Y = X/R$,
$f$ l'application canonique de $X$ sur $Y$. Conformément aux définitions géné-
rales (loc. cité) on dit qu'une $\mathcal{C}$-structure sur $Y$ est \emph{quotient} de celle
de $X$, si c'est une structure finale pour l'application $f$ du $\mathcal{C}$-espace $X$
dans l'ensemble $Y$. Rappelons les faits suivants (loc.cité) : (i) Soit $Y$ un
ensemble quotient du $\mathcal{C}$-espace $X$, admettant une $\mathcal{C}$-structure quotient,
et soit $Z$ un ensemble quotient de $Y$ ; pour que $Z$ admette une $\mathcal{C}$-structure
quotient en tant que quotient de $Y$, il faut et il suffit qu'il en admette
une en tant que quotient de $X$, et ces deux structures quotient sont iden-
tiques. (ii) Un Buchsbaum pour le (ii) du numéro précédent, pour le plai-
sir de l'oeil. (Il y est question de $\mathcal{C}$-espace \emph{somme} de deux autres ;
il semble bien que dans une prochaine rédaction, on devra dire un ou deux
sorites sur les espaces somme : sous-trucs, quotients et produits de som-
mes sont ce qu'on pense).

Faute d'un analogue, pour la prop.7, la notion de $\mathcal{C}$-structure quotient
n'a pas toute la souplesse qu'on lui souhaiterait. Qu'à cela ne tienne,
on va l'assouplir à l'instant.

\emph{Définition 2.} - Soient $X$ un $\mathcal{C}$-espace muni d'une relation d'équivalence
$R$, $f$ l'application canonique de $X$ sur $Y = X/R$. On dit que $Y$ est un $\mathcal{C}$-
espace quotient de $X$, si $Y$ admet une $\mathcal{C}$-structure quotient, satisfaisant
à la condition suivante : la topologie sous-jacente du $\mathcal{C}$-espace $Y$ est la
topologie quotient de celle de $X$, et pour toute partie ouverte $U$ de $Y$,

\newpage

%% ===== Page 15 =====

% typescript page 13 - n° 262

|la $\mathcal{C}$-structure induite sur U par Y est une $\mathcal{C}$-structure quotient 
|de $f^{-1}(U)$.

Cette dernière condition s'énonce aussi ainsi : \emph{Si une application}
g \emph{de} U \emph{dans un} $\hat{\mathcal{C}}$\emph{-espace} Z \emph{est telle que} gf \emph{est un morphisme de} X \emph{dans}
Z\emph{, alors} g \emph{est un morphisme de} U \emph{dans} Z. Cet énoncé reste alors vrai
si Z est seulement un $\hat{\mathcal{C}}$-espace local, comme on vérifie aussitôt ,
(donc Y est aussi un $\hat{\mathcal{C}}$-espace quotient de X). Il est trivial à partir
de la définition que si Y est un $\mathcal{C}$-espace quotient de X, alors toute
partie ouverte U de Y est un $\mathcal{C}$-espace quotient de $f^{-1}(U)$ (et sa struc-
ture quotient est celle induite par Y). De même, il est immédiat que \emph{si
Y est un $\mathcal{C}$-espace quotient de X, Z un $\mathcal{C}$-espace quotient de Y, alors
Z s'identifie aussi à un $\mathcal{C}$-espace quotient de X.}

PROPOSITION 9. - Soient X un $\mathcal{C}$-espace, R une relation d'équivalen-
ce dans X, f l'application canonique de X sur Y = X/R. On munit X/R de
la topologie quotient , et on suppose que X/R est réunion d'ouverts $U_i$
tels que $U_i$ soit un $\mathcal{C}$-espace quotient de $f^{-1}(U_i)$. Alors Y est un $\hat{\mathcal{C}}$ -
espace quotient de X.

En effet, $U_i$ et $U_j$ induisent la même structure sur leur partie ouver-
te $U_i \cap U_j$ , à savoir la $\mathcal{C}$-structure quotient de $f^{-1}(U_i \cap U_j)$. Il existe
donc une $\hat{\mathcal{C}}$-structure sur Y dont la topologie sous-jacente est celle
de Y, et induisant sur les $U_i$ la $\mathcal{C}$-structure qui y est donnée. L'appli-
cation f est un morphisme de X dans Y, puisque ses restrictions aux par-
ties ouvertes $f^{-1}(U_i)$ recouvrant X le sont. Pour prouver que Y muni de
sa $\hat{\mathcal{C}}$-structure est bien un $\hat{\mathcal{C}}$-espace quotient, il reste à montrer que
si g est une application d'une partie ouverte U de Y dans un $\hat{\mathcal{C}}$-espace Z
telle que gf soit un morphisme de $f^{-1}(U)$ dans Z, alors g est un morphisme
de U dans Z. Il suffit pour ceci de prouver que sa restriction à chaque

\newpage

%% ===== Page 16 =====

% - 14 - n° 262

$U \cap U_i$ est un morphisme, pour la $\hat{\mathcal{C}}$-structure induite par U. Or c'est aus-
si celle induite par Y, donc celle induite par $U_i$, de sorte qu'on est
ramené à montrer que l'application gf de $f^{-1}(U \cap U_i)$ dans Z est un morphis-
me, ce qui est bien le cas puisque c'est la restriction d'un morphisme
de $f^{-1}(U)$ dans Z.

Corollaire - Supposons $\mathcal{C}$ recollable. Soient X un $\mathcal{C}$-espace, R une rela-
tion d'équivalence dans X, f l'application canonique de X sur Y = X/R. On
munit Y de la topologie quotient, et on suppose que pour tout couple de
deux points de Y, existe un ouvert U les contenant tel que U soit un $\mathcal{C}$-
espace quotient de $f^{-1}(U)$. Alors Y est un $\mathcal{C}$-espace quotient de X.

On sait déjà en vertu de prop. 9 que Y est un $\hat{\mathcal{C}}$-espace quotient de X,
il reste donc à montrer que Y est en fait un $\mathcal{C}$-espace. On peut supposer
$\mathcal{C}$ parfaite ou semi-parfaite. Dans le deuxième cas, on voit aussitôt que
Y est recouvert par un nombre \emph{fini} d'ouverts U qui sont des $\mathcal{C}$-espaces
quotients des $f^{-1}(U)$. Donc on est ramené à prouver que Y est séparé. Or
tout couple de deux points de Y est contenu dans un ouvert U tel que U
soit (pour la $\hat{\mathcal{C}}$-structure induite par Y) un $\mathcal{C}$-espace quotient de $f^{-1}(U)$
et en particulier séparé. On en conclut aussitôt que Y est séparé , en
vertu du lemme suivant, dont la démonstration est laissée au lecteur,

Lemme - Soit $X \in \hat{\mathcal{C}}$. Supposons que pour tout couple de deux points de X,
existe un ouvert les contenant qui soit séparé pour la $\hat{\mathcal{C}}$-structure in-
duite. Alors X est un $\hat{\mathcal{C}}$-espace séparé.

La notion de "variété quotient" utilisée par Chevalley dans son sémi-
naire n'est pas équivalente à celle de la définition 2, et n'est ni plus
forte, ni moins forte ; elle peut s'énoncer en termes morphiques, mais
la condition est assez canularesque. En renforçant légèrement la défini-
tion de Chevalley, on obtient la notion suivante, que peut être il faudra

\newpage

%% ===== Page 17 =====

% typescript page 15 — n° 262

\begin{center}
- 15 - \hfill n$^{\circ}$ 262
\end{center}

substituer à la définition 2 : $Y = X/R$ est dit un $\mathcal{C}$-espace quotient
strict de $X$, si pour toute partie ouverte $U$ de $X$, $U/R_U$ admet une $\mathcal{C}$-
structure quotient, et l'application canonique $U/R_U \to X/R$ est un iso-
morphisme sur une partie ouverte de $X/R$. (C'est la condition Chevalley,
plus le fait que l'application $X \to X/R$ soit ouverte). Nous n'utilise-
rons cette notion que dans la proposition suivante :

PROPOSITION 10. - Soit $X$ un $\hat{\mathcal{C}}$-espace muni d'une relation d'équivalence
ouverte $R$. Soit $Y = X/R$ muni de la topologie quotient, et soit $f : X \to Y$
l'application canonique. Les conditions suivantes sont équivalentes.

a) $X/R$ est un $\hat{\mathcal{C}}$-espace quotient strict de $X$.

b) Pour tout $(x,y) \in R$, existe un ouvert $W$ contenant $x$ et $y$, tel que
$W/R_W$ soit un $\hat{\mathcal{C}}$-espace quotient strict de $W$.

c) Pour tout $(x,y) \in R$, existent des voisinages ouverts arbitrairement
petits $U$ de $x$ et $V$ de $y$, tels que $f(U) = f(V)$, que $U/R_U$ et $V/R_V$
soient des $\hat{\mathcal{C}}$-espaces quotients de $U$ resp. $V$ (déf. 2) et que l'appli-
cation composée $U/R_U \to f(U) = f(V) \to V/R_V$ soit un isomorphisme.

\underline{Démonstration.} - L'implication a) $\Rightarrow$ b) est triviale, il suffit de pren-
dre un voisinage ouvert $W'$ de $f(x) = f(y)$ tel que $W'$ soit un $\hat{\mathcal{C}}$-espace
pour la structure induite par $Y$, et de poser $W = f^{-1}(W')$. Prouvons
b) $\Rightarrow$ c). Soit $(x,y) \in R$, soient $U_0$ un voisinage ouvert de $x$, $V_0$ un voi-
sinage ouvert de $y$, alors $f(U_0)$ et $f(V_0)$ sont des voisinages ouverts de
$z = f(x) = f(y)$, donc $A = f(U_0) \cap f(V_0)$ est un voisinage ouvert de $z$.
$W$ étant comme la condition b), on peut supposer $U_0 \subset W$, $V_0 \subset W$. Soit
$U = U_0 \cap f^{-1}(A)$, $V = V_0 \cap f^{-1}(A)$, on a alors $f(U) = f(V) = A$ ; d'autre part
comme $W/R_W$ est par hypothèse un $\hat{\mathcal{C}}$-espace quotient strict de $W$, et que
$U$ est un ouvert de $W$, il en résulte que $U/R_U$ est un $\hat{\mathcal{C}}$-espace quotient
de $U$, s'identifiant à $A$ muni de la structure induite par $W/R_W$. De même

\newpage

%% ===== Page 18 =====

% typescript page 16 — n° 262

$V/R_V$ est un $\mathcal{C}$-espace quotient de $V$ s'identifiant à $A$, d'où résulte que
$U$ et $V$ satisfont aux conditions exigées dans a). Prouvons enfin que c)
implique a). Pour tout ouvert $U \subset X$ tel que $U/R_U$ soit un $\mathcal{C}$-espace quo-
tient de $U$, soit $Y_U = f(U)$ muni de la structure de $\mathcal{C}$-espace déduite de
celle de $U/R_U$ par transport de structure. Comme la topologie sous-jacente
du $\mathcal{C}$-espace $U/R_U$ est la topologie quotient de $U$, et que $f$ est ouverte,
la topologie de $\mathcal{C}$-espace $Y_U$ est celle induite par $Y$. L'hypothèse impli-
que que la réunion des $Y_U$ est $Y$. Nous allons montrer que la famille des
$Y_U$ satisfait aux conditions de cohérence de prop. 2 n° 2. Soit
$z \in Y_U \cap Y_V$ , on a $z = f(x) = f(y)$, avec $x \in U$, $y \in V$. On peut (par hypothèse)
trouver des voisinages ouverts $U' \subset U$ de $x$ et $V' \subset V$ de $y$, tels que
$Y_{U'} = Y_{V'}$ (avec identité aussi pour les $\mathcal{C}$-structures), c'est une par-
tie ouverte $W$ de $Y$, donc de $Y_U$ et $Y_V$, contenant $z$, il reste à montrer que
sa $\mathcal{C}$-structure est celle induite par $Y_U$, donc aussi celle induite par
$Y_V$, de sorte que ces deux structures induites coïncident. Cela nous ramène
à montrer que dans le cas où $U \subset V$, la $\mathcal{C}$-structure $S$ de $Y_U$ est identi-
que à la $\mathcal{C}$-structure $S'$ induite par celle de $Y_V$. Il est immédiat qu'elle
est plus fine que la structure induite, il reste à montrer que l'applica-
tion identique de $(Y_U,S')$ dans $Y_U = (Y_U,S)$ est un morphisme. Comme $Y_V$ est
un $\mathcal{C}$-espace quotient de $V$, cela signifie aussi que l'application corres-
pondante $f'$ de $V' = V \cap f^{-1}(Y_U)$ dans $Y_U$ est un morphisme, et pour ceci
il suffit de vérifier que tout point $y$ de $V'$ a un voisinage ouvert
$V_1 \subset V'$ tel que $f'|V_1$ soit un morphisme de $V_1$ dans $Y_U$. Or soit $x \in U$ tel
que $f(x) = f(y)$, il existe alors des voisinages ouverts $V_1 \subset V$ de $y$ et
$U_1 \subset U$ de $x$ tels que $Y_{U_1} = Y_{V_1}$ (identique de $\mathcal{C}$-espace). Or $f'|V_1$ est
un morphisme de $V_1$ dans $Y_{V_1}$ donc aussi dans $Y_{U_1}$ , donc aussi un morphisme

\newpage

%% ===== Page 19 =====

% typescript page 17 — n° 262
de $Y_{V_1}$ dans $Y_U$. Cela achève de montrer que les conditions de prop.2 sont
satisfaites. Il existe par suite sur $Y$ une $\hat{\mathcal{C}}$-structure locale unique pour
laquelle les $Y_U$ sont ouverts et qui induise sur tout $Y_U$ la structure
donnée sur $Y_U$. La topologie sous-jacente est identique à la topologie
quotient sur $Y$. Montrons que $Y$ est un $\hat{\mathcal{C}}$-espace quotient de $X$ pour
cette $\hat{\mathcal{C}}$-structure. Comme $f$ est évidemment un morphisme de $X$ dans $Y$
(en vertu de \underline{CT} 4)), il reste à vérifier que si $W$ est une partie ouverte
de $Y$, $g$ une application de $W$ dans un $\hat{\mathcal{C}}$-espace $Z$ telle que l'application
$gf$ de $f^{-1}(W)$ dans $Z$ soit un morphisme, alors $g$ est un morphisme. Or $W$
est réunion d'ouverts du type $Y_U$, et la structure $Y_U$ étant celle induite
par $Y$ i.e. celle induite par $W$, il suffit de prouver que la restriction
de $g$ à $Y_U$ est un morphisme, i.e. que la restriction de $gf$ à $U$ est un mor-
phisme, ce qui est évidemment le cas.

\underline{Corollaire 1.} - Avec les notations de la prop.10, supposons de plus que
$\mathcal{C}$ soit recollable. Alors les conditions suivantes sont équivalentes :

a) $X/R$ est un $\mathcal{C}$-espace quotient strict de $X$.

b) Pour tout couple $(x,y)$ de points de $X$ (non nécessairement congrus
mod. $R$ !.) existe un ouvert $W$ les contenant, tel que $W/R_W$ soit un $\mathcal{C}$-
espace quotient strict de $W$.

Quand $\mathcal{C}$ est sursaturée, cet énoncé est un cas particulier de la
proposition 10 (l'implication a) $\Rightarrow$ b) étant triviale à priori). Si $\mathcal{C}$
est parfaite ou semi-parfaite, on est ramené, comme dans le corollaire à
prop.9, à prouver que $Y$ est un $\mathcal{C}$-espace local \emph{séparé}, Or il résulte de
la condition b) que deux points quelconques de $Y$ sont contenus dans un
ouvert qui est un $\mathcal{C}$-espace pour la structure induite, \emph{donc} séparé, d'où
résulte que $Y$ est séparé, en vertu du lemme donné plus haut.

\underline{Corollaire 2.} - Soient $X$ un $\mathcal{C}$-espace muni d'un groupe $G$ d'automorphis-
mes,

\newpage

%% ===== Page 20 =====

% typescript page 18 - n° 262

$Y = X/G$ l'espace des orbites, $f:X \to Y$ l'application canonique de $X$ sur $Y$.
On suppose que pour tout $x \in X$, il existe un système fondamental de voi-
sinages ouverts $U$ de $x$ tels que $U/R_U$ soit un $\mathcal{C}$-espace quotient de $U$
(où $R_U$ désigne la relation d'équivalence sur $U$ induite par la relation
d'équivalence définie par $G$). Alors $Y$ est un $\mathcal{C}$-espace quotient de $X$, et
pour tout ouvert $U \subset X$, $U/R_U$ est un $\mathcal{C}$-espace quotient de $U$, s'identifiant
à un sous-espace ouvert de $Y$ muni de la $\mathcal{C}$-structure induite par $Y$.

En effet, la condition c) de la proposition 10 est visiblement satis-
faite. Le fait que $R$ soit ouverte résulte du fait que pour tout ouvert
$U$, le saturé de $U$ pour $R$ est identique à la réunion des ouverts $g.U (g \in G)$
donc lui-même ouvert.

Définition 3. - Soient $X,Y$ deux $\mathcal{C}$-espaces, $f$ un morphisme de $X$ sur $Y$.
Une application $s$ de $Y$ dans $X$ est dite une $\mathcal{C}$-section pour $f$ si c'est un
$\mathcal{C}$-morphisme et si $fs$ est l'application identique de $Y$. On dit que $f$
admet localement une $\mathcal{C}$-section si $Y$ est réunion d'ouverts $U$ tels que
les morphismes $f^{-1}(U) \to U$ induits par $f$ admettent une $\mathcal{C}$-section.

PROPOSITION 11. - Soient $X,Y$ deux $\mathcal{C}$-espaces $f$ un morphisme de $X$ sur $Y$,
$R$ la relation d'équivalence définie par $f$, $f'$ la bijection de $X/R$ sur $Y$
déduite de $f$. Si $s$ est une $\mathcal{C}$-section pour $f$, $s$ est un isomorphisme sur
une partie $Y'$ de $X$ ayant une structure induite par $X$ ; si $X$ est un $\mathcal{C}$-
espace séparé, $Y'$ est une partie fermée de $X$. Si $f$ admet \emph{localement} une
$\mathcal{C}$-section, alors $X/R$ est un $\hat{\mathcal{C}}$-espace quotient de $X$ et $f'$ est un
isomorphisme.

Supposons qu'il existe une $\mathcal{C}$-section $s$, montrons que la $\mathcal{C}$-structure
sur $X/R$ déduite de celle de $Y$ par transport de structure fait de $X/R$ un
$\mathcal{C}$-espace quotient de $X$. Il revient au même de dire, puisque $f$ est
déjà un morphisme, que toute application $g$ d'une partie ouverte $U$ de $Y$

\newpage

%% ===== Page 21 =====

% typescript page 19 — n° 262

dans un $\mathcal{C}$-espace $Z$, telle que $gf$ soit un morphisme de $f^{-1}(U)$ dans $Z$,
est elle-même un morphisme. En effet, on a $g = g(fs) = (gf)s$, qui est
un morphisme comme composé de deux morphismes $U \xrightarrow{s} f^{-1}(U) \xrightarrow{gf} Z$. Mettons sur
$Y' = s(Y)$ la structure de $\mathcal{C}$-espace déduite de celle de $Y$ par transport
de structure, montrons que c'est là une $\mathcal{C}$-structure induite par $X$.
Comme l'application d'injection de $Y'$ dans $X$ est déjà un morphisme
(puisque $s$ l'est) il faut seulement vérifier que si $g$ est une application
d'un $\mathcal{C}$-espace $Z$ dans $Y$ telle que $sg$ soit un morphisme, alors $g$ est
un morphisme ; or on a $g = (fs)g=f(sg)$, c'est donc un morphisme comme
composé des deux morphismes $sg$ et $f$. Soit $p = sf$, c'est un morphisme de
$X$ dans $X$, et $Y'$ est l'ensemble des $x \in X$ tels que $px = x$, i.e. l'image
inverse de la diagonale de $X \times X$ par le morphisme $x \to (x,px)$ de $X$
dans $X \times X$. Si la diagonale de $X \times X$ est fermée, c'est donc une partie fer
mée de $X$. Supposons enfin que $f$ n'ait pas nécessairement de $\mathcal{C}$-section
mais admette localement une $\mathcal{C}$-section. Alors la première assertion de
la proposition est valable, comme il résulte aussitôt de ce qu'on vient
de dire, et de proposition 9.

Corollaire. - Soient $X_i, Y_i$ des $\mathcal{C}$-espaces ($i=1,2$), $f_i$ un morphisme
de $X_i$ sur $Y_i$ admettant localement une $\mathcal{C}$-section. Alors $f_1 \times f_2$ est un
morphisme de $X_1 \times X_2$ sur $Y_1 \times Y_2$ admettant localement une $\mathcal{C}$-section, et
par suite sur $Y_1 \times Y_2$ la $\mathcal{C}$-structure produit des structures quotients
des $Y_i$ est identique à une $\mathcal{C}$-structure quotient du $\mathcal{C}$-espace produit
$X_1 \times X_2$.

8. - \underline{$\mathcal{C}$-espaces réduits à un point.}

(Le rédacteur laisse aux spécialistes de faire un n$^\circ$ étoffé sur les
$\mathcal{C}$-espaces vides). Dans les catégories topologiques que nous aurons à
considérer, l'axiome suivant sera satisfait :

\newpage

%% ===== Page 22 =====

% typescript page 20 — n° 262

\underline{CT} 8) \emph{Il existe un $\mathcal{C}$-espace réduit à un point. Si $e$ est un tel $\mathcal{C}$-espace, alors toute application d'un $\mathcal{C}$-espace $X$ dans $e$ est un morphisme.}

Soient $e$, $e'$ deux $\mathcal{C}$-espaces réduits à un point, soit $f$ (resp. $f'$) l'unique application de $e$ dans $e'$ (resp. de $e'$ dans $e$). D'après \emph{CT} 8), ce sont là des morphismes, réciproques l'un de l'autre, donc en fait des isomorphismes. En particulier, prenant $e=e'$, on voit qu'il n'existe qu'une \emph{seule structure de $\mathcal{C}$-espace sur un ensemble $e$ réduit à un point} : on supposera alors toujours implicitement que $e$ est muni de cette structure. \emph{Si alors $X \in \mathcal{C}$, l'application des projections $X \times e \to X$ est un $\mathcal{C}$-isomorphisme.}

PROPOSITION 12. - Supposons que $\mathcal{C}$ soit une catégorie prétopologique satisfaisant l'axiome \emph{CT} 8) ci-dessus. Soient $X$ un $\mathcal{C}$-espace, $x$ un point de $X$. Pour que $\{x\}$ admette une $\mathcal{C}$-structure induite, il faut et il suffit que toute application constante d'un $\mathcal{C}$-espace $Y$ dans $X$, prenant la valeur $x$, soit un morphisme, ou encore que l'application identique de $\{x\}$ dans $X$ soit un morphisme.

Beweis klar. Le lecteur notera que dans le cas important de la catégorie des espaces algébriques définis sur un corps $k$, et dont le corps des valeurs $K$ est $\neq k$, (cf par. 3), il y a en général des points de $X$ qui ne satisfont pas les conditions de la proposition 12.

Corollaire. - Soit $\mathcal{C}$ une catégorie topologique, satisfaisant à l'axiome \emph{CT} 8) soient $X, Y$ deux $\mathcal{C}$-espaces, $x$ un point de $X$ tel que $\{x\}$ ait une $\mathcal{C}$-structure induite. Alors l'application $y \to (x, y)$ de $Y$ dans $X \times Y$ est un morphisme et même un $\mathcal{C}$-isomorphisme de $Y$ sur une partie de $X \times Y$ munie d'une $\mathcal{C}$-structure induite.

Cette application $i$ est en effet composée de deux applications

\newpage

%% ===== Page 23 =====

% typescript page 21 - n° 262

$Y \to \{x\} \times Y \to X \times Y$, la seconde est un morphisme comme produit de deux
morphismes, et la première est aussi un morphisme, car ses deux appli-
cations composantes $Y \to \{x\}$ et $Y \to Y$ sont des morphismes. La dernière
résulte du fait que $i$ admet un morphisme inverse à gauche, savoir
l'application $(a,b) \to (x,b)$, (qui est bien un morphisme comme produit
de deux morphismes)

9. - \emph{Espaces fibrés.}

Soient $\mathcal{C}$ une catégorie topologique ("catégorie base"), $\underline{\Phi}_0$ et $\underline{\Phi}$ deux
catégories "catégories fibres", $\mathcal{E}$ désignant la catégorie des ensem-
bles, on suppose donnés, en plus du foncteur fondamental $\mathcal{C} \to \mathcal{E}$, des
foncteurs covariants $\underline{\Phi}_0 \to \mathcal{C}$ et $\underline{\Phi}_0 \to \underline{\Phi} \to \mathcal{E}$, de telle façon que
le diagramme suivant de foncteurs
\begin{tikzcd}
& \underline{\Phi}_0 \arrow[dl] \arrow[dr] & \\
\underline{\Phi} \arrow[dr] & & \mathcal{C} \arrow[dl] \\
& \mathcal{E} & 
\end{tikzcd}
soit commutatif . On ne suppose pas que ces trois foncteurs supplémen-
taires soient fidèles, néanmoins, suivant notre habitude, nous considé-
rerons les éléments de $\underline{\Phi}$ et de $\underline{\Phi}_0$ comme des ensembles munis de struc-
tures de type $\underline{\Phi}$ (resp. de type $\underline{\Phi}_0$) (mais un morphisme dans $\underline{\Phi}_0$ pour-
rait ne pas être déterminé sans ambiguïté par l'application d'ensembles
qu'elle définit). Ainsi une structure de type $\underline{\Phi}_0$ porte des structures
"sous-jacentes" de type $\mathcal{C}$ et de type $\underline{\Phi}$. Il sera commode de supposer
que le foncteur
$\underline{\Phi}_0 \to \underline{\Phi}$
est fidèle et transportable, et que si $A,B \in \underline{\Phi}_0$, un $\underline{\Phi}$-morphisme $A \to B$
est un $\underline{\Phi}_0$-morphisme si et seulement si l'application qu'il définit est
un $\mathcal{C}$-morphisme. Le cas le plus important est celui où $\underline{\Phi}_0 = \underline{\Phi}$

\newpage

%% ===== Page 24 =====

% typescript page 22 — n° 262

Le cas $\Phi \neq \Phi_0$ se rencontre en Géométrie algébrique, $\mathcal{C}$ étant la caté-
gorie des $(k,K)$-espaces algébriques, alors $\Phi_0$ sera par exemple, la caté
gorie des $K$-espaces vectoriels de dimension finie définis sur $k$, tandis
que $\Phi$ sera la catégorie des $K$-espaces vectoriels sans plus. Cf;par.2,
n°6 et suivants,

\underline{Définition 4}. - $\mathcal{C}$, $\Phi$ et $\Phi_0$ étant comme ci-dessus, on appelle fibré de
type ($\mathcal{C}$, $\Phi$, $\Phi_0$) (ou simplement fibré de type ($\mathcal{C}$, $\Phi$) si $\Phi = \Phi_0$, ou
enfin simplement $\mathcal{C}$-fibré si $\mathcal{C} = \Phi = \Phi_0$) l'objet $\xi$ formé par : (i) un
morphisme $p$ (la \emph{projection}) d'un $\mathcal{C}$-espace $E$ (l'\emph{espace fibré}) sur un $\mathcal{C}$-
espace $X$ (la \emph{base}) (ii) pour tout $x \in X$, la donnée d'une $\Phi$-structure sur
la fibre $f^{-1}(x)$ de $x$ ; ces données étant assujetties à l'axiome suivant :

\underline{LF} ) $X$ est réunion d'ouverts $U$ tels que l'objet "induit" par $\xi$ sur $U$
soit isomorphe à un produit $U \times F$, où $F \in \Phi_0$.

Il est clair ce qu'il faut entendre par l'objet induit, et aussi quel-
les structures (i) et (ii) définit un produit $U \times F$ ; On met sur $U \times F$ la
structure de $\mathcal{C}$-espace produit, on considère le morphisme canonique
$U \times F \to U$, les fibres de ce fibré ensembliste s'identifient toutes à $F$, et
de ce fait ont une $\Phi$-structure sous-jacente. On ne suppose pas que le
choix de $F$ puisse être pris indépendant de $U$, encore moins qu'il figure
dans la structure de $E$. Si on peut choisir un $F$ indépendant de $U$, on dit
que l'espace fibré est \emph{monotype}, et $F$ est alors appelé une \emph{fibre type} de
$\xi$. Même dans ce cas, $F$ n'est pas nécessairement déterminé par $\xi$, même
à isomorphisme près dans $\Phi_0$. Cependant, si un $x \in X$ est tel que $\{x\}$ ad-
mette une $\mathcal{C}$-structure induite, alors la fibre $E_x$ de $E$ sur $x$ est canoni-
quement munie d'une structure de type $\Phi_0$ (résulte du corollaire à prop.
12), et s'il existe une fibre-type $F$, elle est $\Phi_0$-isomorphe à $E_x$.

\newpage

%% ===== Page 25 =====

% typescript page 23 — n° 262

\emph{Définition 5.} - Soient $\xi = (E,p,X)$ et $\xi' = (E',p',X')$ deux fibrés de
type $(\mathcal{C},\underline{\Phi},\underline{\Phi}_0)$. On appelle $(\mathcal{C},\underline{\Phi})$-\emph{morphisme} de $\xi$ dans $\xi'$
toute application $f$ de $E$ dans $E'$, ayant les propriétés suivantes : $f$
applique fibre dans fibre et est un $\underline{\Phi}$-morphisme sur les fibres ;
l'application $f : E \longrightarrow E'$ et l'application $X \longrightarrow X'$ déduite de $f$ par
passage au quotient est un $\mathcal{C}$-morphisme.

Pour la notion évidente de composition des morphismes, on voit que
les $(\mathcal{C},\underline{\Phi},\underline{\Phi}_0)$-fibrés forment une \emph{catégorie}. Le foncteur
$\xi = (E,p,X) \longrightarrow E$ à valeurs dans $\mathcal{C}$ est \emph{transportable} et \emph{fidèle}
(d'où l'identification de $\xi$ à $E$) ; le foncteur $\xi = (E,p,X) \longrightarrow X$ de la
catégorie des $(\mathcal{C},\underline{\Phi},\underline{\Phi}_0)$-fibrés dans $\mathcal{C}$ est appelé \emph{foncteur-base}.
D'autre part, si $E$ est un $(\mathcal{C},\underline{\Phi},\underline{\Phi}_0)$-fibré sur $X$, on a la notion
de \emph{restriction} de ce fibré à une partie ouverte $U$ de $X$, ce sera encore
un $(\mathcal{C},\underline{\Phi},\underline{\Phi}_0)$-fibré. On vérifie aussitôt qu'on a là les données
pour une \emph{catégorie prélocale} (Top. Elém. chap. I)

PROPOSITION 13. - \emph{Si la catégorie topologique $\mathcal{C}$ est recollable, alors
les fibrés de type $(\mathcal{C},\underline{\Phi},\underline{\Phi}_0)$ forment une catégorie locale.}

Il faut prouver, dans chacun des cas nommés dans la propoisition,
le critère de recollement : On a un $\mathcal{C}$-espace $X$, un recouvrement ou-
vert $(U_i)$ de $X$, sur tout $U_i$ un fibré $E_i$ de type $(\mathcal{C},\underline{\Phi},\underline{\Phi}_0)$, et sur
$U_i \cap U_j$ un isomorphisme $f_{ij}$ de $E_j|U_{ij}$ sur $E_i|U_{ij}$, satisfaisant aux con-
ditions de cohérence $f_{ik} = f_{ij}f_{jk}$. On veut trouver un fibré $E$, sur $X$,
et des isomorphismes $E_i \longrightarrow E|U_i$, de telle façon qu'on ait $f_{ij} = f_i^{-1} f_j$
sur $U_{ij}$ (avec les abus d'écriture habituels). On sait (loc. cité)
qu'on peut trouver un espace topologique $E$ et une application continue
$p$ de $E$ sur $X$, enfin des homéomorphismes $f_i : E_i \longrightarrow E|U_i$ tels qu'on ait

\newpage

%% ===== Page 26 =====

% typescript page 24 - n° 262

$f_{ij} = f_i^{-1} f_j$. De plus, pour tout $x \in X$, mettons sur la fibre $p^{-1}(x) = E_x$
une structure de type $\underline{\Phi}$, en prenant un $U_i \ni x$ et la structure déduite
de celle de $E_{ix} = p_i^{-1}(x)$ par transport de structure à l'aide de $f_i \vert E_{ix}$.
En vertu de $f_j = f_i f_{ij}$, la structure ainsi mise sur $E_x$ est indépendante
du choix de $i$. De plus, soit $E'_i = p^{-1}(U_i)$, munissons $E'_i$ de la $\mathcal{C}$-structure
déduite de la $\mathcal{C}$-structure de $E_i$ par transport de structure à l'aide
de $f_i$. Les $\mathcal{C}$-structures des ouverts $E'_i$ et $E'_j$ de $E$ induisent la même
$\mathcal{C}$-structure dans $E'_i \cap E'_j = p^{-1}(U_{ij})$, puisque $f_i^{-1} f_j = f_{ij}$ est un $\mathcal{C}$-
isomorphisme $E_j \vert U_{ij}$ sur $E_i \vert U_{ij}$. Il s'ensuit qu'il y a sur $E$ une $\mathcal{C}$-
structure \emph{locale} dont la topologie sous-jacente soit celle de $E$,
et induisant sur les ouverts $E'_i$ la $\mathcal{C}$-structure donnée, i.e. telle
que les $f_i$ soient des $\hat{\mathcal{C}}$-isomorphismes de $E_i$ sur les ouverts $E'_i$. Il
est clair maintenant que $E$ est un fibré de type $(\hat{\mathcal{C}}, \hat{\underline{\Phi}}, \hat{\underline{\Phi}}_0)$, et les
$f_i$ des isomorphismes de $(\hat{\mathcal{C}}, \hat{\underline{\Phi}}, \hat{\underline{\Phi}}_0)$-fibrés ayant les propriétés re-
quises. Il reste à montrer que $E$ est en fait un $\mathcal{C}$-espace. C'est
trivial si $\mathcal{C}$ est sursaturée i.e. $\mathcal{C} = \hat{\mathcal{C}}$. Montrons maintenant
que si $X$ et les $E_i$ sont des $\mathcal{C}$-espaces séparés, alors $E$ est un $\mathcal{C}$-
espace séparé, i.e. la diagonale $D$ de $E \times E$ est fermée. Comme $E \times E$ est
réunion des ouverts $E'_i \times E'_j$, il suffit de montrer que $D \cap (E'_i \times E'_j)$ est
fermé dans $E'_i \times E'_j$, ou encore que l'ensemble $A$ des points $(u,v)$ de
$E_i \times E_j$ tels que $f_i(u) = f_j(v)$ est fermé. Or cet ensemble est contenu
dans l'ensemble $B$ des points $(u,v)$ tels que $p_i(u) = p_j(v)$, qui est
fermé puisque c'est image réciproque de la diagonale de $X \times X$ (fermée
par hypothèse) par l'application continue $p_i \times p_j$ de $E_i \times E_j$ dans $X \times X$.
Comme d'autre part $B$ est contenu dans l'ensemble $F_i \times F_j$, où $F_i = E_i \vert U_{ij}$,
$F_j = E_j \vert U_{ij}$, il suffit donc de montrer que $A$ est fermé dans $F_i \times F_j$.
Or $A$ est aussi l'ensemble des points de $F_i \times F_j$ tels que $u = f_{ij}(v)$,

\newpage

%% ===== Page 27 =====

% typescript page 25 - n° 262
i.e, l'image répiproque de la diagonale de $F_i \times F_i$ par l'application
produit $1 \times f_{ij}$, qui est continue puisque c'est un morphisme. Comme
$E_i$ donc $F_i$ est séparé, la diagonale de $F_i \times F_i$ est fermée, d'où s'ensuit
bien que A est fermé dans $F_i \times F_j$. - Ceci prouve en particulier que si
$\mathcal{C}$ est parfaite, alors E, étant un $\mathcal{C}$-espace local \emph{séparé} d'après
ce qui précède, est un $\mathcal{C}$-espace. Enfin, si $\mathcal{C}$ est semi-parfaite,
alors il existe une sous famille, finie de $(V_i)$ recouvrant X, donc E,
étant un $\mathcal{C}$-espace local \emph{séparé} \emph{réunion finie} de parties ouvertes qui
sont des $\mathcal{C}$-espaces, est lui-même un $\mathcal{C}$-espace. Cela achève la dé-
monstration de la prop. 13.

Corollaire. - Soit $\mathcal{C}$ une catégorie topologique, recollable. On
suppose donnés de plus $\underline{\Phi}, \underline{\Phi}_0$ comme ci-dessus. Soient $F \in \underline{\Phi}_0$, $X \in \mathcal{C}$
G le groupe des $\underline{\Phi}$-automorphismes de F, et soit $\mathcal{F}$ le faisceau des
germes d'applications de X dans G, dont les sections sur un ouvert U
sont les applications g de U dans G telles que $(x,y) \rightarrow g(x)y$ et
$(x,y) \rightarrow g(x)^{-1}y$ soient des $\mathcal{C}$-morphismes de $U \times F$ dans F. $\mathcal{F}$ s'iden-
tifie au faisceau des germes d'automorphismes du fibré de type $(\mathcal{C}, \underline{\Phi},$
$\underline{\Phi}_0)$ trivial $U \times F$. Sur un ensemble E, la donnée d'une structure de fi-
bré de type $(\mathcal{C}, \underline{\Phi}, \underline{\Phi}_0)$ homogène de fibre type F et de base X, est
équivalent à la donnée d'une structure de fibré à faisceau structural
de type $(\mathcal{F}, [unclear: U \times F])$.

Il résulte en effet, des définitions que les automorphismes du
fibré $U \times F$ de type $(\mathcal{C}, \underline{\Phi}, \underline{\Phi}_0)$ sont les applications $(x,y) \rightarrow (x,g(x)y)$
où g est une application de U dans $G = \text{Aut}(F)$ satisfaisant aux con-
ditions énoncées, d'où la première partie du corollaire (valable sans
hyptohèse sur $\mathcal{C}$). Or les fibrés homogènes de type $(\mathcal{C}, \underline{\Phi}, \underline{\Phi}_0)$ de
base X, et de fibre type F, sont par définition les fibrés de type

\newpage

%% ===== Page 28 =====

% typescript page 26 - n° 262

($\mathcal{C}, \underline{\Phi}, \underline{\Phi}_0$) sur X localement isomorphes à $X \times F$. Comme les fibrés de
type ($\mathcal{C}, \underline{\Phi}, \underline{\Phi}_0$) forment une catégorie locale, la dernière assertion
est un cas particulier de Top. Elém. chap I.... On voit donc que si
on ne s'intéresse qu'aux fibrés homogènes de fibre donnée F, ceux-ci
sont déterminés par la seule donnée d'un $\mathcal{C}$-espace F muni d'un
certain sous-groupe G du groupe des $\mathcal{C}$-automorphismes de F. La
catégorie $\underline{\Phi}$ n'intervient que pour préciser la notion de morphisme
de fibrés.

Corollaire 2. - Sous les conditions du corollaire 1, l'ensemble des classes
de fibrés homogènes de type ($\mathcal{C}, \underline{\Phi}, \underline{\Phi}_0$) sur X, de fibre type F, s'i-
dentifie à l'ensemble de cohomologie $H^1(X, \underline{\mathcal{F}})$.

10. - \emph{$\mathcal{C}$-espaces fibrés à $\mathcal{C}$-groupe structural.}

On suppose que la catégorie locale $\mathcal{C}$ satisfait l'axiome CT 8) du
n° 8. Conformément aux définitions générales (??) on appelle $\mathcal{C}$-
groupe un ensemble G muni d'une $\mathcal{C}$-structure et d'une structure
de groupe, compatibles dans le sens suivant : les applications $(x,y) \mapsto$
$\mapsto xy$ et $x \mapsto x^{-1}$ de $G \times G$ dans G et de G dans G sont des $\mathcal{C}$-morphismes,
et l'élément unité de G admet une $\mathcal{C}$-structure induite (cf n° 8). Une
application $G \to G'$ est dite un morphisme de $\mathcal{C}$-groupes, si c'est un
morphisme pour les deux structures sous-jacentes. Les $\mathcal{C}$-groupes for-
ment ainsi une catégorie $\mathcal{G}(\mathcal{C})$, munie de deux foncteurs fidèles et
transportables, à valeurs dans la catégorie $\mathcal{G}$ des groupes et dans $\mathcal{C}$
respectivement, et satisfaisant aux conditions envisagées au début du
n° 9. La notion de "$\mathcal{C}$-fibré en groupes" est donc définie, mais ce
n'est pas à elle qu'ici nous en avons.

Soit G un $\mathcal{C}$-groupe. Toujours conformément aux définitions générales
un $\mathcal{C}$-espace à $\mathcal{C}$-groupe d'opérateurs G est un ensemble E, muni des

\newpage

%% ===== Page 29 =====

% typescript page 27 — n° 262

structures de $\mathcal{C}$-espace et d'ensembles à groupe d'opérateurs $G$, com-
patibles dans le sens suivant : l'application $(g, x) \rightarrow g.x$ de $G \times E$ dans
$E$ est un $\mathcal{C}$-morphisme. Les morphismes pour de tels objets se défini-
sent comme ci-dessus en séparant les structures, soit qu'on garde $G$
fixe, soit qu'on le bouge. Dans chaque cas on obtient encore une caté-
gorie avec foncteurs à valeurs dans $\mathcal{C}$ resp. la catégorie des ensem-
bles à groupe d'opérateurs, satisfaisant toujours aux conditions du
début du n° 9, et permettant donc de définir la notion de "$\mathcal{C}$-fibré
à $\mathcal{C}$-groupe d'opérateurs" mais ce n'est pas à celle-là non plus que
nous en avons.

Supposons qu'on se soit donné un $\mathcal{C}$-espace $F$ à $\mathcal{C}$-groupe d'opé-
rateurs $G$, $G$ opérant fidèlement sur $F$. Soit $X$ un $\mathcal{C}$-espace, soit $g$ le
faisceau des germes morphismes de $X$ dans $G$, c'est un faisceau de grou-
pes, qui opère de façon évidente par germes d'automorphismes du $\mathcal{C}$-
espace fibré trivial $X \times F$.

\underline{Définition 6.} - Les notations sont celles qui précèdent. On appelle
$\mathcal{C}$-espace fibré à groupe structural $G$, fibré $F$ (muni de sa struc-
ture de $\mathcal{C}$-espace à opérateurs) et base $X$, un ensemble $E$ muni de la
structure d'$\mathcal{C}$-espace fibré ensembliste sur $X$, à faisceau structural $g$,
de type $(g, X \times F)$.

Cela signifie donc qu'on se donne une application $p$ de $E$ sur $X$,
et un faisceau           de germes d'isomorphismes de $X \times F$ sur $E$, qui
soit stable par composition à droite avec les germes de $g$, et qui
soit principal sous $G$. En vertu des propriétés, si les $\mathcal{C}$-fibrés
forment une catégorie locale, comme $g$ est un faisceau de germes
d'automorphismes du $\mathcal{C}$-fibré $X \times F$, un fibré $E$ à groupe structural $G$,
fibre $F$ est muni canoniquement d'une structure de $\mathcal{C}$-fibré ; il en

\newpage

%% ===== Page 30 =====

% typescript page 28 — n° 262

est ainsi en particulier si $\mathcal{C}$ est sursaturée, ou parfaite, ou semi-
parfaite (prop.13). Plus généralement, avec les notations du début du
n°9, supposons que $F \in \underline{\Phi}_0$, et que $G$ soit un groupe de $\underline{\Phi}$-automorphismes
de $F$, alors $g$ est un faisceau de germes d'automorphismes de $X \times F$ consi-
déré comme un fibré de type $(\mathcal{C}, \underline{\Phi}, \underline{\Phi}_0)$ ; ainsi, sous ces conditions,
un fibré à groupe structural $G$, fibre $F$ a une structure sous-jacente
de $\mathcal{C}$-fibré de type $(\mathcal{C}, \underline{\Phi}, \underline{\Phi}_0)$. Voici un cas important où on a une
réciproque :

PROPOSITION 14. - \emph{Les notations étant celles du début, du n° 9, et celles
précédant la définition 6, supposons que $F \in \underline{\Phi}_0$, que $G$ soit exactement
le groupe des $\underline{\Phi}$-automorphismes de $F$, et que la condition suivante
soit satisfaite : si $g$ est une application d'un $\mathcal{C}$-espace $U$ dans $G$,
telle que $(x,y) \to g(x).y$ et $(x,y) \to g(x)^{-1}.y$ soient des $\mathcal{C}$-morphismes
de $U \times F$ dans $F$, alors $g$ est un $\mathcal{C}$-morphisme de $U$ dans $G$. Supposons en-
fin $\mathcal{C}$ recollable. Alors pour tout ensemble $E$, la donnée sur $E$ d'une
structure de $\mathcal{C}$-espace fibré à groupe structural $G$, fibre $F$ est équiva-
lente à la donnée d'une structure de $\mathcal{C}$-espace fibré de type $(\mathcal{C}; \underline{\Phi}, $
$\underline{\Phi}_0)$ de fibre type $F$.}

Ceci résulte en effet aussitôt du corollaire 1 à la proposition 13.

On notera, que pour un $\mathcal{C}$-espace $F$ à $\mathcal{C}$-groupe fidèle $G$ d'opérateurs,
il existe toujours des catégories $\underline{\Phi}_0$, $\underline{\Phi}$ du type envisagé n° 9, et
une structure de type $\underline{\Phi}_0$ sur $F$, telle que $G$ soit exactement le groupe
des $\underline{\Phi}$-automorphismes de $F$ : Il suffit de considérer la catégorie $\underline{\Phi}$
des ensembles $Y$ munis à la fois d'une structure de "$G$-ensemble de type
$F$" donnée par un élément de l'ensemble quotient de l'ensemble des bi-
jections de $F$ sur $Y$ par les opérations de $G$), et la catégorie $\underline{\Phi}_0$ des
ensembles munis à la fois d'une structure $\underline{\Phi}$ et d'une structure de

\newpage

%% ===== Page 31 =====

% typescript page 29 — n° 262

$\mathcal{C}$-espace, soumis à la condition que les bijections "permises" de $F$ dans $Y$ soient des $\mathcal{C}$-isomorphismes. Cela montre que, pour peu que la condition, supplémentaire sur $G,F$ envisagée dans la prop. 14 est satisfaite, la notion de fibré à groupe structural se ramène (avantageusement) aux notions du n° 9.

Prenons par exemple $F=G$, $G$ opérant sur lui-même par \emph{translations gauches}. Si $g$ est une application de $U$ dans $G$ telle que l'application $(x,y) \to g(x).y$ de $U \times G$ dans $G$ soit un morphisme, on en conclut, en composant avec l'application $x \to (x,e)$ de $U$ dans $U \times G$, qui est un morphisme en vertu du corollaire à prop. 12, que $g$ est bien un morphisme de $U$ dans $G$. Prenons pour $\underline{\Phi}_0$ la catégorie des $\mathcal{C}$-espaces à $\mathcal{C}$-groupe d'opérateurs à droite, $F$, pour $\underline{\Phi}$ la catégorie des ensembles à groupe d'opérateur à droite $G$, et considérons $F=G$ comme un élément de cette catégorie, en faisant opérer $G$ sur $F$ par \emph{translations à droite}. Les $\mathcal{C}$-fibrés de type $(\mathcal{C}, \underline{\Phi}, \underline{\Phi}_0)$, i.e. les $\mathcal{C}$-espaces fibrés à $\mathcal{C}$-groupe d'opérateurs $G$ à droite, qui admettent la fibre type qu'on vient de définir, sont appelés les \emph{$\mathcal{C}$-espaces principaux} (à droite) de \emph{groupe} $G$. On déduit donc de ce qui précède :

\emph{Corollaire : Les $\mathcal{C}$-espaces fibrés principaux de groupe $G$ s'identifient aux $\mathcal{C}$-espaces fibrés à groupe structural $G$, fibre $G$, où le groupe $G$ opère par translations à gauche. (On suppose $\mathcal{C}$ recollable).}

En vertu des faits généraux bien connus (Top. Elém. chap. I) la catégorie prélocale des $\mathcal{C}$-espaces fibrés à groupe structural $G$ fibre $F$ sur des ouverts de $X$, i.e. des fibrés ensemblistes sur des ouverts $U$ de $X$, à faisceau structural $g$, de type $(g, U \times F)$, ne dépend, à "isomorphisme canonique" près, que du faisceau $g$ et non du fibré ensembliste $X \times F$ où on fait opérer ce faisceau, donc ne dépend que de $G$ et non de la fibre

\newpage

%% ===== Page 32 =====

% typescript page 30 — n° 262

$F$. Ceci montre donc \emph{à priori} la correspondance entre $\mathcal{C}$-fibrés à grou-
pe structural et fibre $F$, et $\mathcal{C}$-fibrés \emph{principaux} (de groupe $G$) asso-
ciés. On en conclut, compte tenu de prop. 13 :

\emph{Corollaire 2} - Si $\mathcal{C}$ est recollable, alors la catégorie prélocale des
$\mathcal{C}$-fibrés sur $X$ à groupe structural $G$, fibre $F$, est une catégorie lo-
cale.

On a oublié, dans les remarques précédant le corollaire 2, et dans
ledit corollaire, de préciser que les morphismes d'un $\mathcal{C}$-fibré $E$ à groupe
structural $G$, fibre $F$ dans un autre (de même groupe structural, de même
fibre et de même base, sont ici par définition les isomorphismes de fi-
brés ensemblistes à faisceau structural $g$, de type $(g, X \times F)$. Prenant
comme morphismes dans $\Phi$ les isomorphismes d'ensembles à groupe d'opé-
rateurs $G$ fixé (applications "compatibles" avec les opérations de $G$),
ce qui détermine aussi la notion de morphismes dans $\Phi_0$, on voit que la
correspondance envisagée dans prop. 14 est fonctorielle. Ceci montre en
plus clairement que la notion de $\mathcal{C}$-espace fibré à groupe structural
n'est pas très commode telle quelle quand on veut avoir des notions
plus variées de morphismes de fibrés ; ainsi la notion naturelle de
morphismes de $\mathcal{C}$-fibrés vectoriels (cf par. 2) ne s'exprime pas en
termes de groupes structuraux. Aussi le langage du n° 9 sera-t-il le
plus souvent préférable à celui développé ici.

Nous ne reviendrons pas ici sur la construction "géométrique" d'un
$\mathcal{C}$-fibré à groupe structural $G$ et fibre $F$, associé à un fibré princi-
pal et vice-versa, espérant qu'elle découlera automatiquement des trucs
généraux exposés dans Top. Elém. Chap.I

Enfin il faut encore préciser la version globale de la notion de

\newpage

%% ===== Page 33 =====

% typescript page 31 — n° 262

$\mathcal{C}$-fibré principal à $\mathcal{C}$-groupe structural $G$. Soit d'abord $E$ un $\mathcal{C}$-
fibré principal à $\mathcal{C}$-groupe structural $G$, de base $X$ ; alors $E$ admet
localement une $\mathcal{C}$-section (nº 7, déf. 3) : on peut en effet supposer
que $E$ est de la forme $X \times G$, et on prend la section définie par l'appli-
cation constante $x \to e$ de $X$ dans $G$, qui est une $\mathcal{C}$-section en ver-
tu du corollaire à prop. 12. Donc (nº 7, prop. 11) $X$ s'identifie à un $\mathcal{C}$-
espace quotient de $E$. Par suite, \emph{le $\mathcal{C}$-espace fibré principal $E$ est
connu, à isomorphisme près, quand on connait la structure sous-jacente
de $\mathcal{C}$-espace à $\mathcal{C}$-groupe d'opérateurs à droite.}

\underline{Définition 7} - Un $\mathcal{C}$-espace $E$ à $\mathcal{C}$-groupe d'opérateurs $G$ à droite
est dit \emph{localement trivial} (ou encore on dit que $X$ est \emph{fibré locale-
ment trivial par $G$}) s'il est isomorphe au $\mathcal{C}$-espace à $\mathcal{C}$-groupe
d'opérateurs défini par un $\mathcal{C}$-espace fibré principal $E_1$ à $\mathcal{C}$-groupe
structural $G$.

S'il en est ainsi, on peut supposer alors que $E_1 = E$ en tant que
$\mathcal{C}$-espace à opérateurs, et que la base de $E_1$ est $\mathcal{C}$-espace quotient
$X = E/G$.

PROPOSITION 15 - \emph{Soit $E$ un $\mathcal{C}$-espace à $\mathcal{C}$-groupe d'opérateurs $G$ à
droite. Pour que $E$ soit isomorphe à un $\mathcal{C}$-espace fibré principal tri-
vial, il faut et il suffit qu'il existe un $\mathcal{C}$-morphisme $f$ de $E$ dans $G$
ayant les propriétés suivantes : (i) $f(xg) = f(x).g$ pour $x \in E$, $g \in G$
(ii) l'ensemble $X$ des $x \in E$ tels que $f(x) = e$ admet une $\mathcal{C}$-structure
induite.}

La condition est nécessaire, car si $E = X \times G$, on pose $f(x,g) = g$ ;
la condition (i) est vérifiée trivialement, et la condition (ii) résul-
te du corollaire à prop. 12. La condition est suffisante, car considé-
rons les applications $(x,g) \to x.g$ de $X \times G$ dans $E$, et $y \to (y.f(y)^{-1},
f(y))$ de $E$ dans $X \times G$, ce sont là des $\mathcal{C}$-morphismes, compatibles avec

\newpage

%% ===== Page 34 =====

% typescript page 32 - n° 262

les opérations de G, et on constate aussitôt qu'elles sont inverses
l'une de l'autre.

\underline{Corollaire 1} - Soit E un $\mathcal{C}$-espace à $\mathcal{C}$-groupe d'opérateurs G à droite
Supposons $\mathcal{C}$ \emph{sursaturée}. Pour que E soit fibré localement trivial par
G (déf.7) il faut et il suffit que pour tout $x \in E$ existe un voisinage
ouvert U stable par G, et un $\mathcal{C}$-morphisme f de U dans G satisfaisant
les conditions (i) et (ii) de la proposition 15.

\underline{Corollaire 2} - Soit E un $\mathcal{C}$-espace à $\mathcal{C}$-groupe d'opérateurs G à
droite. Supposons $\mathcal{C}$ \emph{parfaite} ou \emph{semi-parfaite}. Pour que E soit fibré
localement trivial par G (déf. 7), il faut et il suffit que la rela-
tion d'équivalence définie par G ait un graphe fermé, et que pour tout
$x \in E$ existe un voisinage ouvert U de $x$ et un $\mathcal{C}$-morphisme f de U dans
G satisfaisant les conditions (i) et (ii) de prop. 15.

Démontrons les corollaires : la nécessité étant évidente dans cha-
que cas grâce à prop.15, il suffit de démontrer la suffisance des con-
ditions envisagées. Il résulte des hypothèses et de prop.15 que tout
$x \in E$ admet un voisinage ouvert U stable par G, tel que $U/G$ soit un $\mathcal{C}$-
espace quotient de U ; il en résulte ($n^{\circ} 7$, prop.9) que $E/G$ est un $\mathcal{C}$-
espace quotient de E. C'est donc un $\mathcal{C}$-espace quotient si $\mathcal{C}$ est sur-
saturée. Dans le cas où $\mathcal{C}$ est parfaite ou semi-parfaite, montrons que
$X = E/G$ est séparé, d'où résultera aussitôt que X est un $\mathcal{C}$-espace
quotient de E. Or il résulte de prop.11, corollaire que $X \times X$ est un $\mathcal{C}$-
espace quotient de $E \times E$, en particulier sa topologie est la topologie
quotient de celle de $E \times E$. Comme l'image inverse de la diagonale D de
$X \times X$ est le graphe de la relation d'équivalence définie par G dans E,
donc fermée, il s'ensuit que D est fermée, donc X est séparé. Donc en
tous cas $X = E/G$ est un $\mathcal{C}$-espace quotient de E.

\newpage

%% ===== Page 35 =====

% typescript page 33 — n° 262

Munissons alors $E$ de sa $\mathcal{C}$-structure, de sa projection sur le $\mathcal{C}$-
espace $X = E/G$ (projection qui est un $\mathcal{C}$-morphisme) et des structures
d'ensembles à groupe d'opérateurs $G$ sur les fibres, définies par les
opérations de $G$ dans $E$, je dis qu'alors $E$ devient un $\mathcal{C}$-fibré au sens
du n$^\circ$ 9, déf. 4. Cela résulte en effet de prop. 15 appliquée aux ouverts
$U$ envisagés dans les corollaires 1 et 2.

Remarque - Nous verrons que dans les cas les plus importants, la con-
dition (ii) de la proposition 15 résulte déjà de la condition (i).
C'est le cas évidemment si $\mathcal{C}$ est la catégorie topologique des espaces
topologiques, puisque toute partie d'un espace topologique admet une
structure induite. C'est aussi le cas si $\mathcal{C}$ est la catégorie des
ensembles algébriques définis sur un corps $k$, car toute partie fermée
d'un $\mathcal{C}$-espace admet alors une structure induite. C'est enfin aussi
le cas si $\mathcal{C}$ est la catégorie des variétés $m$-fois différentiables, ou
analytiques réelles, ou analytiques complexes, comme nous le verrons
ultérieurement.

11. - \underline{$\mathcal{C}$-fibrés associés}.

Revenons à la situation générale du début du n$^\circ$ 9, symbolisé par un
diagramme de foncteurs fidèles et transportables.

\begin{tikzcd}
& \Phi_0 \arrow[dl] \arrow[dr] & \\
\mathcal{C} \arrow[dr] & & \mathcal{E} \arrow[dl] \\
& \mathcal{E}' &
\end{tikzcd}

et supposons un diagramme analogue faisant intervenir des catégories
$\underline{\Phi}$, $\underline{\Phi}_0$ :

\begin{tikzcd}
& \underline{\Phi}_0 \arrow[dl] \arrow[dr] & \\
\underline{\Phi} \arrow[dr] & & \mathcal{E} \arrow[dl] \\
& \mathcal{E}' &
\end{tikzcd}

\newpage

%% ===== Page 36 =====

% typescript page 34 — n° 262

Supposons donnés enfin des foncteurs, covariants pour fixer les idées
(mais rien ne serait essentiellement changé s'ils étaient tous deux
contravariants)
\[ F : \underline{\Phi} \to \Psi \quad , \quad F_0 : \underline{\Phi}_0 \to \Psi_0 \]
donnant lieu aux trois relations de commutation naturelles relativement
aux foncteurs intervenant dans les diagrammes ci-dessus. Supposons que
$\mathcal{C}$ soit recollable. Nous allons alors définir un foncteur, dénoté de
façon abrégée par $F'$, de la catégorie des fibrés de type $(\mathcal{C}, \underline{\Phi}, \underline{\Phi}_0)$
dans la catégorie des fibrés de type $(\mathcal{C}, \Psi, \Psi_0)$, à un espace fibré $\xi$
étant associé un espace fibré de même base, dit \underline{associé} à $\xi$ et aux fonc
teurs $(F,F_0)$. (Dans le cas où $\underline{\Phi} = \underline{\Phi}_0$, on a $F = F_0$, et on ne parle plus
que de $F$). Soit $\xi = (E,p,X)$ le fibré envisagé, pour tout $x \in X$ la fibre
de $F'(\xi)$ est par définition $F(\xi_x)$, où $\xi_x = p^{-1}(x) \in \underline{\Phi}_0$ est la fibre de $\xi$
sur $x$ : on prendra donc pour $F'(\xi)$ l'ensemble somme des $\xi'_x = F(\xi_x)$,
pour projection $p' : E' \to X$ l'application dont les fibres sont les $\xi'_x$,
enfin on définit sur $E'$ une structure de fibré de type $(\mathcal{C}, \Psi, \Psi_0)$
par recollement au moyen de la famille suivante de structures de ce
type sur des $p'^{-1}(U_i)$ : pour tout isomorphisme $f_i : p^{-1}(U_i) \to U_i \times A_i$ (où
$U_i$ est un ouvert de $X$, et $A_i \in \underline{\Phi}_0$), on déduit une famille d'isomorphis
mes $\xi_x \to A_i$ dans $\underline{\Phi} \; (x \in U_i)$, d'où en appliquant le foncteur $F$ une
famille d'isomorphismes $\xi'_x \to F(A_i)$ dans $\Psi$, d'où une bijection
$f'_i : p'^{-1}(U_i) \to U_i \times F(A_i)$, compatible avec les applications de projec
tion sur $U_i$ et les structures de type $\Psi$ sur les fibres. Il faut véri
fier que les cartes $f'_i$ ainsi obtenues sont compatibles, pour ceci,
il faut imposer au foncteur $F$ la condition suivante :

\underline{Soient $U \in \mathcal{C}$ et $A,B \in \underline{\Phi}_0$, enfin soit $g : U \to \underline{\mathrm{Hom}}_{\underline{\Phi}}(A,B)$ une appli-}
\underline{cation telle que l'application $(x,a) \to g(x)a$ de $U \times A$ dans $B$ soit un}

\newpage

%% ===== Page 37 =====

% typescript page 35 - n° 262

$\mathcal{C}$-morphisme, alors l'application $x \longrightarrow g'(x) = F(g(x))$ de $U$ dans
$\mathrm{Hom}_{\Psi}(F(A),F(B))$ \emph{a la propriété analogue}.

Supposant cette condition satisfaite, les cartes envisagées sur les
${p'}^{-1}(U_i)$ se recollent, et définissent donc en vertu de prop.13 une
structure de fibré de type $(\mathcal{C}, \Psi, \Psi_0)$ sur $E'$, et par définition
$F'(\xi)$ désigne $E'$ muni de la structure en question.

Soient $\xi = (E,p,X)$ et $\xi_1 = (E_1,p_1,X_1)$ deux fibrés de type $(\mathcal{C}, \underline{\Phi}, \underline{\Phi}_0)$
et $f$ un morphisme du premier dans le second, $f$ est donc défini par un
$\mathcal{C}$-morphisme $f_X: X \to X_1$, et par une application $x \longrightarrow f(x)$, où pour
tout $x \in X$, $f(x)$ désigne un élément de $\mathrm{Hom}_{\underline{\Phi}}(\xi_x, \xi_{1,f_X(x)})$. Posons
$$\xi' = F'(\xi), \quad \xi'_1 = F'(\xi_1), \quad f'(x)=F(f(x)) \in \mathrm{Hom}_{\Psi}(\xi'_x, \xi'_{1,f_X(x)})$$

Utilisant la condition imposée ci-dessus à $F$, on constate que l'appli-
cation $f': E' \longrightarrow E'_1$ définie par $f_X$ et les $f'(x)$ définis par la formule
précédente, est un $\mathcal{C}$-morphisme, donc en fait un morphisme de fibrés de
type $(\mathcal{C}, \Psi, \Psi_0)$. On le désigne par $F'(f)$. On constate que les fonc-
tions $F'(\xi)$ et $F'(f)$ définissent bien un foncteur comme annoncé i.e.
que $F'(\text{identité})=\text{identité}$, $F'(gf)=F'(g)F'(f)$, ce qui achève la définition
que nous avions en vue.

Parmi les propriétés imprtantes de la notion de fibré associé, il
faut signaler la \emph{transitivité des fibrés associés}, exprimée par la for-
mule fonctorielle,
$$(GF)' = G'F'$$
(où la variance de $G$ et $F$ est ce qu'elle veut). De plus, $F'$ \emph{est un
foncteur par rapport au couple} $(F,F_0)$, i.e. si on a un couple analogue
$(G,G_0)$ de foncteurs $\underline{\Phi} \to \Psi$ et $\underline{\Phi}_0 \to \Psi_0$, et un homomorphisme foncto-
riel $u: F \to G$ tel que, pour $A \in \underline{\Phi}_0$, l'homomorphisme $u(A): F(A) \to G(A)$ soit
en fait un $\Psi_0$-homomorphisme de $F_0(A)$ dans $G_0(A)$, il correspond à $u$

\newpage

%% ===== Page 38 =====

% typescript page 36 - n° 262

$$u' : F' \longrightarrow G'$$

dont la définition, évidente, est laissée à un lecteur fictif, (que,
suivant l'usage, nous supposons infiniment scrupuleux). Ce même lec-
teur vérifiera aussi les importantes formules

$$(\mathrm{identité})' = \mathrm{identité}, \quad (vu)' = v'u'$$

qui justifient notre assertion que $F'$ est un foncteur en $(F, F_o)$.

On a oublié de dire que les foncteurs construits dans ce numéro,
sont des foncteurs de catégories locales sur $\mathcal{C}$ (Top. Elem. Chap I);
et les homomorphismes fonctoriels $u'$ sont des homomorphismes de fonc-
teurs de catégories locales.

En particulier, soit $G$ un $\mathcal{C}$-groupe, et désignons comme au n°10,
par $\Phi_o$ la catégorie des $\mathcal{C}$-espaces à $\mathcal{C}$-groupe d'opérateurs \emph{à droite}
$G$ isomorphes à $G$ (opérant sur lui-même par la représentation régulière
droite), les morphismes étant les isomorphismes, et par $\Phi$ la catégo-
rie des ensembles à groupe d'opérateurs $G$ à droite, les morphismes
étant les isomorphismes. Soient $\Psi_o$, $\Psi$ comme au début de ce numéro.
Notons que $\mathrm{Hom}_{\Phi}(G,G)$ est isomorphe à $G$, opérant sur lui-même par la
représentation \emph{régulière} gauche. Donc un couple de foncteurs $(F, F_o)$
satisfaisant aux conditions envisagées définit un objet $A \in \Psi_o$, savoir
$A = F_o(G)$, et des opérations à gauche de $G$ par $\Psi$-automorphismes de $A$
(Transformées par $F$ des $\Phi$-automorphismes de $G$), et la condition sou-
lignée sur $F$ (appliquée au cas où $U=G$, et où $g : G \times G \to G$ est définie
par la loi de groupe de $G$) implique que en fait $A = F_o(G)$ est un $\mathcal{C}$-es-
pace à opérateurs à gauche. Inversement, la donnée d'un $A \in \Psi_o$, avec
une structure de $\mathcal{C}$-espace à groupe d'opérateurs $G$ \emph{à gauche}, les
opérations de $G$ étant des $\Psi$-automorphismes, définit de façon bien
connue un foncteur $\Phi \to \Psi$, et de même un foncteur $\Phi_o \to \Psi_o$ noté

\newpage

%% ===== Page 39 =====

% typescript page 37 — n° 262

$\underline{E} \to (\underline{E},A)_G$ (en tant qu'ensemble, $(\underline{E},A)_G$ est le quotient de $\underline{E} \times A$ par le
groupe $G$ opérant par $g.(e,a) = (e.g^{-1},g.a)$) ; ces foncteurs satisfont
aux trois relations de compatibilité qu'il faut, et de plus à la condi-
tion soulignée , comme on vérifie aisément. Ainsi, \emph{un couple de fonc-
teurs $(F,F_0)$ ayant les propriétés envisagées dans ce numéro, s'identifie
aussi à un objet de $\underline{\Psi}_0$ qui est en même temps un $\mathcal{C}$-espace à $\mathcal{C}$-groupe
d'opérateurs $G$, les opérations de $G$ étant des $\underline{\Psi}$-automorphismes}. Cela
précise du point de vue actuel la notion de fibré associé à un $\mathcal{C}$-fibré
principal et une représentation de $G$ par des automorphismes d'un ensem-
ble-fibre $A$. On notera aussi que dans l'interprétation précédente de
foncteurs $(F,F_0)$ par des $\mathcal{C}$-espaces à opérateurs, les homomorphismes
fonctoriel (permettant, comme on l'a dit, de construire des homomorphis-
mes fonctoriels pour les fibrés associés à un même fibré) correspondent
aux morphismes $A \to B$, étant entendu que morphisme signifie morphisme
pour la structure de $\mathcal{C}$-espaces à groupe d'opérateurs $G$, \underline{et} pour la
structure de type $\underline{\Psi}_0$.

Un cas intéressant est celui où on prend pour $\underline{\Psi}_0$ la catégorie des
$\mathcal{C}$-groupes, pour $\underline{\Psi}$ la catégorie des groupes, et pour $A$ le groupe
$G$ lui-même (donc $A \in \underline{\Psi}_0$) opérant sur lui-même par représentation inté-
rieure : $int(g)g' = gg'g^{-1}$. On trouve donc, pour tout $\mathcal{C}$-fibré prin-
cipal $P$ à groupe structural $G$, un $\mathcal{C}$-fibré en groupes tous isomorphes
à $G$, dont la fibre en $x \in X$ s'identifie d'ailleurs au groupe des permu-
tations de la fibre $P_x$ qui permutent aux opérations de $G$ ; ce fibré en
groupes est appelé le \emph{fibré adjoint}. Il est défini chaque fois qu'on a
un $\mathcal{C}$-fibré à groupe structural $G$ (non nécessairement principal, car
on peut commencer par passer au fibré principal associé). Ses $\mathcal{C}$-sections
sur $X$ s'identifient aux automorphismes de $\underline{E}$.

\newpage

%% ===== Page 40 =====

% - 38 - n° 262

Bien entendu, le passage d'un fibré à groupe structural au fibré
principal associé rentre aussi dans le cadre exposé dans ce numéro,
pourvu que l'on soit sous les conditions de prop.14, permettant de
ramener la notion de fibré à groupes structural aux notions du $n^\circ$ 9.

11. bis - \emph{Extension et restriction du groupe structural.}

Dans tout ce numéro, $\mathcal{C}$ désigne une catégorie locale recollable.
Soit $G$ un $\mathcal{C}$-groupe , $F$ un sous-$\mathcal{C}$-groupe fermé (i.e. un sous-ensemble
fermé, muni d'une $\mathcal{C}$-structure induite, et qui est en même temps un
sous-groupe, ce qui implique que $F$ est lui-même un $\mathcal{C}$-groupe). Soit
$\underline{\Phi}_0(G)$ la catégorie des $\mathcal{C}$-espaces à $\mathcal{C}$-groupe d'opérateurs $G$ à droi-
te, $\underline{\Phi}(G)$ la catégorie des ensembles à groupe d'opérateurs $G$ à droite.
On a $G \in \underline{\Phi}_0(G)$, et d'autre part $G$, opérant sur $G \in \underline{\Phi}(G)$ par translations
à gauche, s'identifie à $\mathrm{Aut}_{\underline{\Phi}}(G)$. Donc l'injection $i:F \to G$ définit un
homomorphisme $F \to \mathrm{Aut}_{\underline{\Phi}}(G)$. Donc cela permet, en vertu de $n^\circ 10$, pour
tout $\mathcal{C}$-fibré principal $P$ à $\mathcal{C}$-groupe structural $F$, de construire le
fibré principal $P'$ à groupe structural $G$ associé à $P$ et $i$ ; on dit que
$P'$ est \emph{déduit de $P$ par extension du $\mathcal{C}$-groupe structural $F$ à $G$}. On va
définir un morphisme fonctoriel

(1) $P \to P'$

compatible avec les opérations de $F$ sur les deux membres (et bien sur
avec les restrictions aux parties ouvertes de $X$). Pour ceci, soit $\underline{\Phi}_0$ la
catégorie des $\mathcal{C}$-espaces à $\mathcal{C}$-groupes d'opérateurs à droite (groupes
non précisés), $\underline{\Phi}$ la catégorie des ensembles à groupes d'opérateurs à
droite. Un morphisme $(E,G) \to (E',G')$ dans $\underline{\Phi}_0$ est un couple $(f,u)$
d'un $\mathcal{C}$-morphisme $f:E \to E'$ et d'un morphisme $u:G \to G'$ de $\mathcal{C}$-groupes
tel que $f(x.g)=f(x).u(g)$ pour $x \in E$, $g \in G$, et on définit de façon ana-
logue les morphismes dans $\underline{\Phi}$. Alors les deux membres de (1) peuvent être

\newpage

%% ===== Page 41 =====

% typescript page 39 — n° 262

considérés comme des fibrés de type $(\mathcal{C}, \underline{\Phi}, \underline{\Phi}_0)$ associés à $P$, et la
représentation naturelle de $F$ dans $\mathrm{Aut}_{\underline{\Phi}}(F)$ resp. dans $\mathrm{Aut}_{\underline{\Phi}}(G)$ par
translations à gauche (où $F$ et $G$ sont considérés comme des éléments de
$\underline{\Phi}_0$, leurs groupes d'opérateurs étant $F$ resp. $G$). Le couple $(1,i)$, où
$1$ est l'application identique $F \to F$, est un morphisme $F \to G$ dans $\underline{\Phi}_0$
compatible avec les automorphismes définis par les éléments de $F$. En ver-
tu des sorites du n$^\circ$ 11, ce morphisme définit donc bien un morphisme (1)
pour les fibrés associés, succès brillant de l'axiomatique. Comme $i$ est
injectif, il en est de même de (1), d'après un des nombreux sorites
oubliés. En fait, $F$ étant un sous-ensemble fermé du $\mathcal{C}$-espace $G$, admet-
tant une $\mathcal{C}$-structure induite, il en est de même de $P$ considéré comme
sous-ensemble de $P'$ (re-sorite oublié, s'appuyant sur prop.7)

Inversement, partons d'un $\mathcal{C}$-fibré principal à $\mathcal{C}$-groupe structural $G$.
Définition 7 bis - Soit $P_1$ un $\mathcal{C}$-fibré principal à $\mathcal{C}$-groupe structu-
ral $G$. On appelle \emph{restriction du groupe structural à $F$} tout $\mathcal{C}$-fibré à
$\mathcal{C}$-groupe structural $F$, satisfaisant aux conditions suivantes : $P$ est
contenu dans $P_1$, et il existe un isomorphisme sur $P_1$ du $\mathcal{C}$-fibré princi
pal $P'$ déduit de $P$ par extension du groupe structural à $G$, se réduisant
à l'identité sur $P$.

Une telle restriction existe donc si et seulement si $P_1$ est isomor-
phe à un $\mathcal{C}$-espace fibré principal obtenu par extension à $G$ du groupe
structural d'un $\mathcal{C}$-espace fibré principal à $\mathcal{C}$-groupe structural $F$.
Ceci peut encore s'exprimer de la façon suivante. Soit $\underline{O}_X(G)$ le fais-
ceau de groupes des germes de $\mathcal{C}$-morphismes de $X$ dans $G$, définissons de
même $\underline{O}_X(F)$. L'homomorphisme d'injection $i:F \to G$ définit un homomorphisme
de faisceaux $\underline{O}_X(F) \to \underline{O}_X(G)$, d'où une application
$$ H^1(X, \underline{O}_X(F)) \to H^1(X, \underline{O}_X(G)) $$

\newpage

%% ===== Page 42 =====

% typescript page 40 — n° 262

Il est immédiat que, pour que $P_1$ admette une restriction du groupe
structural à $F$, il faut et il suffit que sa classe $c(P_1) \in H^1(X, \underline{O}_X(G))$
soit dans l'image de $H^1(X, \underline{O}_X(F))$ (compte tenu d'une compatibilité géné-
rale entre fibrés associés et images de classes de cohomologie, que
nous avons oubliée au numéro précédent).

Nous allons donner une interprétation plus géométrique de la no-
tion de restrictions du groupe structural. Notons d'abord qu'une telle
restriction est bien déterminée quand on connait la partie $P$ de $P_1$
qu'elle définit. En effet, il résulte de ce qu'on a vu que sa $\mathcal{C}$-
structure est la $\mathcal{C}$-structure induite par $P_1$, sa projection sur $X$ est
induite par celle de $P_1$, et les opérations de $F$ sur $P$ sont induites pa
par les opérations de $F$ dans $P_1$. D'ailleurs, $P$ est l'image réciproque,
par l'application canonique $P_1 \to P_1/F$, d'une partie $s(X)$, où $s$ est une
section de $P_1/F$ (considéré comme fibré ensembliste sur $X$) sur $X$, bien
déterminée par $P$. On est ainsi amené à examiner quelles sont les sec-
tions de $P_1/F$ qui correspondent à une restriction de groupe structu-
ral. Je dis que \emph{ce sont exactement celles qui, localement, sont images
de $\mathcal{C}$-sections de $P_1$}. La nécessite est évidente, puisque on peut sup-
poser alors que $P_1$ est le fibré $P'$ déduit d'un fibré $P = X \times F$ par exten-
sion du groupe structural : alors $P_1 = X \times G$, et la section $s$ est l'image
canonique de la section de $P_1$ définie par l'application constante
$x \to e$ de $X$ dans $G$, section qui est un $\mathcal{C}$-morphisme. Pour la suffisan-
ce, supposons d'abord que $s$ provienne d'une section régulière $s_1$ de
$P_1$, cette section définit un isomorphisme de $P_1$ avec $X \times G$, et identi-
fiant $P_1$ à $X \times G$, on voit que $s$ correspond à la partie $X \times F$ de $X \times G$, qui
est bien une restriction du groupe structural. Si on suppose seule-
ment que localement, $s$ provient d'une $\mathcal{C}$-section de $P_1$, il en résulte

\newpage

%% ===== Page 43 =====

% typescript page 41 — n° 262

d'abord que la partie $P$ de $P_1$ correspondante à $s$ admet une $\mathcal{C}$-structure
induite (car son intersection avec les images inverses des ouverts $U$ de
$X$ assez petits ont une structure induite, de sorte qu'on peut appliquer
la prop.7). Munissons alors $P$ de cette $\mathcal{C}$-structure induite, de la pro-
jection sur $X$ induite par celle de $P_1$, enfin des opérations de $F$ sur les
fibres induites par les opérations de $F$ dans $P_1$ : il est immédiat que $P$
devient alors un $\mathcal{C}$-espace fibré principal à $\mathcal{C}$-groupe structural $F$,
i.e. satisfait la condition LT de déf.4. On prouve aussi facilement que
le $\mathcal{C}$-fibré $P'$ déduit de $P$ par extension du groupe structural à $G$ s'i-
dentifie bien à $P_1$ ; on définit une bijection canonique $P' \to P_1$, compa-
tible avec les projections sur $X$ et les opérations de $G$, c'est d'ailleurs
un $\mathcal{C}$-isomorphisme au-dessus des ouverts $U$ de $X$ assez petits, c'est donc
un $\mathcal{C}$-isomorphisme, donc un isomorphisme pour les structures de $\mathcal{C}$-fibré
principaux. Nous avons prouvé :

PROPOSITION 15 bis - Supposons $\mathcal{C}$ recollable, et soit $P_1$ un $\mathcal{C}$-fibré
principal à $\mathcal{C}$-groupe structural $G$. Les restrictions du groupe structu-
ral de $P_1$ à $F$ correspondent biunivoquement aux sections du fibré ensem-
bliste $P_1/F$ sur $X$ qui localement sont des images de sections $\mathcal{C}$-réguliè-
res de $P_1$.

Supposons maintenant que le couple $(F,G)$ satisfasse à la condition
suivante :

(2) Pour tout $Y \in \mathcal{C}$, $(Y \times G)/F$ est un $\mathcal{C}$-espace quotient de $Y \times G$ ($n^\circ 7$,
déf.2) et s'identifie au produit $Y \times (G/F)$.

Ce produit a un sens, puisque $G/F$ est un $\mathcal{C}$-espace quotient de $G$, comme
on voit en prenant $Y$ réduit à un point).

\emph{Corollaire 1.} Soient $G$ un $\mathcal{C}$-groupe, $F$ un sous-$\mathcal{C}$-groupe satisfaisant
à la condition (2) ci-dessus, $P_1$ un $\mathcal{C}$-fibré principal à $\mathcal{C}$-groupe

\newpage

%% ===== Page 44 =====

% typescript page 42 - n° 262
structural $G$, de base $X$. Alors $P_1/F$ est un $\mathcal{C}$-espace quotient de $P_1$,
et muni de la projection $P_1/F \to X$ c'est un $\mathcal{C}$-espace fibré à fibre-
type $G/F$. Les sections de $P_1/F$ associées aux restrictions du groupe
structural de $P_1$ à $F$ sont des $\mathcal{C}$-sections.

Soit $p$ la projection de $P_1$ sur $X$. Il résulte immédiatement de (2)
que pour tout ouvert $U \subset X$ assez petit, $p^{-1}(U)/F$ est un $\mathcal{C}$-espace quo-
tient de $p^{-1}(U)$, isomorphe à $U \times (G/F)$ (isomorphisme compatible avec les
projections naturelles sur $U$). En vertu de prop. 9 il en résulte que
$Q = P_1/F$ est un $\mathcal{C}$-espace quotient de $P_1$. Il faut donc montrer que si
$\mathcal{C}$ est parfaite ou semi-parfaite, $Q$ est un $\mathcal{C}$-espace local. séparé. Or
on montre que $Q \times Q$ est un $\mathcal{C}$-espace quotient de $P_1 \times P_1$ par $F \times F$, tout
d'abord si $P_1 = X \times G$ en utilisant (2) deux fois, et tenant compte de la
transitivité des $\mathcal{C}$-espaces quotients, puis dans le cas général en
utilisant prop. 9. En particulier la topologie de $Q \times Q$ est la topologie
quotient de celle de $P_1 \times P_1$, et pour prouver que la diagonale $\Delta$ de $Q \times Q$
est fermée, il suffit de prouver que le graphe $B$ de la relation d'équi-
valence définie par $F$ dans $P_1$ est une partie fermée de $P_1 \times P_1$. Or si $A$
est le graphe de la relation d'équivalence définie par $G$ dans $P_1$, et si
$u$ est l'application de $A$ dans $G$ telle que $(x,y) \in A$ implique $y = x.u(x,y)$,
$B$ est l'image réciproque de $F$ par l'application $u$. Notre assertion
résulte alors du lemme suivant, dont la démonstration est facile et
laissée au lecteur :

Lemme 1. - Soit $P_1$ un $\mathcal{C}$-fibré principal à $\mathcal{C}$-groupe structural $G$,
base $X$. Soient $A$ le graphe de la relation d'équivalence définie par $G$
dans $P_1$, $u$ l'application de $A$ dans $G$ telle que $(x,y) \in A$ implique
$y = x.u(x,y)$. Alors $u$ est un $\mathcal{C}$-morphisme, et pour que $X$ soit séparé, il
faut et il suffit que $A$ soit fermé.

\newpage

%% ===== Page 45 =====

% typescript page 43 — n° 262

Cela achève de prouver que, sous les conditions du corollaire 1,
$P_1/F$ est un $\mathcal{C}$-espace quotient de $P_1$. La démonstration montre en même
temps que c'est un $\mathcal{C}$-espace fibré sur $X$ (et même un $\mathcal{C}$-espace fibré
à $\mathcal{C}$-groupe structural $G$, pourvu que $G$ opère fidèlement dans $G/F$). Le
fait que les sections de ce fibré associées aux restrictions du groupe
structural sont des $\mathcal{C}$-sections résulte alors immédiatement de prop.
15 bis.

\underline{Corollaire 2} - Soient $G$ un $\mathcal{C}$-groupe , $F$ un sous-$\mathcal{C}$-groupe tel que $G$
fibré par $F$ (opérant à droite) soit localement trivial (n° 10 définition
7). Soit $P_1$ un $\mathcal{C}$-fibré principal à $\mathcal{C}$-groupe structural $G$. Alors
$P_1/F$ est un $\mathcal{C}$-espace quotient de $P_1$ et un $\mathcal{C}$-espace fibré sur $X$, à
fibre-type $G/F$, dont les $\mathcal{C}$-sections correspondent biunivoquement aux
restrictions du groupe structural $G$ de $P_1$ à $F$. \underline{Tout d'abord}, l'hypothè-
se faite sur le couple $(F,G)$ implique l'hypothèse (2) du corollaire 1,
en vertu de n° 7, cor. à prop. 11. On peut donc appliquer le dit coro-
laire 1, de sorte qu'il reste seulement à montrer que \underline{toute} $\mathcal{C}$-section
de $P_1/F$ sur $X$ correspond à une restriction du groupe structural à $F$, i.e.
(prop. 15 bis) est localement l'image d'une $\mathcal{C}$-section de $P_1$. Or on
peut supposer pour ceci que $P_1 = X \times G$, donc $P_1/F = X \times (G/F)$, de sorte
qu'une $\mathcal{C}$-section de $P_1/F$ s'identifie à un $\mathcal{C}$-morphisme $f$ de $X$ dans $G/F$.
Comme $G$ est un $\mathcal{C}$-fibré, sur $G/F$, on voit aussitôt que $f$ provient, lo-
calement, d'un $\mathcal{C}$-morphisme de $X$ dans $G$, ce qui prouve notre asser-
tion.

\underline{Remarque} - La condition (2) est toujours vérifiée quand $\mathcal{C}$ est la caté-
gorie des ensembles algébriques sur un corps de définition $k$, mais en
général, dans cette catégorie, il ne sera pas vrai que $G$ fibré par $F$
soit localement trivial, et le corollaire 2 ne sera plus valable.

\newpage

%% ===== Page 46 =====

% typescript page 44 — n° 262

Exemple : $G = \underline{k}^*$, $F$ sous-groupe fini de $G$ non réduit à $e$, $X = G/F$,
$P_1 = X \times G$ d'où $P_1/F = X \times (G/F)$, $s$ la section de $P_1/F$ définie par l'appli-
cation identique de $X$ sur $G/F$ ; cette section provient localement d'une
section $\mathcal{C}$-régulière de $P_1$ si et seulement si $G$ fibré par $F$ est locale-
ment trivial, ce qui n'est pas le cas ici. (On notera que le raisonne-
ment qu'on vient de faire prouve une réciproque du corollaire 2, qui
mérite de figurer dans son énoncé).

Enfin, il faut aussi définir la notion d'extension et restriction
du groupe structural dans des espaces fibrés à fibre \emph{fixée}, $F$ où $G$ opère
fidèlement. Le rédacteur en laisse le soin à son successeur, et ne s'en
permettra pas moins d'utiliser dans la suite cette notion.

12. - \emph{Sections de $\mathcal{C}$-fibrés.}

Pour la notion de $\mathcal{C}$-section, voir définition 3, n° 7. Pour ne pas
s'exposer à d'innombrables confusions, nous n'avons pas cru devoir adop-
ter la suggestion de dire "section" au lieu de " $\mathcal{C}$-section", car
quand il y a dix structures différentes sur un fibré (i.e. dix catégories
$\mathcal{C}$ différentes) il y a ambiguïté décuple ; mais on pourra convenir qu'on
se permet, par abus de langage, de dire simplement "section" quand aucune
confusion ne pourra en résulter. L'ensemble des $\mathcal{C}$-sections d'un $\mathcal{C}$-
fibré (ou simplement d'un $\mathcal{C}$-morphisme $f:E \to X$) est un foncteur contra-
variant par rapport à ce fibré, dont les propriétés d'exactitude (\emph{à gauche})
doivent être explicitées. Les diverses lois de composition sur les fi-
bres passent aux ensembles de sections, et sont respectées par le fonc-
teur précédent. Sorites sur les restrictions de sections.

Soit $P$ un $\mathcal{C}$-fibré principal à droite sur $X$, de groupe $G$, et soit
$F$ un $\mathcal{C}$-espace à groupe d'opérateurs $G$ à gauche, soit enfin $E$ le

\newpage

%% ===== Page 47 =====

% typescript page 45 - n° 262

$\mathcal{C}$-fibré à groupe structural G associé à P et F. Alors les $\mathcal{C}$-sections
de E s'identifient aux $\mathcal{C}$-morphismes f de P dans F, tels que
$f(e.g) = g^{-1}.f(e)$ pour $e \in P, g \in G$. Cette identification est compatible avec
les lois de composition résultant des lois de composition sur F, et est
bien entendu fonctorielle.

13. - \underline{Images réciproques de $\mathcal{C}$-fibrés.}

(Si dans la notion de catégorie locale, on impose la donnée d'une
notion d'image réciproque, il faudra incorporer ce numéro au numéro 9,
où on a dit que les fibrés de type $(\mathcal{C}, \underline{\Phi}, \underline{\Phi}_0)$ forment une catégorie
locale).

Comme au numéro précédent, nous supposons la catégorie $\mathcal{C}$ recollable,
et nous supposons données des catégories $\underline{\Phi}, \underline{\Phi}_0$ donnant lieu à un diagram-
me fonctoriel, et permettant donc de définir la catégorie locale des fi-
brés de type $(\mathcal{C}, \underline{\Phi}, \underline{\Phi}_0)$.

PROPOSITION 15. - Sous les conditions précédentes, soit $E = (E,p,X)$ un
fibré de type $(\mathcal{C}, \underline{\Phi}, \underline{\Phi}_0)$, et soit $f:X' \to X$ un morphisme de $\mathcal{C}$-espaces,
Soit E' la partie de $X' \times E$ formée des $(x',e)$ avec $f(x') = p(e)$. Alors E'
admet une $\mathcal{C}$-structure induite, et si on considère le morphisme $p': E' \to X'$
induit par la projection $X' \times E \to X'$, et sur les fibres ${p'}^{-1}(x')$ les struc-
tures de type $\underline{\Phi}$ déduites par transport de structure des structures sur
les $p^{-1}(f(x'))$, E' devient un fibré de type $(\mathcal{C}, \underline{\Phi}, \underline{\Phi}_0)$.

Utilisant les propositions de localisation 7 et 14, on est ramené
au cas où E est un produit $X \times F$ ($F \in \underline{\Phi}_0$). Alors E' s'identifie en tant
qu'ensemble à $X' \times F$, montrons que la $\mathcal{C}$-structure produit est une $\mathcal{C}$-
structure induite par $X' \times E = X' \times X \times F$ : en effet, l'application d'injection
$(x',e) \to (x',f(x'),e)$ est un $\mathcal{C}$-morphisme, et l'application
$(x',x,e) \to (x',e)$ est un $\mathcal{C}$-morphisme inverse à gauche, d'où notre

\newpage

%% ===== Page 48 =====

% typescript page 46 — n° 262

assertion. Les autres assertions sont alors triviales.

$\xi' = (E',p',X')$ est appelé le \emph{fibré de type $(\mathcal{C}, \underline{\Phi}, \underline{\Phi}_0)$ image
réciproque} de $\xi$, et sera parfois noté $f^{-1}(\xi)$. On aurait pu aussi commen-
cer par définir le \emph{fibré produit} $X' \times E$ sur $X' \times X$, et prendre le \emph{fibré
induit} par ce dernier sur l'image de l'injection $X' \to X' \times X$, de composantes
$1_{X'}$ et $f$. Il convient de prouver les propriétés formelles usuelles de
la notion d'image réciproque (cf. Top. Elém. chap. I), notamment le fait
que l'\emph{image réciproque} $f^{-1}(\xi)$ \emph{est un foncteur en $\xi$}, et la \emph{transitivité
des images réciproques}. Enfin, une \emph{propriété de commutativité de l'opéra-
tion image réciproque, et de l'opération "fibré associé"} étudiées au numé-
ro précédent. La donnée d'un morphisme de fibrés $\xi = (E,p,X) \to \xi' $
$= (E',p',X')$, correspondant à un $\mathcal{C}$-morphisme donné $f : X \to X'$, est
équivalente à la donnée d'un morphisme à \emph{base fixée} de $\xi$ dans $f^{-1}(\xi')$ ;
prenant $E=X$, $p=\text{identité}$, on trouve la condition usuelle de relèvement
d'un $\mathcal{C}$-morphisme $f:X \to X'$ dans la base d'un $\mathcal{C}$-fibré $E'$, à savoir
que le fibré image réciproque admette une $\mathcal{C}$-section.

Soient encore $\xi = (E,p,X)$ un fibré de type $(\mathcal{C}, \underline{\Phi}, \underline{\Phi}_0)$, $f: X' \to X$
un $\mathcal{C}$-morphisme, $\xi' = (E', p', X')$ le fibré image réciproque de $\xi$ par
$f$. Alors toute section ensembliste de $E$ donne par image réciproque une
section ensembliste de $E'$, et on constate aussitôt que \emph{$\mathcal{C}$-sections sont
transformées en $\mathcal{C}$-sections}. (Seule la structure de $\mathcal{C}$-fibrés est inter-
venue). L'application ainsi obtenue est un homomorphisme de foncteurs
par rapport à $\xi$.

14. - \emph{Catégories de modèles}

Soit $\mathcal{M}$ une catégorie, munie d'un foncteur $i_{\mathcal{M}}$ (noté aussi $i$ pour
abréger, s'il ne peut en résulter de confusion) covariant à valeurs dans
la catégorie $\mathcal{T}$ des espaces topologiques. Nous allons montrer comment

\newpage

%% ===== Page 49 =====

% typescript page 47 — n° 262

on peut lui associer une catégorie prétopologique, notée $\hat{\mathcal{M}}$, par une
construction généralisant celle du n° 2. Pour simplifier, nous allons
cependant supposer que les conditions suivantes sont satisfaites (qui
correspondent à CT 1 à CT 4)

\underline{M 1)} \emph{Le foncteur $i_{\mathcal{M}}$ est fidèle.}

Nous dirons alors qu'un $\mathcal{M}$-morphisme $M \to N$ est un \emph{$\mathcal{M}$-isomorphisme},
si la $\mathcal{M}$-structure de $M$ est une structure initiale pour l'application $f$,
i.e. si pour tout $P \in \mathcal{M}$, et toute application $g$ de $P$ dans $M$, $g$ est un
$\mathcal{M}$-morphisme si et seulement si $f g$ l'est. - Appelons \emph{modèles} les élé-
ments de $\mathcal{M}$.

\underline{M 2)} \emph{Tout modèle $M$ admet une base d'ouverts de la forme $f(N)$, où $f$ est un
$\mathcal{M}$-isomorphisme d'un modèle $N$ dans $M$.}

\underline{M 3)} \emph{Un $\mathcal{M}$-isomorphisme, d'un modèle $M$ dans un modèle $N$, dont l'image
est ouverte, est une application ouverte.}

\underline{M 4)} \emph{Soient $f$ une application d'un modèle $M$ dans un autre $N$, $(f_i)$ une
famille de $\mathcal{M}$-isomorphismes de modèles $M_i$ dans $M$. On suppose que les
$f_i(M_i)$ sont des ouverts recouvrant $M$, et que les $f f_i$ sont des $\mathcal{M}$-
morphismes. Alors $f$ est un $\mathcal{M}$-morphisme.}

Une catégorie $\mathcal{M}$, munie d'un foncteur $i=i_{\mathcal{M}}$ à valeurs dans $\mathcal{C}$, et
satisfaisant aux conditions \underline{M 1)} à \underline{M 4)} précédentes est dite une \emph{catégorie
de modèles}. Par exemple, une catégorie prélocale est évidemment une caté-
gorie de modèles. Si $\mathcal{M}$ et $\mathcal{M}'$ sont des catégories de modèles, on dit
que $\mathcal{M}'$ est une \emph{sous-catégorie de modèles} de $\mathcal{M}$ si la classe $\mathcal{M}'$ est con-
tenue dans la classe $\mathcal{M}$, si les $\mathcal{M}'$-morphismes entre deux objets $M, N \in \mathcal{M}'$
sont identiques aux $\mathcal{M}$-morphismes, la loi de composition de $\mathcal{M}'$-morphi-
smes étant celle des $\mathcal{M}$-morphismes (i.e. $\mathcal{M}'$ est une sous-catégorie de
$\mathcal{M}$ ), si le foncteur $i_{\mathcal{M}'}$ est induit par $i_{\mathcal{M}}$, et si enfin la

\newpage

%% ===== Page 50 =====

% typescript page 48 — n° 262

condition suivante est satisfaite :

\underline{Si $M,N$ sont dans $\mathcal{M}'$ et si $f: M \to N$ est un $\mathcal{M}$-isomorphisme de $M$ dans $N$,}
\underline{tel que $f(M)$ soit une partie ouverte, alors $f$ est même un $\mathcal{M}'$-isomorphisme.}

PROPOSITION 16. - Soit $\mathcal{C}$ une catégorie de modèles, et soit $\mathcal{M}$ une classe
contenue dans $\mathcal{C}$. Pour que ce soit une sous-catégorie de modèles, il
faut et il suffit que $\mathcal{M}$ satisfasse la condition suivante :

\underline{M}' 2) \underline{Tout $M \in \mathcal{M}$, admet une base d'ouverts de la forme $f(N)$, où $f$ est}
\indent \indent \underline{un $\mathcal{C}$-isomorphisme d'un $N \in \mathcal{M}$ dans $M$.}

La condition est évidemment nécessaire. Pour montrer qu'elle est suffi-
sante, il suffit manifestement de prouver que, si $f: M \to N$ est un
$\mathcal{M}$-isomorphisme tel que $f(M)$ soit une partie ouverte de $N$, alors $f$
est de même un $\mathcal{C}$-isomorphisme de $M$ dans $N$. (Bien entendu, on met dans
$\mathcal{M}$ la structure de catégorie avec foncteur fondamental induits par $\mathcal{C}$)
Alors les axiomes \underline{M} 1) à \underline{M} 4) pour $\mathcal{M}$ résultent trivialement de ces
mêmes axiomes pour $\mathcal{C}$, donc $\mathcal{M}$ sera une catégorie de modèles, et même
une sous-catégorie de $\mathcal{C}$. Nous devons donc prouver, avec les notations
ci-dessus : si $X \in \mathcal{C}$ et si $g$ est une application de $X$ dans $M$ telle que
$f g$ est un $\mathcal{C}$-morphisme, alors $g$ est un $\mathcal{C}$-morphisme. Or, $f(M)$ étant
ouvert, il résulte de \underline{M}' 2) que $f(M)$ est la réunion d'ouverts de la
forme $f_i(M_i)$, où $M_i \in \mathcal{M}$ et où les $f_i$ sont des $\mathcal{C}$-isomorphismes des $M_i$
dans $N$. Comme $f g$ est continue (étant un $\mathcal{C}$-morphisme), les images réci-
proques $X_i$ des $f_i(M_i)$ par $f g$ sont des parties ouvertes de $X$. D'après
l'axiome \underline{M} 2) sur $\mathcal{C}$, $X_i$ est donc la réunion de parties ouvertes de la
forme $h_j(Y_j)$, où $h_j$ est un $\mathcal{C}$-isomorphisme de $Y_j$ dans $X$. D'après \underline{M} 4)
pour $\mathcal{C}$, il suffit de prouver que les $g h_j$ sont des $\mathcal{C}$-morphismes. Soit

\newpage

%% ===== Page 51 =====

% typescript page 49 — n° 262

$f'_i$ l'application de $M_i$ dans $M$ définie par $ff'_i = f_i$, comme $f$ est un $\mathcal{M}$-isomorphisme (dans), et $ff'_i$ un $\mathcal{M}$-morphisme, $f'_i$ est un $\mathcal{M}$-morphisme, i.e. un $\mathcal{C}$-morphisme. De plus, on a par construction,

\[
\begin{tikzcd}
Y_j \arrow[r, "h_j"] \arrow[d, "g_{ij}"'] & X \arrow[r, "g"] & M \arrow[d, "f"] \\
M_i \arrow[urr, "f'_i"'] \arrow[rr, "f_i"'] & & N
\end{tikzcd}
\qquad
\text{$f f'_i g_{ij} = f g h_j$ pour une application $g_{ij}$ convenable de $Y_j$ dans $M_i$.}
\]

Il suffit donc de prouver que les $g_{ij}$ sont des $\mathcal{C}$-morphismes ou encore, puisque les $f_i$ sont des $\mathcal{C}$-isomorphismes, que $f_i g_{ij}$ est un $\mathcal{C}$-morphisme. Or $f_i g_{ij} = (fg)h_j$, donc est bien un $\mathcal{C}$-morphisme.

D'autres axiomes utiles sur les catégories de modèles sont les suivants, qui généralisent les axiomes de même numéro pour les catégories locales :

M 5) \underline{Pour deux modèles $M,N$, il existe un modèle $P$ tel que $l(P)= M \times N$ et}
\underline{pour lequel la $\mathcal{M}$-structure sur $M \times N$ soit la structure produit des}
\underline{$\mathcal{M}$-structures des facteurs.}
(Bien entendu ce $P$ est alors unique). On dit que $\mathcal{M}$ est une catégorie de modèles multiplicative. Une sous-catégorie de modèles de $\mathcal{M}$ est dite \underline{sous-catégorie multiplicative de modèles} si elle est stable sous l'opération $M \times N$.

M 6) \underline{$\mathcal{M}$ est une catégorie de modèles avec produits, et pour tout modèle}
\underline{$M$, la diagonale de $M \times M$ est fermée.} (Alors $\mathcal{M}$ est dite séparée).

M 7) \underline{Pour tout modèle $M$ et tout ouvert $U \subset M$, $U$ est réunion finie d'ouverts}
\underline{de la forme $f(N)$, où $f$ est un $\mathcal{M}$-isomorphisme de $N$ dans $M$.}

\underline{15. - Catégorie topologique définie par une catégorie de modèles.}

Théorème 1. - Soit $\mathcal{M}$ une catégorie de modèles. Alors on peut trouver une catégorie prétopologique sursaturée $\hat{\mathcal{M}}$, telle que $\mathcal{M}$ en soit une sous-catégorie de modèles, et que tout $X \in \hat{\mathcal{M}}$ admette une base d'ouverts

\newpage

%% ===== Page 52 =====

% typescript page 50 — n° 262

isomorphes à des modèles $M \in \mathcal{M}$. $\hat{\mathcal{M}}$ est déterminé à un isomorphisme
unique près, se réduisant à l'identité sur $\mathcal{M}$. Si $\mathcal{M}$ satisfait à
l'axiome \underline{M 5)}, alors $\hat{\mathcal{M}}$ satisfait à l'axiome \underline{CT 5)} i.e. est une caté-
gorie topologique, et $\mathcal{M}$ en est une sous-catégorie \emph{multiplicative} de
modèles.

Il s'ensuit facilement compte tenu de prop.5 :

Corollaire 1 - Si $\mathcal{M}$ est une catégorie multiplicative de modèles sépa-
rés alors $\mathcal{C}(\mathcal{M}) = \hat{\mathcal{M}}^s$ est une catégorie locale parfaite (dite asso-
ciée à $\mathcal{M}$), dont $\mathcal{M}$ est une sous-catégorie multiplicative de modèles.

De plus, on vérifie aussitôt :

Corollaire 2. - Sous les conditions du corollaire précédent, suppo-
sons la condition \underline{M 7)bis)} vérifiée. Alors les $X \in \mathcal{C}(\mathcal{M})$ qui sont
réunions finies de modèles forment une sous-catégorie multiplicative
semi-parfaite de $\mathcal{C}(\mathcal{M})$ contenant $\mathcal{M}$ (dite \emph{associée} à $\mathcal{M}$)

Remarques.- La prop.6 montre que les axiomes \underline{M 1)} à \underline{M 4)} sont aussi
nécessaires pour la validité du théorème 1. La démonstration du théorè-
me 1, notamment de la première assertion, esquissée ci-dessus, est
assez pénible. On notera qu'elle devient triviale si on sait déjà que
$\mathcal{M}$ est une sous-catégorie de modèles d'une catégorie prétopologique
$\mathcal{C}$ ; ce sera le cas dans toutes les applications raisonnables auxquelles
le rédacteur peut penser en ce moment, mais il semblerait antibourba-
chique de trainer cette hypothèse superflue.

\emph{Démonstration du théorème 1.}

Soit $X$ un espace topologique. On appelle $\mathcal{M}$-\emph{carte}, de $X$ ou simple-
ment carte, un homéomorphisme $f$ de $U$ sur un modèle $M$. (N.B. dans la
suite de ce numéro nous réserverons le nom de "modèle" aux éléments

\newpage

%% ===== Page 53 =====

% typescript page 51 — n° 262

de $\mathcal{M}$). $U$ est dit l'ouvert de la carte, $K$ le modèle de la carte. Deux
cartes $f:U \to K$ et $g:V \to N$ de $X$ seront dites \emph{compatibles} si pour toute
application continue $h$ d'un modèle $P$ dans $U \cap V$, la relation "$fh$ est un
$\mathcal{M}$-morphisme" est équivalente à la relation "$gh$ est un $\mathcal{M}$-morphisme".
Un ensemble $(f_i)$ de cartes $f_i : U_i \to M_i$ de $X$ est dite un \emph{$\mathcal{M}$-atlas},
ou simplement un atlas, si les $U_i$ recouvrent $X$, et si ces cartes sont
deux à deux compatibles. On a alors :

PROPOSITION 17 - Soit $(f_i, U_i, M_i)$ un \emph{$\mathcal{M}$-atlas}. Si deux cartes sont compa-
tibles avec chacune des $f_i$, elles sont compatibles entre elles.

Ceci résultera évidemment du résultat suivant :

Lemme - Considérons une carte $f:U \to M$ compatible avec les cartes $f_i$. Soit
$g:P \to U$ une application continue d'un modèle $P$ dans $U$. Pour tout $i$, soit
$P_i$ la partie ouverte $g^{-1}(U_i)$ de $P$, et soit $(h_{ij})$ une famille de $\mathcal{M}$-isomor-
phismes $h_{ij}$ de modèles $P_{ij}$ dans $P$, tels que les $h_{ij}(P_{ij})$ forment un recou-
vrement ouvert de $P_i$. (Une telle famille existe d'après l'axiome $\underline{\text{M}}$ 2)).
Alors les $f_i g h_{ij} : P_{ij} \to M_i$ sont définis. Pour que $fg$ soit un $\mathcal{M}$-morphis-
me, il faut et il suffit que les $f_i g h_{ij}$ le soient.

En vertu de l'axiome $\underline{\text{M}}$ 4), $fg$ est un $\mathcal{M}$-morphisme si et seulement si
les $f g h_{ij}$ le sont. Or, puisque $\text{Im}(g h_{ij}) \subset U \cap U_i$, et que les cartes $f$ et $f_i$
sont compatibles, $f(g h_{ij})$ est un $\mathcal{M}$-morphisme si et seulement si $f_i(g h_{ij})$
l'est. D'où la conclusion).

Corollaire à la Proposition 17 - Supposons que $\mathcal{M}$ soit un ensemble.
Alors tout atlas d'un espace $X$ est contenu dans un atlas \emph{maximal} unique.

Il suffit de prendre l'ensemble de toutes les cartes de $X$ compatibles
avec toutes les cartes de l'atlas donné.

Corollaire 2 Considérons la relation entre des atlas $A$ et $A'$ de $X$ :

\newpage

%% ===== Page 54 =====

% typescript page 52 — n° 262

"toute carte élément de $A$ est compatible avec toute carte élément de $A'$
C'est là une relation d'équivalence.

En effet, si $(A,A')$ et $(A',A'')$ satisfont la relation, alors en
vertu de proposition 8 appliquée à l'atlas $A'$, une carte de $A$ et une
carte de $A''$, toute carte de $A$ est équivalente à toute carte de $A''$. -

On a envie de définir une $\mathcal{M}$-structure locale sur $X$ comme étant un
atlas maximal, mais cela suppose pratiquement que $\mathcal{M}$ est un ensemble,
et il serait cucu de se borner à ce cas. Aussi choisira-t-on, au moyen
du symbole de Hilbert, un représentant dans chaque "classe d'équivalen-
ce" d'atlas sur $X$, (ce n'est pas beau, mais on a l'habitude) ; c'est
ça une $\mathcal{M}$-structure locale. La classe des espaces $X$ munis d'une
$\mathcal{M}$-structure locale (ou $\mathcal{M}$-espaces locaux) sera notée $\hat{\mathcal{M}}$. On va en
faire une catégorie.

Tout d'abord, soit $X$ un espace muni d'un atlas $(f_i,U_i,M_i)$, consi-
dérons comme dans le lemme une application continue $g$ d'un modèle $P$
dans $X$, et définissons les $P_{ij}$, $h_{ij}$ comme dans le lemme. Je dis que la
condition : "les $f_i g h_{ij}$ sont des $\mathcal{M}$-morphismes" est indépendante du
choix des $h_{ij}$, et ne change pas si on remplace l'atlas donné par un
atlas équivalent. Montrons en effet d'abord qu'elle équivaut à la con-
dition suivante, plus forte en apparence :

(M) pour tout $i$ et tout \emph{$\mathcal{M}$-morphisme} $h:P' \to P$, tel que $gh(P') \subset U_i$,
l'application $f_i gh$ de $P'$ dans $M_i$ est un $\mathcal{M}$-morphisme.

Il suffit de montrer qu'elle l'implique . Or en vertu de $\underline{M}$ 2), $P'$
est réunion d'ensembles ouverts de la forme $h'(P'')$, où $h'$ est un
$\mathcal{M}$-isomorphisme de $P''$ dans $P'$ tel que $hh'(P'')$ soit contenu dans un
des $h_{ij}(P_{ij})$, et en vertu de $\underline{M}$ 4) il suffit de montrer que chaque

\newpage

%% ===== Page 55 =====

% typescript page 53 - n° 262

chaque $(f_i gh)h'$ est un $\mathcal{M}$-morphisme. Or $f_i ghh' = (f_i g)(hh')$, et $hh'=h_{ij}h''$
où $h''$ est une application de $P''$ dans $P_{ij}$. Comme $h_{ij}$ est un $\mathcal{M}$-isomorphis-
me et $h_{ij}h''$ un $\mathcal{M}$-morphisme, $h''$ est un $\mathcal{M}$-morphisme. Il vient alors
$f_i ghh' = (f_i gh_{ij})h''$, qui est un $\mathcal{M}$-morphisme puisque $f_i gh_{ij}$ en est un
par hypothèse, ainsi que $h''$. La condition (M) ci-dessus contient encore
l'atlas donné, montrons qu'elle ne change pas cependant si on remplace
celui-ci par un atlas équivalent $(U'_j, f'_j, M'_j)$. En effet, comme les
$U_i \cap U'_j$ forment un recouvrement ouvert de $X$, il résulte de la condi-
tion \underline{M} 2) que $P$ est réunion d'ouverts de la forme $h(P')$, où $h$ est un
$\mathcal{M}$-isomorphisme de $P'$ dans $P$ tel que $gh(P')$ soit contenu dans un des
$U_i \cap U'_j$. Comme $f_i$ et $f'_j$ sont compatibles, il revient au même de dire que
$f_i (gh)$ soit un $\mathcal{M}$-morphisme, ou que $f'_j(gh)$ en soit un, d'où notre asser-
tion. Nous dirons alors simplement que $g$ est un morphisme, si la condi-
tion (M) énoncée plus haut est vérifiée pour un atlas définissant la
structure de $X$.

Soient maintenant $X$ et $Y$ deux $\mathcal{M}$-espaces locaux, $f$ une application
de $X$ dans $Y$. Nous dirons que $f$ est un \emph{morphisme}, si quel que soit le mor-
phisme $g$ d'un modèle $P$ dans $X$, $fg$ est un morphisme au sens de la défini-
tion qui précède. (Il n'y a pas lieu de comparer ici ces trois notions
de morphismes, définis dans des situations disjointes). Il est trivial
qu'avec cette notion de morphisme, $\hat{\mathcal{M}}$ devient une catégorie, de plus $\hat{\mathcal{M}}$
est muni d'un foncteur évident, noté $i$, à valeurs dans la catégorie $\mathcal{T}$
des espaces topologiques. Les axiomes CT 1) et CT 2) sont trivialement
vérifiés. Vérifions \underline{CT} 3). Soit $X$ un $\mathcal{M}$-espace local, $U$ une partie ouver
te de $X$, montrons qu'on peut trouver un atlas définissant la structure de
$X$ tel que $U$ soit réunion d'ouverts de cartes de cet atlas. Cela signifie

\newpage

%% ===== Page 56 =====

% typescript page 54 — n° 262

qu'on a le

\emph{Lemme 2} - Les ouverts de $X$ qui sont des ouverts de cartes de $X$ (compa-
tibles avec sa structure de $\mathcal{M}$-espace local) forment une base de la
topologie de $X$.

Pour ceci, il suffit de montrer que si $U$ est l'ouvert d'une carte
$f:U \to M$ de $X$ (par quoi on sous-entend toujours : compatible avec la
structure donnée sur $X$), les ouverts de cartes contenus dans $U$ forment
une base de la topologie de $U$. Pour ceci on applique l'axiome $\underline{M}$ 2) à $M$,
et on utilise, pour les $\mathcal{M}$-isomorphismes $g_i:N_i \to M$ obtenus, l'axiome
$\underline{M}$ 3) : les ouverts $f^{-1}(g_i(N_i))$ forment une base de la topologie de $U$,
et sont les ouverts des cartes $g_i^{-1}f$.

Ainsi tout ouvert $U$ de $X$ est réunion des ouverts de cartes de $X$
dedans contenus, on obtient ainsi des atlas de $U$, évidemment tous équi-
valents, d'où une $\mathcal{M}$-structure locale sur $U$, montrons que c'est là
une structure induite : ceci signifie en effet que pour toute applica-
tion continue $g$ d'un modèle $P$ dans $U$, il revient au même de dire qu'elle
est un morphisme de $M$ dans $X$ ou de $M$ dans $U$ ; et cela résulte aussitôt
de la définition. La condition $\underline{CT}$ 4) se vérifie aussi immédiatement.
Ainsi $\hat{\mathcal{M}}$ est une catégorie prélocale, et on a fait ce qu'il fallait
pour qu'elle soye saturée.

Si $M$ est un modèle, il admet un atlas des plus naturels, réduit à
l'application identique. Il définit donc etc.. Le rédacteur, qui en a
marre de faire le travail des autres, refuse de faire trotter l'âne
plus longtemps sur cet air de bombinage. D'ailleurs, le théorème peut
certainement se déduire d'un théorème plus général, où intervient un
foncteur \emph{non fidèle}, et qu'il faudra donner, qui permettra aussi de

\newpage

%% ===== Page 57 =====

% typescript page 55 - n° 262

construire des catégories diverses comme des catégories de variétés
arithmétiques (tels les schémas de Chevalley) des catégories de fibrés
etc.. à partir de petits morceaux. Le théorème 2 du numéro suivant doit
aussi se généraliser dans cette direction.

16 - \emph{Foncteurs topologiques}.

Soient $\mathcal{C}$ et $\mathcal{C}'$ deux catégories prétopologiques. Un foncteur cova-
riant $F$ de $\mathcal{C}$ dans $\mathcal{C}'$ est dit un \emph{foncteur topologique} s'il satisfait
aux deux conditions suivantes :

FT 1) \emph{Si $f: U \rightarrow X$ est un $\mathcal{C}$-isomorphisme de $U$ sur un ouvert de $X$, alors }
$F(f): F(U) \rightarrow F(X)$ \emph{est un $\mathcal{C}'$-isomorphisme de $F(U)$ sur un ouvert de }
$F(X)$.

Si donc $U$ est un ouvert de $X \in \mathcal{C}$, nous pourrons identifier $F(U)$ à un
ouvert de $F(X)$ (avec la structure induite).

FT 2) \emph{Si $X \in \mathcal{C}$ et si $U, V, U_i$ sont des ouverts de $X$, on a $F(U \cap V) =$}
$= F(U) \cap F(V)$, $F(\bigcup U_i) = \bigcup F(U_i)$.

On dit que $F$ est \emph{multiplicatif}, si $\mathcal{C}$ et $\mathcal{C}'$ sont des catégories topologi-
ques, et si $F$ satisfait à l'axiome supplémentaire:

FT 3) \emph{Si $X,Y \in \mathcal{C}$, l'application naturelle $F(X \times Y) \rightarrow F(X) \times F(Y)$ est un $\mathcal{C}'$-}
\emph{isomorphisme de $F(X \times Y)$ sur $F(X) \times F(Y)$.}

On dit que $F$ est un \emph{foncteur local} s'il satisfait à \emph{FT 1)} et à

FL ) \emph{On a $i_{\mathcal{C}'} F = i_{\mathcal{C}}$ (identité fonctorielle)}

Alors la condition \emph{FT 2)} est trivialement satisfaite, donc un foncteur
local est un foncteur topologique. Exemple de foncteur local : le fonc-
teur d'injection d'une sous-catégorie prélocale de $\mathcal{C}'$ dans $\mathcal{C}'$.

Soit $\mathcal{M}$ une classe contenue dans $\mathcal{C}$ et telle que tout $X \in \mathcal{C}$ admette
une base d'ouverts isomorphes à des éléments de $\mathcal{M}$ (appelés \emph{modèles}) :

\newpage

%% ===== Page 58 =====

% typescript page 56 — n° 262

$\mathcal{M}$ est donc une catégorie de modèles (prop.16) et $\hat{\mathcal{C}}$ s'identifie à
$\hat{\mathcal{M}}$ (résulte de l'assertion d'unicité dans le théorème 1, qui est d'ail-
leurs à peu près triviale). Nous dirons aussi que $\mathcal{M}$ est une catégorie
de modèles engendrant $\mathcal{C}$. Si $F$ est un foncteur topologique de $\mathcal{C}$ dans
$\mathcal{C}'$, sa restriction à $\mathcal{M}$ satisfait manifestement aux conditions sui-
vantes :

\underline{FTM 1)} \emph{Si} $f: M \to N$ \emph{est un} $\mathcal{M}$\emph{-isomorphisme d'un modèle sur une partie
ouverte d'un autre, alors} $F(f)$ \emph{est un isomorphisme de} $F(M)$ \emph{sur
une partie ouverte de} $F(N)$.

\underline{FTM 2)} \emph{Soient} $f': M' \to M$, $f'': M'' \to M$, $f_i: M_i \to M$ \emph{des} $\mathcal{M}$\emph{-isomorphismes
de modèles sur des parties ouvertes de} $M$\emph{, tels que
$\bigcup_i f_i(M_i) \supset f'(M') \cap f''(M'')$. Alors $\bigcup_i F(f_i)(F(M_i)) \supset F(f')(F(M')) \cap$
$\cap F(f'')(F(M''))$.}

De plus, si $F$ est multiplicatif, on a :

\underline{FTM 3)} \emph{Si} $M$ \emph{et} $N$ \emph{sont deux modèles, l'application naturelle} $F(M \times N) \to$
$F(M) \times F(N)$ \emph{est un} $\mathcal{C}'$\emph{-isomorphisme.}

Si $F$ est un foncteur local, on a :

\underline{FLL} \quad $i_{\mathcal{C}'} F = F i_{\mathcal{M}}$ , (\emph{identité fonctorielle sur} $\mathcal{M}$).

Théorème 2 - Soit $\mathcal{C}$ une catégorie prétopologique, $\mathcal{M}$ une catégorie de
modèles engendrant $\mathcal{C}$, $F$ un foncteur de $\mathcal{M}$ dans une catégorie prétopo-
logique sursaturée $\mathcal{C}'$. Si $F$ satisfait aux conditions \underline{FTM 1)} et \underline{FTM 2)},
il existe un foncteur prétopologique $\hat{F}$ de $\mathcal{C}$ dans $\mathcal{C}'$ dont la restric-
tion à $\mathcal{M}$ soit $F$. Si $F$ est multiplicatif (i.e. satisfait à \underline{FTM 3)}) il
en est de même de $\hat{F}$. Si $\hat{F}'$ satisfait à la même condition que $\hat{F}$, il
existe un isomorphisme unique de $\hat{F}$ sur $\hat{F}'$ se réduisant à l'identité
sur $F$.

\newpage

%% ===== Page 59 =====

% typescript page 57 — n° 262

La démonstration ressemble comme un frère, en moins long, à celle
du théorème 1, aussi ne la donnerons nous pas. Il faut encore regarder
à quelle condition $\hat{F}$ transforme des $\mathcal{C}$-espaces séparés en $\mathcal{C}'$-espaces
séparés. On montre facilement qu'il suffit de le vérifier pour les $\mathcal{C}$-
espaces séparés qui sont réunions de \emph{deux} ouverts isomorphes à des modè-
les. Voici un cas important où on peut se dispenser de faire la vérifi-
cation.

Corollaire 1 - Si F satisfait à \underline{FLM}, alors on peut trouver un $\hat{F}$ qui soit
un foncteur local. Plus généralement, supposons donné un homomorphisme
fonctoriel $i_{\mathcal{M}} \to i_{\mathcal{C}'} F$ tel que pour tout modèle M, l'application conti-
nue $M \to F(M)$ soit bijective, Si F est multiplicatif, alors $\hat{F}$ transforme
$\mathcal{C}$-espaces séparés en $\mathcal{C}'$-espaces séparés. Dans ce cas, la condition
du th. 2 reste valable si on suppose $\mathcal{C}$ séparée et $\mathcal{C}'$ parfaite.

En effet, l'homomorphisme fonctoriel donné définit un homomorphisme
fonctoriel analogue relatif à $\hat{F}$, permettant d'identifier les ensembles
X et $\hat{F}(X)$, la topologie de X étant moins fine. Par suite, si la diagonale
de $X \times X$ est fermée pour la topologie de $X \times X$, elle l'est à fortiori pour
celle de $\hat{F}(X) \times \hat{F}(X) = \hat{F}(X \times X)$.

Corollaire 2 - Pour un foncteur variable F de $\mathcal{M}$ dans $\mathcal{C}'$ satisfaisant aux
conditions \underline{FTM 1}) et \underline{FTM 2}), $\hat{F}$ est un "foncteur" en F, et de façon préci-
se : tout homomorphisme du foncteur F dans un foncteur analogue G se
prolonge de façon unique en un homomorphisme de $\hat{F}$ en $\hat{G}$.

La démonstration directe à partir du théorème 2 ne semble pas bien
longue. Mais on aimerait bien avoir une fois pour toutes des énoncés qui
permettent de voir que certains termes sont des "foncteurs" par rapport
à certaines lettres, et qui se substitueraient avantageusement aux énoncés
étriqués de \underline{Ens}, 1 éd., chap. 4, Appendice.

\newpage

%% ===== Page 60 =====

% typescript page 58 — n° 262

Corollaire 3 -(transitivité) : $\widehat{GF} = \hat{G} \hat{F}$, si $F$ prend ses valeurs dans
une sous-catégorie de modèles $\mathcal{M}'$ de $\mathcal{C}'$ engendrant $\mathcal{C}'$, et si $G$ est un
foncteur, analogue à $F$, défini sur $\mathcal{M}'$.

Corollaire 4 - Si $\mathcal{M}$ est une sous-catégorie de modèles de $\mathcal{C}'$ (cf. prop.
16), et si $F$ est le foncteur d'injection de $\mathcal{M}$ dans $\mathcal{C}'$, alors $F$ iden-
tifie $\mathcal{C}$ à une sous-catégorie prétopologique de $\mathcal{C}'$, et même à une sous-
catégorie topologique si $\mathcal{M}$ est une sous-catégorie \emph{multiplicative} de
modèles dans $\mathcal{C}'$.

La démonstration directe de ce corollaire est d'ailleurs à peu
près triviale (voir remarque après le th. 1)

\newpage

%% ===== Page 61 =====

% typescript page 59 — n° 262

\section*{§ 2 - \emph{Catégories de variétés.}}

1. - \emph{Espaces $K$-annelés.}

\emph{Définition 1} - Soit $K$ un anneau avec unité. On appelle \emph{espace $K$-annelé}
un espace topologique $X$ muni d'un sous-faisceau d'anneaux $\underline{O}_X$ du faisceau
des germes d'applications de $X$ dans $K$, contenant les germes unité.

$\underline{O}_X$ est appelé le \emph{faisceau fondamental} de $X$. En vertu des définitions
relatives aux faisceaux, la donnée d'une structure $K$-annelée sur $X$ équi-
vaut à la donnée, pour tout ouvert $U \subset X$, d'un sous-anneau $\underline{O}_X(U)$ de
l'anneau des applications de $U$ dans $K$, contenant l'unité, de telle façon
que les conditions suivantes soient satisfaites : (i) si $V$ est un ouvert
contenu dans $U$, la restriction à $V$ de toute $f \in \underline{O}_X(U)$ est dans $\underline{O}_X(V)$ ;
(ii) si inversement $f$ est une application de $U$ dans $K$, et si tout point
de $U$ admet un voisinage ouvert $V$ tel que la restriction de $f$ à $V$ soit
dans $\underline{O}_X(V)$, alors $f$ est dans $\underline{O}_X(U)$. - Les fonctions dans $\underline{O}(U)$ sont appe-
lées \emph{fonctions régulières} sur l'ouvert $U$ de l'espace $K$-annelé $X$.

\emph{Définition 2} - Soient $X, Y$ deux espaces $K$-annelés. On appelle \emph{morphisme}
d'espaces $K$-annelés de $X$ dans $Y$, une application continue $f$ de $X$ dans $Y$
telle que pour toute fonction régulière $g$ sur une partie ouverte $V$ de $Y$,
la fonction $gf$ sur $f^{-1}(V)$ soit régulière.

Si $X, Y, Z$ sont trois espaces $K$-annelés, $f: X \to Y$ et $g: Y \to Z$ des mor-
phismes d'espaces $K$-annelés, alors $gf$ est un morphisme de $X$ dans $Z$. Par
suite, les espaces $K$-annelés forment une \emph{catégorie}, et le foncteur $i$ sur
cette dernière, qui attache à tout espace $K$-annelé l'espace topologique
sous-jacent, et à un morphisme $X \to Y$ d'espaces $K$-annelés ce même morphis-
me considéré comme application de l'espace topologique sous-jacent à $X$
dans celui sous-jacent à $Y$, est un \emph{foncteur fidèle}.

\newpage

%% ===== Page 62 =====

% typescript page 60 - n° 262

PROPOSITION 1 - Soit $(f_i)$ une famille d'applications continues d'un
ensemble $X$ dans des espaces $K$-annelés $X_i$. Alors il existe sur $X$ une
structure d'espace $K$-annelé \emph{initiale} pour la famille $(f_i)$. La topolo-
gie sous-jacente de $X$ est la moins fine qui rende continue les applica-
tions $f_i$, et le faisceau fondamental de $X$ est le plus petit sous fais-
ceau d'anneaux du faisceau des germes d'applications de $X$ dans $K$ pour
lequel les fonctions de la forme $g_i f_i$ sur $f_i^{-1}(U_i)$, (où $g_i$ est une fonc-
tion régulière sur l'ouvert $U_i$ de $X_i$), sont des sections.

En effet, on voit aussitôt en vertu des définitions que muni de la
topologie et du faisceau ainsi explicités, $X$ est un espace $K$-annelé, et
que pour toute application $f$ d'un espace $K$-annelé $Y$ dans $X$, $f$ est un
morphisme si et seulement si les $f_i f$ le sont.

Corollaire 1 - Soit $X$ un espace $K$-annelé, $Y$ une partie de $X$. Alors $Y$
admet une structure $K$-annelée \emph{induite}. La topologie sous-jacente est
celle induite par $X$, et une fonction sur un ouvert $U$ de $Y$ est régulière
si et seulement si tout $y \in U$ admet un voisinage ouvert $V$ dans $X$ tel que
$V \cap Y \subset U$ et que $f|V \cap Y$ induite par une fonction régulière sur $V$. En
particulier, si $Y$ est ouvert, une fonction $f$ sur $U$ est régulière au
sens de $Y$ si et seulement si elle l'est au sens de $X$, i.e. le faisceau
fondamental de $Y$ est induit par le faisceau fondamental de $X$. .

Corollaire 2 - Les espaces $K$-annelés forment une catégorie topologique
sursaturée.

Nous avons déjà signalé que les axiomes CT 1) et CT 2) sont véri-
fiés, il en est de même de CT 3) et CT 5) en vertu de proposition 1
(existence de structures induites sur les ouverts, et de structures
produits), et aussi CT 4) affirmant que la notion de morphisme est de

\newpage

%% ===== Page 63 =====

% typescript page 61 — n° 262

nature locale. Enfin, si $X$ est un espace topologique muni d'un recouvre-
ment ouvert $(U_i)$, et si sur tout $U_i$ on se donne une structure $K$-annelée,
les structures induites par $U_i$ et $U_j$ dans $U_i \cap U_j$ étant identiques, alors
il résulte de la fin du corollaire 1 que les faisceaux fondamentaux de
$U_i$ et $U_j$ ont même restriction à $U_i \cap U_j$, donc qu'il existe un faisceau
unique sur $X$ induisant sur chaque $U_i$ le faisceau fondamental de $U_i$. Muni
de ce faisceau, $X$ est évidemment un espace $K$-annelé, et d'après la fin
du corollaire 1, il induit sur chaque $U_i$ la structure $K$-annelée donnée
sur $U_i$. Donc la catégorie des espace $K$-annelés est bien sursaturée.

PROPOSITION 2 - Soit $(f_i)$ une famille d'applications d'espaces $K$-annelés
$X_i$ dans un ensemble $X$. Alors il existe sur $X$ une structure $K$-annelée fi-
nale. La topologie sous-jacente est la plus fine des topologies rendant
continues les $f_i$. Une fonction $g$ sur un ouvert $U \subset X$ est régulière si et
seulement si les fonctions $gf_i$ sur les ouverts $f_i^{-1}(U) \subset X_i$ sont régu-
lières.

En effet, si $X$ est muni de la topologie la plus fine rendant les $f_i$
continues, il est immédiat que les fonctions sur un ouvert $U \subset X$ qui sa-
tisfont à la condition énoncée sont exactement les sections d'un faisceau
sur $X$ définissant une structure d'espace annelé. Il résulte alors des
définitions que la structure $K$-annelée ainsi définie sur $X$ est bien struc-
ture finale pour la famille des $f_i$.

Corollaire - Soit $X$ un espace $K$-annelé, $R$ une relation d'équivalence
dans $X$, $g$ l'application canonique de $X$ sur $X/R$. Alors $X/R$ admet une
structure $K$-annelée quotient. La topologie sous-jacente est la topologie
quotient, et une fonction sur un ouvert $U \subset X/R$ est régulière si et seu-
lement si la fonction $fg$ est régulière sur l'ouvert $g^{-1}(U)$ de $X$.

\newpage

%% ===== Page 64 =====

% typescript page 62 — n° 262

2. - \emph{Catégories de variétés sur K.}

\emph{Définition 3} - Soit $K$ un anneau commutatif avec unité. \emph{Une catégorie de
variétés} sur $K$ est une catégorie locale recollable $\mathcal{V}$ (par.1, n$^\circ$4),
munie d'un objet $\underline{K} \in \mathcal{V}$ dont l'ensemble sous-jacent est $K$, et satisfai-
sant aux axiomes CV 1), CV 2) et CV 3) énoncés plus bas.

Les objets de $\mathcal{V}$ seront appelés \emph{$\mathcal{V}$-variétés} ou \emph{$\mathcal{V}$-espaces} (ce
dernier terme semblant préférable, ne créant pas de risques de confu-
sions), et on employera librement la terminologie introduite pour les
catégories topologiques, parlant de $\mathcal{V}$-structures sur un ensemble donné
etc... Nous noterons, suivant l'abus de notation courant, $\underline{K}$ par la
lettre $K$ ; sauf avis expresse du contraire, $K$ désigne donc l'anneau $K$,
muni en plus de la $\mathcal{V}$-structure $\underline{K}$ . Nous dirons qu'une application $f$
d'un $\mathcal{V}$-espace $X$ dans un $Y$, est \emph{$\mathcal{V}$-régulière}, ou simplement \emph{régu-
lière}, si c'est un $\mathcal{V}$-morphisme.

CV 1) \emph{Il existe un $\mathcal{V}$-espace réduit à un point. Si $e$ est un tel $\mathcal{V}$-espace
toute application $f$ d'un $\mathcal{V}$-espace $X$ dans $e$ est un morphisme, et
une application $g$ de $e$ dans $X$ est un morphisme si et seulement si
$f(e)$ est fermé dans $X$.}

Il résulte immédiatement de cet axiome que tout ensemble $e'$ réduit à un
point admet une structure de $\mathcal{V}$-espace et une seule.

CV 2) \emph{Pour tout $\mathcal{V}$-espace $X$, les morphismes de $X$ dans $K$ forment un sou
anneau, contenant l'unité, de l'anneau de toutes les applications
de $X$ dans $K$.}

Montrons que cet axiome équivaut au suivant :

CV 2') \emph{La partie de $K$ réduite à l'élément unité $1$ est fermée. Les appli-
cations $(x,y) \to x-y$ et $(x,y) \to xy$ de $K \times K$ dans $K$ sont des mor-
phismes.}

\newpage

%% ===== Page 65 =====

% typescript page 63 - n° 262

\underline{CV 2)} implique \underline{CV 2')} : $\{1\}$ est l'image d'un $\mathcal{V}$-espace $e$ réduit à
un point par un morphisme $i$ de $e$ dans $K$ (en vertu de \underline{CV 2}) donc fermé en
vertu de \underline{CV 1}) ; les fonctions $(x,y) \to x$ et $(x,y) \to y$ sur $K \times K$ sont des
morphismes, donc en vertu de \underline{CV 2}) il en est de même des fonctions $x-y$ et
$xy$. Réciproquement, supposons \underline{CV 2'}) satisfaite, il en résulte tout
d'abord que pour tout $\mathcal{V}$-espace $X$, l'application $x \to 1$ de $X$ dans $K$ est
un morphisme, puisque composé de $X \to e$ et $i : e \to K$ qui sont des morphis-
mes en vertu de \underline{CV 1}). Soient maintenant $f,g$ deux morphismes de $X$ dans $K$
montrons que $f-g$ et $f g$ sont encore des morphismes : or ce sont respecti-
vement les applications composées $S \circ (f \times g) \circ \Delta$ et $P \circ (f \times g) \circ \Delta$, où $\Delta$ est
l'application diagonale $X \to X \times X$, et $S$ et $P$ sont les applications de $K \times K$
dans $K$ envisagées dans l'énoncé de \underline{CV 2'}) ; comme $S, P, \Delta$ sont des mor-
phismes, la conclusion voulue s'ensuit.

PROPOSITION 3 - Soit $X$ un $\mathcal{V}$-espace. A tout ouvert $U \subset X$, attachons l'en-
semble des morphismes de $U$ dans $K$. Alors $X$ devient un espace $K$-annelé
(définition 1). Si $X,Y$ sont deux $\mathcal{V}$-espaces, tout morphisme de $X$ dans
$Y$ est un morphisme pour les structures $K$-annelées sous-jacentes (défi-
nition 2).

Cela résulte aussitôt des définitions.

\underline{CV 3)} \underline{Soient $X,Y$ deux $\mathcal{V}$-espaces, $f$ une application de $X$ dans $Y$ qui soit}
\underline{un morphisme pour les structures $K$-annelées sous-jacentes, alors}
\underline{$f$ est un $\mathcal{V}$-morphisme.}

Cet axiome signifie aussi que le foncteur naturel de $\mathcal{V}$ dans la catégo-
rie des espaces $K$-annelés, défini par la proposition 2, identifie à $\mathcal{V}$
une sous-catégorie prétopologique (par.1, n$^{\circ}$1) de cette dernière, en
effet, il suffit de vérifier que si $U$ est un ouvert d'un $\mathcal{V}$-espace $X$,

\newpage

%% ===== Page 66 =====

% typescript page 64 - n° 262

alors la structure $K$-annelée associée à la $\mathcal{V}$-structure induite est
induite par celle de $X$, ce qui résulte aussitôt des définitions et
prop. 1, cor. 1. On fera attention cependant que dans tous les cas que
nous aurons à envisager, on n'a pas là une sous-catégorie topologique,
le foncteur d'inclusion n'étant pas multiplicatif.

Soit $e$ un $\mathcal{V}$-espace réduit à un point. Identifiant les applica-
tions de $e$ dans $K$ aux éléments de $K$, l'ensemble des morphismes de $e$
dans $K$ est une partie de $K$, qui d'après \underline{CV} 2) est un sous-anneau avec
unité, noté $k$. Ce dernier ne dépend évidemment pas du choix de $e$, on
l'appelle \emph{l'anneau de définition} $k$ de la catégorie $\mathcal{V}$, $K$ s'appelant le
\emph{domaine des valeurs}. D'après \underline{CV} 1), $k$ est aussi l'ensemble des éléments
$x$ de $K$ tels que $\{x\}$ soit fermé.

PROPOSITION 4 - Soit $X$ un $\mathcal{V}$-espace. L'anneau des fonctions régulières
sur $X$ contient l'anneau des applications constantes de $X$ dans $k$, et est
donc une algèbre sur le corps de définition $k$ de $\mathcal{V}$. Donc le faisceau
fondamental $\underline{O}(X)$ de $X$ est un faisceau de $k$-algèbres.

En effet, une application constante de $X$ dans $k$ se factorise en
une application de $X$ dans $e$, et une application de $e$ dans $k$. Par défi-
nition de $k$, cette dernière est un morphisme de $e$ dans $K$, il en est de
même de $X \to e$ en vertu de \underline{CV} 1), donc aussi du composé.

Corollaire 1.- Soient $m,n$ deux entiers $> 0$. Toute application de $K^m$
dans $K^n$ qui est polynomiale à coefficients dans $k$, est un morphisme, en
particulier les applications $x \to \lambda x$ ($\lambda \in k$) et $x \to x+a$ ($a \in K^n$) de $K^n$
dans lui-même sont des morphismes. Inversement, supposons $K$ et $k$ des
corps, $K$ infini, pour qu'une application polynomiale $f: K^m \to K^n$ soit un
morphisme, il faut que ses coefficients soient dans $k$.

\newpage

%% ===== Page 67 =====

% typescript page 65 — n° 262

Pour la première assertion, il suffit de la vérifier quand $n=1$
auquel cas c'est un cas particulier de prop. 4. De même pour la récipro-
que, on peut supposer $n=1$, comme chaque élément de $k$ est une partie
fermée de $K$, chaque élément de $k^n$ est une partie fermée de $K^n$, donc une
fonction régulière $f$ sur $K^n$ prend sur $k^n$ des valeurs qui sont dans $k$
(en vertu de \underline{CV} 1). Si donc $f$ est polynômiale, ses coefficients sont
nécessairement dans $k$.

\underline{Corollaire 2} - Les conditions suivantes sur $\mathcal{V}$ sont équivalentes :
a) $k=K$ b) Toute application constante d'un $\mathcal{V}$-espace dans un autre est
un morphisme c) pour tout $\mathcal{V}$-espace $X$, les points de $X$ sont des par-
ties fermées.

a) implique que toute application constante d'un $\mathcal{V}$-espace $X$ dans
$K$ est un morphisme, donc aussi toute application constante de $X$ dans
un $\mathcal{V}$-espace $Y$, en vertu de \underline{CV} 3). b) implique que toute application
du $\mathcal{V}$-espace $e$ réduit à un point, dans un $\mathcal{V}$-espace $X$, est un morphis-
me, donc a une image fermée en vertu de \underline{CV} 1), d'où c). Enfin c) impli-
que a), puisque $k$ est l'ensemble des points de $K$ qui sont fermés.

Nous allons donner une méthode universelle pour construire des ca-
tégories de variétés :

PROPOSITION 5 - Soit $\mathcal{M}$ une sous-catégorie de modèles (cf. par.1, prop.
16) de la catégorie des espaces $K$-annelés, donc définie par une sous-
classe $\mathcal{M}$ de cette dernière, satisfaisant à la condition :

(i) tout modèle admet une base d'ouverts isomorphes à des modèles.

On suppose de plus les conditions suivantes satisfaites :

(ii) Il existe des modèles réduits à un point, et si $e$ est un tel modèle
toute application d'un modèle $M$ dans $e$ est un morphisme, et une

\newpage

%% ===== Page 68 =====

% typescript page 66 - n° 262

application de $e$ dans $M$ est un morphisme si et seulement si son image
est fermée.

(iii) On se donne un modèle $\underline{K}$, dont l'ensemble sous-jacent est $K$, tel
que pour tout modèle $M$, les applications régulières de $M$ dans $K$ soient
précisément les morphismes de $M$ dans $\underline{K}$.

Alors la catégorie topologique parfaite $\mathcal{V}$ associée à $\mathcal{M}$ (par.1,
n$^\circ$ 15, th.1 cor.1) est une catégorie parfaite de variétés sur $K$, et
s'identifie à la catégorie topologique séparée, associée à la sous-
catégorie prétopologique de la catégorie topologique des espaces K-anne-
lés formée des espaces K-annelés admettant un système fondamental d'ou-
verts isomorphes à des modèles. Pour tout $X \in \mathcal{V}$, les morphismes de $X$
dans $K$ sont les fonctions régulières sur l'espace K-annelé $X$. Si $\mathcal{M}$
satisfait l'axiome de finitude $\underline{\text{M}}\ 7')$ (par.1, n$^\circ$ 14) alors les objets
de $\mathcal{V}$ qui sont réunions \emph{finies} d'ouverts isomorphes à des modèles for-
ment une catégorie semi-parfaite de variétés sur $K$.

D'après le cor. 4 au th.2, la catégorie topologique sursaturée
associée à $\mathcal{M}$ s'identifie à la sous-catégorie prétopologique de la caté-
gorie topologique des espaces K-annelés, formée des espaces K-annelés
admettant un système fondamental d'ouverts isomorphes à des modèles.
Donc la catégorie parfaite $\hat{\mathcal{M}}$ associée à $\mathcal{M}$, qui est une sous-caté-
gorie topologique de $\hat{\mathcal{M}}$ (par.1, prop.4) donc une sous-catégorie préto-
pologique de la catégorie topologique des espaces K-annelés, est bien
celle indiquée dans l' énoncé de prop.5. Reste à, prouver que cette
catégorie $\mathcal{V}$ satisfait aux axiomes $\underline{\text{CV}}\ 1)$, $\underline{\text{CV}}\ 2)$ et $\underline{\text{CV}}\ 3)$. Tout d'abord
ce qui précède implique que les morphismes d'un $X \in \hat{\mathcal{M}}$ dans $\underline{K}$ sont
précisément les fonctions régulières sur l'espace K-annelé $X$ (car il

\newpage

%% ===== Page 69 =====

% typescript page 67 - n° 262

en est ainsi si $X$ est un modèle en vertu de la condition (iii), et les
ouverts de $X$ isomorphes à des modèles recouvrent $X$). Cela implique les
\emph{axiomes} \emph{CV} 2) et \emph{CV} 3). Quant à \emph{CV} 1), c'est une conséquence immédiate
de (ii). Les mêmes raisonnements s'appliquent à $\mathcal{V}_f$, si $\mathcal{M}$ satisfait
à l'axiome de finitude M 7 bis) (par. 1, n$^\circ$ 14) : alors $\mathcal{V}_f$ est une sous-
catégorie topologique semi-parfaite de $\mathcal{M}$ contenant $\mathcal{M}$ (par. 1 th. 1,
cor. 2), donc une sous-catégorie prétopologique de la catégorie des espa-
ces $K$-annelés, d'où aussitôt les axiomes \emph{CV} 1), \emph{CV} 2) et \emph{CV} 3).

3. - \emph{Structures induites.}

Nous supposons donnée dans ce numéro et le suivant une catégorie $\mathcal{V}$
de variétés sur un anneau $K$. Soit $X$ un $\mathcal{V}$-espace, $Y$ une partie de $X$.
Nous savons (Ensembles, Chap. IV, 2ème éd.) ce que signifie l'assertion
que $Y$ admet une structure de $\mathcal{V}$-espace induite, i.e. une structure ini-
tiale pour l'application identique de l'ensemble $Y$ dans le $\mathcal{V}$-espace.
Cependant, dans certains cas importants, il convient de restreindre cette
notion (pour éviter des structures induites de nature pathologique, com-
me il s'en présente par exemple pour une partie "très irrégulière" d'une
variété différentiable ou analytique - ayant alors une structure induite
de variété discrète - ou une courbe avec point de rebroussement $x=t^2$,
$y=t^3$, dans une variété analytique complexe, (cf exercices ultérieurs).

\emph{Définition} 4 - Soit $X$ un $\mathcal{V}$-espace, $Y$ une partie de $X$, $i$ l'injection
canonique de $Y$ dans $X$. $Y$ est dite un \emph{$\mathcal{V}$-espace plongé}, s'il admet une
structure $S$ de $\mathcal{V}$-espace induite satisfaisant à la condition suivante :
$Y$ étant muni de la topologie $T$ sous-jacente à $S$, pour qu'un germe
d'application de $Y$ dans $K$ en un point $y \in Y$ soit un germe de fonction
régulière, il faut et il suffit qu'il soit l'image réciproque par $i$ d'un

\newpage

%% ===== Page 70 =====

% typescript page 68 — n° 262

germe de fonction régulière sur $X$ en $y$. On dit que $Y$ est un sous-$\mathcal{V}$-
espace de $X$ si c'est un sous-$\mathcal{V}$-espace plongé, et si de plus la topolo-
gie $T$ sous-jacente à la structure de $\mathcal{V}$-espace induite est aussi la
topologie induite par $X$.

On peut aussi dire que $Y$ est un sous-$\mathcal{V}$-espace plongé s'il existe
une structure $S$ de $\mathcal{V}$-espace induite, et si la structure $K$-annelée
sous-jacente est identique à la structure $K$-annelée initiale pour
\emph{l'application continue $i$ de l'espace topologique} $Y$ (muni de $T$) dans le
$\mathcal{V}$-espace $X$ ; et que $Y$ est un sous-$\mathcal{V}$-espace, si la structure $K$-anne-
lée sous-jacente à $S$ est identique à la structure $K$-annelée induite par
$X$ (prop.1) \emph{i.e.} à la structure $K$-annelée initiale pour \emph{l'application} $i$
de \emph{l'ensemble} $Y$ dans le $\mathcal{V}$-espace $X$. [En exercice, on montrera qu'il
peut exister des variétés plongées qui ne sont pas des sous-variétés
(géodésiques du tore). On soumet aussi à Bourbaki la question s'il ne
faut pas imposer à un sous-$\mathcal{V}$-espace plongé d'être connexe pour la
topologie $T$ sous-jacente à la structure de $\mathcal{V}$-espace induite (sinon le
cas pathologique de parties totalement discontinues et analogues d'une
bonne variété ne serait pas exorcisée) ; mais bien entendu, il serait
absurde d'exiger la même chose pour un sous-$\mathcal{V}$-espace.] Une partie
ouverte de $X$ est toujours un sous-$\mathcal{V}$-espace, comme il résulte de
l'axiome \underline{CT} 3)des catégories topologiques (par.1, n° 1).

PROPOSITION 5 - Supposons que $\mathcal{V}$ satisfasse à l'axiome suivant :

\underline{CV} 4) \emph{Soit $f$ un morphisme d'un $\mathcal{V}$-espace $Y$ dans un autre $X$, $y$ un point
de $Y$ tel que tout germe de fonction régulière $g$ en $y$ provienne
d'un germe de fonction régulière $h$ en $x=f(y)$ par $g=h \circ f$. Alors il
existe un voisinage ouvert $V$ de $y$, et un voisinage ouvert $U$ de $x$,}

\newpage

%% ===== Page 71 =====

% - 69 - n° 262

\emph{tel que la restriction de $f$ à $U$ soit un isomorphisme de $U$ sur un sous-$\mathcal{V}$-espace fermé de $V$.}

Alors $i$ est l'application d'injection d'un $\mathcal{V}$-espace plongé $Y$
dans un $\mathcal{V}$-espace $X$, $i$ est localement un homéomorphisme dans, et si $Y$
est un sous-$\mathcal{V}$-espace de $X$, $Y$ est une partie localement fermée de $X$,
donc un sous-$\mathcal{V}$-espace fermé d'une partie ouverte $U$ de $X$.

La démonstration de cette proposition est triviale. (Bien entendu,
l'axiome \underline{CV} 4) est bien plus fort que la conclusion de cette proposi-
tion, puisque d'un hypothèse de nature "ponctuelle" il permet de tirer
une conclusion valable dans tout un voisinage du point envisagé ; c'est
là une forme affaiblie du théorème des fonctions implicites, qui est
valable aussi pour les variétés algébriques.

PROPOSITION 6 - Soit $X$ un $\mathcal{V}$-espace, $Y$ un $\mathcal{V}$-espace plongé(resp. un
sous-$\mathcal{V}$-espace) dans $X$. munissons $Y$ de la structure de $\mathcal{V}$-espace
induite et soit $Z$ un $\mathcal{V}$-espace plongé(resp. un sous-$\mathcal{V}$-espace) dans
$Y$. Alors $Z$ est un sous-$\mathcal{V}$-espace plongé (resp. un sous-$\mathcal{V}$-espace) de
$X$ et sur $Z$ les structures de $\mathcal{V}$-espaces induites par $X$ et $Y$ sont les
mêmes.

En effet, d'après la transitiviré des structures induites (Ens.
chap. IV, 2ème éd.) $Z$ admet dans $X$ une structure de $\mathcal{V}$-espace induite,
identique à celle induite par $Y$. Par hypothèse, la structure $K$-annelée
sous-jacente à cette dernière est l'image réciproque, par morphisme
d'injection $i:Z \to Y$ de l'espace (resp. l'ensemble) $Z$ dans $Y$, de la struc
ture $K$-annelée de $Y$, et cette dernière est de même l'image réciproque
par le morphisme d'injection $j:Y \to X$, de la structure $K$-annelée de $X$.
D'après loc. cité, la structure $K$-annelée de $Z$ est donc l'image

\newpage

%% ===== Page 72 =====

% typescript page 70 - n° 262

réciproque, par le morphisme $j i : Z \to X$ de l'espace (resp. l'ensemble)
$Z$ dans $X$, de la structure K-annelée de $X$, i.e. $Z$ est bien une sous-
variété plongée (resp. une sous-variété) dans $X$.

PROPOSITION 7 - Supposons que $\mathcal{V}$ satisfasse l'axiome \underline{CV} 4), et de plus
l'axiome suivant :

\underline{\underline{CV} 5) Soient $X_i$ ($i=1,2$) deux $\mathcal{V}$-espaces, $Y_i$ ($i=1,2$) un sous-$\mathcal{V}$-espace fer-}
\underline{mé de $X_i$. Alors $Y_1 \times Y_2$ est un sous-$\mathcal{V}$-espace de $X_1 \times X_2$.}

Soient alors $X_i$ ($i=1,2$) deux $\mathcal{V}$-espaces, $Y_i$ ($i=1,2$) un $\mathcal{V}$-espace plon-
gé (resp. un sous- $\mathcal{V}$-espace) de $X_i$. Alors $Y_1 \times Y_2$ est un sous- $\mathcal{V}$-espace
plongé (resp. un sous- $\mathcal{V}$-espace) de $X_1 \times X_2$, et la structure induite sur
$Y_1 \times Y_2$ est identique à la structure produit des structures de $\mathcal{V}$-espaces
induites sur $Y_1$ par $X_1$ et sur $Y_2$ par $X_2$.

Le fait que $Y_1 \times Y_2$ admette une structure de $\mathcal{V}$-espace induite par
$X_1 \times X_2$, identique à la structure produit des structures induites envisa-
gées sur $Y_1$ et $Y_2$, résulte de la transitivité générale des structures
initiales (\emph{Ens.}, \emph{loc. cité}). Par hypothèse, l'assertion de la proposi-
tion est vraie si $Y_1$ et $Y_2$ sont des sous- $\mathcal{V}$-espaces fermés. Si $Y_i$ est
un sous- $\mathcal{V}$-espace quelconque, alors d'après prop. 5 c'est un sous- $\mathcal{V}$ -
espace fermé d'une partie ouverte $U_i$ de $X_i$. Donc $Y_1 \times Y_2$ est un sous- $\mathcal{V}$ -
espace de $U_1 \times U_2$, qui lui-même est un sous- $\mathcal{V}$-espace de $X_1 \times X_2$ (puisque
une partie ouverte de $X_1 \times X_2$), donc d'après prop. 6, $Y_1 \times Y_2$ est aussi un
sous- $\mathcal{V}$-espace de $X_1 \times X_2$. Supposons maintenant que $Y_i$ soit seulement
un sous- $\mathcal{V}$-espace plongé dans $X_i$ ($i=1,2$), montrons que $Y_1 \times Y_2$ est un
sous- $\mathcal{V}$-espace plongé dans $X_1 \times X_2$. Soit donc $(y_1,y_2) \in Y_1 \times Y_2$, il faut
prouver que tout germe de fonction régulière sur $Y_1 \times Y_2$ en $(y_1,y_2)$
provient d'un germe de fonction régulière sur $X_1 \times X_2$. Il suffit de

\newpage

%% ===== Page 73 =====

% typescript page 71 — n° 262

prouver ceci quand on remplace chaque $Y_i$ par une partie ouverte $U_i$
de $Y_i$ contenant $y_i$ ; or, en vertu de l'axiome $\underline{CV}\ 4)$, on peut trouver
un tel $U_i$ qui soit un sous-$\mathcal{V}$-espace de $X_i$. D'après ce qui précède,
$U_1 \times U_2$ est alors un sous-$\mathcal{V}$-espace de $X_1 \times X_2$, d'où aussitôt la con-
clusion voulue.

PROPOSITION 8 - Soit $X$ un $\mathcal{V}$-espace, $Y$ une partie de $X$. Si $X$ est une
sous-variété plongée (resp. une sous-variété) de $X$, alors pour toute
partie ouverte $U$ de $X$, $U \cap Y$ est une sous-variété plongée (resp. une
sous-variété ) dans $U$, et sa $\mathcal{V}$-structure induite par $U$ est iden-
tique à celle induite par $Y$ (dont $U \cap Y$ est une partie ouverte, pour
la topologie de $Y$ associée à sa $\mathcal{V}$-structure induite par $X$). Récipro-
quement, si $X$ admet un recouvrement ouvert $(U_i)$ tel que pour tout $i$,
$Y \cap U_i$ soit un $\mathcal{V}$-espace plongé (resp. un sous-$\mathcal{V}$-espace de $U_i$, alors
$Y$ est un sous-$\mathcal{V}$-espace plongé (resp. un sous-$\mathcal{V}$-espace de $X$),
$Y' = U \cap Y$ étant une partie ouverte de $Y$, est un sous-$\mathcal{V}$-espace de $Y$,
donc d'après prop.6 c'est un $\mathcal{V}$-espace plongé (resp. un sous-$\mathcal{V}$-espace
de $X$ et sa structure induite par $Y$ est la même que celle induite par $X$ .
Il reste à montrer que si un sous-$\mathcal{V}$-espace plongé (resp. un sous-
$\mathcal{V}$-espace)$Y'$ de $X$ est contenu dans un ouvert, c'est un sous-$\mathcal{V}$ -
espace plongé (resp. un sous-$\mathcal{V}$-espace) de ce dernier : alors les
structures induites sur $Y'$ par $X$ et $U$ sont les mêmes d'après prop.6
et la première partie de prop.7 sera prouvée. Or il est trivial à par-
tir des définitions que, $Y'$ étant une partie quelconque de l'ouvert $U$,
$Y'$ admet une structure induite dans $X$ si et seulement si elle admet
une structure induite dans $U$, et qu'alors $Y'$ est un sous-$\mathcal{V}$-espace
plongé (resp. un sous-$\mathcal{V}$-espace) dans $U$, si et seulement si il l'est
dans $X$.

\newpage

%% ===== Page 74 =====

% typescript page 72 — n° 262

dans $X$. Supposons maintenant donné un recouvrement ouvert $(U_i)$ de $X$, et
soit $Y$ une partie de $X$ telle que tout $Y \cap U_i$ ait une structure induite dans
$U_i$. D'après le par.1, n$^\circ$5, prop. 7, $Y$ admet une $\mathcal{V}$-structure locale in-
duite, et cette dernière est séparée puisque $X$ est séparé (par. 1, n$^\circ$ 4,
prop.6, corollaire) donc $Y$ est un $\mathcal{V}$-espace puisque $\mathcal{V}$ est parfait ou semi-
parfaite, et que dans ce dernier cas on pourrait supposer $(U_i)$ fini. Soit
$Y_i = Y \cap U_i$, alors la structure induite par $Y$ sur sa partie ouverte $Y_i$ est
identique à la structure induite sur $Y_i$ par $U_i$. Il résulte alors immédia-
tement des définitions que si chaque $Y_i$ est un sous-$\mathcal{V}$-espace plongé
dans $U_i$, $Y$ est un sous-$\mathcal{V}$-espace plongé dans $X$. Si de plus, $Y_i$ est un
sous-$\mathcal{V}$-espace de $U_i$, la topologie de $Y_i$ est celle induite par $U_i$, donc
celle induite par $X$, d'où résulte aussitôt que la topologie de $Y$ est aus-
si celle induite par $X$, donc que $Y$ est un sous-$\mathcal{V}$-espace de $X$.

PROPOSITION 9 - Soit $f$ un morphisme d'un $\mathcal{V}$-espace $Y$ dans un autre $X$.
Pour que $f$ soit un isomorphisme de $Y$ sur un $\mathcal{V}$-espace plongé (resp. un
sous-$\mathcal{V}$-espace) dans $X$, il faut et il suffit qu'il satisfasse aux con-
ditions suivantes :

(i) $f$ est injectif, (resp. $f$ est un homéomorphisme de $Y$ dans $X$).

(ii) Pour tout $y \in Y$, tout germe de fonction régulière en $y$ est l'ima-
ge inverse, par $f$, d'un germe de fonction régulière en $x=f(y)$.

(iii) (S'il s'agit de $\mathcal{V}$-espaces plongés) : toute application $g$ d'un
$\mathcal{V}$-espace $Z$ dans $Y$, telle que $fg$ soit un morphisme est \emph{continue}.

La nécessité de ces conditions est triviale. Pour la suffisance, il
faut seulement montrer que si $g$ est une application d'un $\mathcal{V}$-espace $Z$
dans $Y$, telle que $fg$ soit un morphisme, alors il en est de même de $g$.
Or $g$ est continu en vertu de (iii) (resp. en vertu de (ii)), de plus, il

\newpage

%% ===== Page 75 =====

% typescript page 73 — n° 262

résulte de (ii) que $g$ est un morphisme pour les structures $\underline{K}$-annelées
sous-jacentes à $Z$ et $Y$, donc c'est un morphisme en vertu de \underline{CV 3)}.

(N.B) La condition (iii) n'est pas surabondante, même si $Y$ est connexe
comme on voit par exemple dans le cas où $\mathcal{V}$ est la catégorie des varié-
tés différentiables : dans la définition "classique" des variétés plon-
gées par la condition que l'ensemble sous-jacent soit une partie de $X$,
et l'application d'injection $f:Y \to X$ satisfasse à (ii), i.e. soit de
rang égal à $\text{dim.} Y$ en tout point, \emph{l'ensemble des points de $Y$ ne suffit
plus à déterminer sa structure de variété plongée}. Cependant, la condi-
tion (iii) est vérifiée dans les cas importants en pratique, par exem-
ple, si $Y$ est une feuille maximale dans un feuilletage d'une variété
différentiable, définie par un champ complétement intégrable d'éléments
de contact - du moins si $X$ est paracompacte. On laisse à Bourbaki de
trancher la question terminologique ainsi soulevée.

4. - \emph{$\mathcal{V}$-espaces quotients}

Soit $X$ un $\mathcal{V}$-espace muni d'une relation d'équivalence $R$, soit
$Y=X/R$, rappelons (par.1, n°7, déf. 2) que $Y$ est appelé un $\mathcal{V}$-espace
quotient s'il existe sur $Y$ une structure de $\mathcal{V}$-espace dont la topolo-
gie sous-jacente soit la topologie quotient de celle de $X$, et telle que
pour tout ouvert $U \subset Y$, la $\mathcal{V}$-structure induite sur $U$ soit une $\mathcal{V}$-
structure quotient (au sens morphisque) de la $\mathcal{V}$-structure sur $f^{-1}(U)$
induite par $X$ ($f$ désignant l'application canonique de $X$ sur $Y$). En
particulier, les morphismes de $U$ dans $\underline{K}$ sont les applications $g:U \to \underline{K}$
telles que $gf: f^{-1}(U) \to \underline{K}$ soient des morphismes. D'où aussitôt, compte
tenu de \underline{CV 3)} :

\newpage

%% ===== Page 76 =====

% typescript page 74 - n° 262

PROPOSITION 10 - Soient $X$ un $\mathcal{V}$-espace muni d'une relation d'équivalence
$R$, $Y$ l'espace quotient $X/R$, $f:X \to Y$ l'application canonique, mettons sur
$Y$ la structure $K$-annelée quotient de celle de $X$, ($n^\circ 1$, prop.2, cor.1),
pour laquelle les fonctions régulières sur un ouvert $U \subset Y$ sont les fonc-
tions $g:U \to \underline{K}$ telles que $gf: f^{-1}(U) \to \underline{K}$ soit régulière. Pour que $Y$ soit
un $\mathcal{V}$-espace quotient de $X$, (par.1, $n^\circ 7$, déf.2) il faut et il suffit
que la structure précédente fasse de $Y$ un $\mathcal{V}$-espace, et dans ce cas
cette structure est aussi une $\mathcal{V}$-structure quotient.

A part ça, le rédacteur n'a rien à ajouter aux sorites de par.1,
$n^\circ 7$.

5. - \emph{Un axiome supplémentaire.}

Dans toute la suite de ce paragraphe, nous supposerons que la caté-
gorie de variétés $\mathcal{V}$ satisfait à l'axiome suivant :

\underline{CV 6)} \emph{Pour tout $\mathcal{V}$-espace $X$ et toute fonction $f$ régulière sur $X$ telle}
\emph{que $f(x) \neq 0$ pour tout $x \in X$, il existe une fonction régulière $g$}
\emph{sur $X$ telle que $fg = 1$.}

Cet axiome est équivalent au suivant :

\underline{CV 6')} \emph{$K$ est un corps, et l'application $x \to 1/x$ de l'ouvert $K-(0)$ dans $K$}
\emph{est une fonction régulière.}

CV 6) implique CV 6'), comme on voit en prenant pour $f$ l'applica-
tion identique de $K-(0)$ dans $K$, et la réciproque est vraie, si on note
qu'avec les notations de \underline{CV 6)}, la fonction $1/f$ est composée du morphis-
me $f: X \to K-(0)$ et du morphisme $x \to 1/x$ de $K-(0)$ dans $K$, donc un morphis-
me. Si $X$ est un $\mathcal{V}$-espace, nous désignerons par $\underline{O}_X$ le faisceau des ger-
mes de fonctions régulières sur $X$.

PROPOSITION 11 - Supposons que $\mathcal{V}$ satisfasse à l'axiome \underline{CV 6)}. Alors
l'anneau de base $k$ de $\mathcal{V}$ est un corps. Si $X$ est un $\mathcal{V}$-espace, les anneaux

\newpage

%% ===== Page 77 =====

% typescript page 75 - n° 262

$\underline{O}_{X,x} \quad (x\in X)$ sont des anneaux locaux, dont l'idéal maximal est le noyau de l'augmentation $\underline{O}_{X,x} \to K$.

Beweis klaar.

6. - \emph{$\mathcal{V}$-structures sur les espaces vectoriels}

Rappelons que dans toute la suite, nous supposerons \underline{CV 6)} satisfait

PROPOSITION 12 - Soient $V_0$ un espace vectoriel de dimension finie $n$ sur $k$, $V$ l'espace vectoriel sur $K$ déduit de $V_0$ par extension des scalaires, $e = (e_i)_{1 \le i \le n}$ une base de $V_0$, et $u_e : K^n \to V$ l'isomorphisme d'espaces vectoriels sur $K$ défini par $u_e((\lambda_i)) = \sum \lambda_i e_i$. Alors la $\mathcal{V}$-structure sur $V$ déduite de celle de $K^n$ par transport de structure à l'aide de $u_e$ ne dépend pas du choix de $e$, et $V_0$ s'identifie à l'ensemble des points fermés de $V$.

En effet, si $e'$ est une autre base de $V_0$, on a $u_{e'} = u_e v$, où $v$ est un automorphisme de $K^n$ défini par une matrice à coefficients dans $k$, et dont l'inverse a ses coefficients dans $k$. Il s'ensuit que $v$ et $v^{-1}$ sont des $\mathcal{V}$-morphismes de $K^n$ dans $K^n$ (prop. 4, cor 1) donc $v$ est un $\mathcal{V}$-automorphisme, ce qui prouve que $u_e$ et $u_{e'}$ définissent la même $\mathcal{V}$-structure sur $V$. Pour la dernière assertion, on peut supposer $V_0 = k^n$, d'où $V=K^n$. En vertu de \underline{CV 1)}, les points fermés dans un produit de $\mathcal{V}$-espaces sont les points dont les projections sont fermées, d'où s'ensuit que $k^n$ est l'ensemble des points fermés de $K^n$.

Par la suite, si $X$ est un $\mathcal{V}$-espace, nous désignerons par $X_{(k)}$ l'ensemble des points de $X$ qui sont fermés. Avec les notations précédentes on a donc $V_0 = V_{(k)}$. De façon générale, nous appelerons \emph{$\mathcal{V}$-espace vectoriel} un espace vectoriel $V$ de dimension finie sur $K$, muni d'une $\mathcal{V}$-structure telle que $V_{(k)}$ soit un sous-espace vectoriel sur $k$ de $V$, et

\newpage

%% ===== Page 78 =====

% typescript page 76 - n° 262

$V$ s'identifie (avec sa structure vectorielle et sa $\mathcal{V}$-structure) à
l'espace déduit de $V_{(k)}$ par extension des scalaires.

Corollaire 1 - Soient $V$ et $W$ deux $\mathcal{V}$-espaces vectoriels, pour toute
$u \in L_k(V_{(k)},W_{(k)})$ soit $u^K$ l'application $K$-linéaire de $V$ dans $W$ qui prolon-
ge $u$. Alors l'application $u \mapsto u^K$ est une bijection de $L_k(V_{(k)},W_{(k)})$
sur l'ensemble des applications $K$-linéaires de $V$ dans $W$ qui sont des
$\mathcal{V}$-morphismes.

On peut supposer $V=K^n$, $W=K^m$, alors $u^K$ est un $\mathcal{V}$-morphisme en vertu
de prop.4, cor.1, et réciproquement, si $v$ est une application $K$-linéaire
de $V$ dans $W$ qui est un $\mathcal{V}$-morphisme, elle transforme points fermés en
points fermés en vertu de \underline{CV} 1), donc applique $k^n$ dans $k^m$, donc est de
la forme $u^K$. En particulier :

Corollaire 2 - Si $V$ est un $\mathcal{V}$-espace vectoriel, le groupe des auto-
morphismes de $V$ s'identifie au groupe $GL_k(V_{(k)})$.

Le corollaire 1 montre qu'il n'y a pas d'inconvénient à identifier
la structure $\mathcal{V}$-espace vectoriel avec la structure de $K$-espace vecto-
riel de dimension finie \underline{défini sur} $k$ (i.e. muni d'un sous-espace vecto-
riel sur $k$ dont il se déduise par extension des scalaires), les notions
de morphisme pour ces deux structures étant en effet, les mêmes. Notons
alors que si $V$ et $W$ sont des espaces vectoriels de dimension finie défi-
nis sur $k$, il en est de même de $V'$, de $V\otimes W$, de $L(V,W)$, etc...
et on a

$$ (V')_{(k)}=(V_{(k)})' \quad , \quad (V\otimes_K W)_{(k)}=V_{(k)}\otimes_k W_{(k)} \quad , \quad L_K(V,W)_{(k)}=L_k(V_{(k)},W_{(k)}) $$

etc.... à cause des propriétés de compatibilité bien connues entre les
opérations $V \rightarrow V'$, $(V,W) \rightarrow V\otimes W$, $(V,W) \rightarrow L(V,W)$ etc.. avec l'opération
d'extension des scalaires. Aussi nous sous-entendrons toujours implicite-
ment que des espaces vectoriels comme les précédents, construits à partir

\newpage

%% ===== Page 79 =====

% typescript page 77 — n° 262

d'espaces vectoriels définis sur k, sont aussi considérés comme des
espaces vectoriels définis sur k.

PROPOSITION 13 - Soit V un $\mathcal{V}$-espace vectoriel. Alors le groupe $Gl(V)$
est une partie ouverte du $\mathcal{V}$-espace vectoriel $L(V,V)$, et c'est un $\mathcal{V}$-
groupe pour la $\mathcal{V}$-structure induite.

On peut supposer $V=k^n$. Comme $Gl(V)$ est la partie de $L(V,V)$ formée
des u tels que $\text{dét } u \neq 0$, que la fonction $\text{dét.}u$ est une application
polynômiale à coefficient dans k de $L(V,V) = \underline{K}^{n^2}$ dans $\underline{K}$, donc un morphis
me (prop. 4 cor.1), enfin que $\underline{K}-(0)$ est une partie ouverte de $\underline{K}$, il
s'ensuit que $Gl(V)$ est bien une partie ouverte. La multiplication
$G \times G \rightarrow G$ est induite par la multiplication dans $L(V,V)$, qui est une
application polynômiale à coefficients dans k de $\underline{K}^{n^2} \times \underline{K}^{n^2}$ dans $\underline{K}^{n^2}$ donc un morphisme
(prop.4, cor.1), il en est donc de même de $G \times G \rightarrow G$. Il est immédiat
que l'élément unité de G est un point fermé de G, puisque c'est un point
fermé de $L(V,V)$ (ses coordonnées étant toutes 0 ou 1, donc dans k).
Enfin montrons que l'application $x \rightarrow x^{-1}$ de G dans G est un morphisme,
cette application est de la forme $f(x)(\text{dét.}x)^{-1}$, où $x \rightarrow f(x)$ est une
application induite d'une application polynômiale, à coefficients
dans le corps premier, de $L(V,V)$ dans lui-même (à une matrice u corres-
pondant la matrice dont les éléments sont les déterminants des mineurs
correspondants de u). En vertu de prop.4, cor.1) f est un morphisme,
il en est de même de $x \rightarrow (\text{dét.}x)^{-1}$ en vertu de la même proposition et
de CV 6), donc aussi de $x \rightarrow f(x)(\text{dét.}x)^{-1}$, composée de morphismes
$G \rightarrow L(V,V) \times \underline{K} \rightarrow L(V,V)$ (le deuxième morphisme étant donné par
$(u, \lambda) \rightarrow \lambda u)$. C'est donc bien un morphisme de G dans $L(V,V)$, donc
aussi de G dans G. - Par suite, nous supposerons toujours implicitement

\newpage

%% ===== Page 80 =====

% typescript page 78 — n° 262

$Gl(V)$ muni de la structure de $\mathcal{V}$-groupe qu'on vient d'expliciter.

Corollaire - Muni du $\mathcal{V}$-groupe $Gl(V)$ d'opérateurs à gauche, $V$ devient un $\mathcal{V}$-espace à opérateurs (cf. par. 1, n°10).

En effet, l'application $G \times V \to V$ est un morphisme, car elle est induite par l'application $L(V,V) \times V \to V$ qui est un morphisme en vertu de prop. 4 cor. 1.

Soit $G$ un $\mathcal{V}$-groupe, $V$ un $\mathcal{V}$-espace vectoriel, on appelle \emph{$\mathcal{V}$-représentation} linéaire de $G$ dans $V$ un morphisme de $\mathcal{V}$-groupes de $G$ dans $Gl(V)$. Pour qu'une représentation linéaire $u$ de $G$ dans $V$ soit une $\mathcal{V}$-représentation, il faut et il suffit que ce soit un $\mathcal{V}$-morphisme de $G$ dans $L(V,V)$, i.e. que les fonctions coefficients $x \to \langle u(x).a, a' \rangle$ soient régulières sur $G$ pour $a \in V_{(k)}$, $a' \in V'_{(k)}$. L'application de prop. 4, cor. 1 prouve immédiatement :

PROPOSITION 14 - Une représentation linéaire d'un $\mathcal{V}$-groupe $G$, qui est la somme ou le produit tensoriel d'un nombre fini de $\mathcal{V}$-représentations, ou la contragrédiente d'une $\mathcal{V}$-représentation, est elle-même une $\mathcal{V}$-représentation. En particulier, si on a des $\mathcal{V}$-représentations de $G$ dans des $\mathcal{V}$-espaces vectoriels $V$ et $W$, on en déduit une $\mathcal{V}$-représentation dans le $\mathcal{V}$-espace vectoriel $L(V,W) = V' \otimes W$.

7. - \emph{$\mathcal{V}$-fibrés vectoriels. Relation avec les faisceaux de modules.}

Définition 5 - Soit $\underline{\mathcal{V}}_1$ la catégorie des $\mathcal{V}$-espaces vectoriels, $\underline{\Phi}$ la catégorie des espaces vectoriels de dimension finie sur $K$, et considérons les foncteurs identiques $\underline{\Phi} \to \underline{\mathcal{V}}_1$, $\underline{\Phi} \to \underline{\Phi}$ et $\underline{\Phi} \to \underline{\mathcal{E}}$ (où $\underline{\mathcal{E}}$ désigne la catégorie des ensembles). Un fibré de type ($\underline{\mathcal{V}}_1$, $\underline{\Phi}$, $\underline{\Phi}_0$) (cf. par. 1, n° 9, déf. 4) est appelé un \emph{$\mathcal{V}$-fibré vectoriel}. Un morphisme au sens de la catégorie des fibrés de type ($\mathcal{V}$, $\underline{\Phi}$, $\underline{\Phi}_0$) (loc. cité, déf. 5) est appelé un morphisme de \emph{$\mathcal{V}$-fibrés vectoriels}.

\newpage

%% ===== Page 81 =====

% typescript page 79 — n° 262

Ainsi, un $\mathcal{V}$-fibré vectoriel est défini par la donnée de deux $\mathcal{V}$-
espaces $X$ (la base) et $E$ (le fibré vectoriel), d'un morphisme surjectif
$p : E \to X$, et, pour tout $x \in X$, d'une structure d'espace vectoriel de dimen-
sion finie sur $K$ dans la fibre $p^{-1}(x)$, ces données étant assujetties à la
condition suivante : pour toute partie ouverte assez petite $U$ de $X$,
$p^{-1}(U)$ est isomorphe (au sens des structures précédentes, l'espace de base
étant $U$) à un produit $U \times V$, où $V$ est un $\mathcal{V}$-espace vectoriel.

PROPOSITION 15 - La donnée d'un $\mathcal{V}$-fibré vectoriel dont toutes les fibres
aient même dimension $n$, est équivalente à la donnée d'un $\mathcal{C}$-espace fibré
à groupe structural $\mathrm{Gl}(n, K)$, à fibre $K^n$ (considérée comme un $\mathcal{V}$-espace
à $\mathcal{V}$-groupe d'opérateurs $\mathrm{Gl}(V)$, cf. cor. à prop.13).

En effet, on est sous les conditions d'application du par.1, n$^\circ$ 10
prop.11. Il en résulte qu'on peut parler du \emph{fibré principal associé} à
un $\mathcal{V}$-fibré vectoriel, pourvu que toutes les fibres soient des vecto-
riels de même dimension. (On se demande s'il ne faut pas limiter le nom
de fibré vectoriel à ce cas). On peut aussi, dans ce cas, considérer le
$\mathcal{C}$-fibré \emph{adjoint}, qui est un $\mathcal{C}$-fibré en groupes tous isomorphes à
$\mathrm{Gl}(n, K)$, et dont les sections s'identifient aux automorphismes du fibré
vectoriel envisagé.

Rappelons qu'une structure de $\mathcal{V}$-fibré principal à $\mathcal{C}$-groupe
structural $G = \mathrm{Gl}(n, K)$, à fibre $K^n$, est par définition une structure de
fibré ensembliste à faisceau structural $\underline{O}_X(G)$ (faisceau des germes d'appli-
cations régulières de $X$ dans $G$), de type $(\underline{O}_X(G), G)$. Or on vérifie aussi-
tôt, à l'aide de \underline{CV} 6), qu'on a un isomorphisme naturel
$$
\underline{O}_X(\mathrm{Gl}(n, K)) \simeq \mathrm{Gl}(n, \underline{O}_X)
$$
le deuxième membre désignant le faisceau de germes de matrices inversible
dans le faisceau d'anneaux $\underline{O}_X$ ; et d'autre part $\mathrm{Gl}(n, \underline{O}_X)$ s'identifie

\newpage

%% ===== Page 82 =====

% typescript page 80 - n° 262

aussi au faisceau des germes d'automorphismes du faisceau $\underline{O}_X^n$ de $\underline{O}_X$-
modules. Donc une structure de $\underline{O}_X$-Module localement isomorphe à $\underline{O}_X^n$
peut aussi s'exprimer comme une structure de fibré ensembliste à fais-
ceau structural $\underline{Gl}(n,\underline{O}_X)$, de type $(\underline{Gl}(n,\underline{O}_X),\underline{O}_X^n)$. On en conclut, en
vertu de Top. Elém. chap.1 que, à condition de restreindre les morphismes
à être des isomorphismes, la catégorie des $\mathcal{V}$-fibrés vectoriels de base
fixée, $X$, à fibre isomorphe à $K^n$, est naturellement équivalente à la
catégorie des $\underline{O}_X$-modules localement libres, de rang n en chaque point.
Cette correspondance se précise de la façon suivante :

PROPOSITION 16 - Les $\mathcal{V}$-fibrés vectoriels de fibre $K^n$ sur $X$ correspon-
dent canoniquement aux faisceaux de $\underline{O}_X$-Modules localement isomorphes à
$\underline{O}_X^n$. Au $\mathcal{V}$-fibré vectoriel $E$ correspond le faisceau $\underline{O}_X(E)$ des germes
de sections régulières de $E$, et au faisceau $M$ de $\underline{O}_X$-modules correspond le
$\mathcal{V}$-fibré vectoriel unique $E(M)$ dont la fibre en $x \in X$ est
\[ E(M)_x = M_x \otimes_{\underline{O}_{X,x}} K = (M_x/\underline{m}_x M_x) \otimes_{\underline{k}_x} K \]
($\underline{m}_x$ désigne l'idéal maximal de $\underline{O}_{X,x}$, $\underline{k}_x$ son corps résiduel), et dont
les sections régulières sur un ouvert $U$ sont celles définies par les sec-
tions de $M$ sur $U$.

Dans le cas le plus important, on aura $\underline{k}_x = K$, d'où $E(M)_x = M_x/\underline{m}_x M_x$. -
On a oublié de dire que les \emph{$\mathcal{V}$-morphismes} sont aussi appelés \emph{applications
régulières}, d'où la notion de \emph{section $\mathcal{V}$-régulière} d'un $\mathcal{V}$-fibré. La
structure de $\underline{O}_X$-Module sur un faisceau $\underline{O}_X(E)$ se définit de façon évidente
le produit d'une fonction régulière $U \rightarrow K$ et d'une section régulière
de $E$ sur $U$ étant encore une section régulière de $E$ sur $U$, (il suffit de
vérifier cette assertion si $E$ est un produit $U \times K^n$, auquel cas elle est
conséquence triviale de prop.4, cor.1). On vérifie de même que si on
part d'un $\underline{O}_X$-module localement isomorphe à $\underline{O}_X^n$, il existe sur l'ensemble

\newpage

%% ===== Page 83 =====

% typescript page 81 — n° 262

$E(M)$ réunion des $M_x \otimes_{\underline{O}_{X,x}} K$ une structure unique de $\mathcal{V}$-fibré vectoriel
dont les sections régulières sur un ouvert arbitraire sont celles qui
proviennent de sections de $M$ : utilisant par. 1, n$^\circ$ 9, prop. 13 (qui
s'applique, $\mathcal{V}$ étant recollable), on est ramené à prouver cette asser-
tion dans le cas où $M = \underline{O}_X^n$, où la démonstration est laissée au lecteur.
Les deux correspondances envisagées dans prop.16 sont des foncteurs de
catégories locales sur $X$ et pour achever la démonstration il n'y a plus
qu'à montrer que les deux foncteurs composés sont isomorphes naturellement
aux foncteurs identiques. Le rédacteur n'a pas cru utile de copier ici
la démonstration qu'il en a faite pour lui-même par acquit de conscience.
Il avoue d'ailleurs volontiers que d'une part, l'énoncé de la proposition
16 est un peu vague (de sorte que d'aucuns ne sauront pas très bien ce
qu'il s'agit de démontrer), et que d'autre part, si on voulait le rendre
parfaitement précis, en explicitant les deux catégories, les deux fonc-
teurs, et les deux équivalences de foncteurs, et faisant toutes les véri-
fications requises pour satisfaire la définition (rédaction ultérieure ?)
d'une équivalence de catégories locales, il y aurait de quoi désespérer
les plus masochistes. Il se posera de plus en plus une question de rédac-
tion, et il semble que le parti raisonnable soit le suivant : donner au
stade fonctoriel général (Ensemble, et Top. élém. Chap.I) des définitions
parfaitement en forme, et aussi suffisamment de démonstrations parfaite-
ment correctes,  mais expliquer quelque part en laïus qu'on est obligé
en pratique, par la suite, de se dispenser d'un certain nombre de véri-
fications fastidieuses - tout comme on est obligé, en pratique, de ne pas
s'en tenir au langage formalisé de Ens. Chap.I, et d'utiliser des nota-
tions ambiguës.

\newpage

%% ===== Page 84 =====

% typescript page 82 — n° 262

\underline{Corollaire 1} - Les correspondances établies dans prop.16 sont des foncteurs locaux pour la notion de morphisme de $\mathcal{V}$-fibrés vectoriels de déf. 5, et les homomorphismes de $\underline{O}_X$-Modules.

On notera d'ailleurs que la proposition et son corollaire s'étendent pour donner une équivalence entre la catégorie des $\mathcal{V}$-fibrés vectoriels (à fibre de dimension non précisée) et la catégorie des $\underline{O}_X$-Modules localement libres (i.e. tels que pour tout ouvert U assez petit, existe un entier n tel que $M|U$ soit isomorphe à $\underline{O}_X^n|U$). Ces catégories sont d'ailleurs des catégories additives (Alg. Hom; ChapI), et les foncteurs envisagés sont des foncteurs additifs.

Soit E un $\mathcal{V}$-fibré vectoriel à fibre $K^n$. On appelle \emph{$\mathcal{V}$-repère linéaire} ou \emph{repère linéaire régulier}, de E, un système $(f_i)_{i \in I}$ de n sections régulières de E, telles que pour tout $x \in X$, les $f_i(x)$ soient des éléments linéairement indépendants de la fibre $E_x$, et en forment donc une base. Comme $\underline{O}_{X,x}$ est un anneau local et que $\underline{O}_X(E)_x$ est un modu-le libre de rang n sur celui-ci, il revient au même de dire que pour tout $x \in X$, les germes $f_{i,x} \in \underline{O}_X(E)_x$ forment une base du $\underline{O}_{X,x}$-module $\underline{O}_X(E)_x$, ou encore que les $f_i$, considérées comme sections du faisceau $\underline{O}_X(E)$, définissent un isomorphisme de $\underline{O}_X^I$ sur $\underline{O}_X(E)$, ou enfin que l'application $(x, (\lambda_i)_{i \in I}) \to \sum \lambda_i f_i(x)$ est un isomorphisme du $\mathcal{V}$-fibré trivial $X \times K^I$ sur E. Ainsi, pour un ensemble I fixé à n éléments, les \emph{$\mathcal{V}$-repères linéaires} $(f_i)_{i \in I}$ de E sont en correspondance biunivoque avec les isomorphismes de $X \times K^I$ sur E (ou encore, si on veut, avec les sections régulières du fibré principal associé à E quand $I=[1,n]$). En particulier, E admet des repères linéaires sur tout ouvert assez petit.

\newpage

%% ===== Page 85 =====

% typescript page 83 — n° 262

8. - \emph{Opérations usuelles sur les fibrés vectoriels.}

\emph{Fibrés associés à des fibrés vectoriels}. Les notations sont comme au n°7
Considérons les foncteurs contavariants $\underline{\Phi}_0 \to \underline{\Phi}_0$ et $\underline{\Phi} \to \underline{\Phi}$ associant
à un $\mathcal{V}$-espace vectoriel $V$ (resp. à un espace vectoriel de dimension
finie sur $K$) l'espace dual $V'$. On est dans les conditions du par.1, n°11
ce qui permet, à tout $\mathcal{V}$-fibré vectoriel $E$, d'associer le \emph{$\mathcal{V}$-fibré
vectoriel dual}, noté le plus souvent $E'$. Le fibré $E'$ est un foncteur par
rapport à $E$, comme on a vu dans loc. cité. On a de plus un isomorphisme
de foncteurs locaux.

\[
\underline{O}_X(E') = (\underline{O}_X(E))'
\]

où $E$ est un $\mathcal{V}$-fibré vectoriel de base $X$, et où pour tout $\underline{O}_X$-module $A$,
$A'$ désigne le $\underline{O}_X$-module des germes d'homomorphismes de $A$ dans $\underline{O}_X$. La véri-
fication de compatibilités comme la précédente est le type de celles que
Bourbaki doit se permettre de "laisser au lecteur".

On définit de façon toute analogue le produit tensoriel de deux ou
plusieurs $\mathcal{V}$-fibrés vectoriels, à cela près qu'il s'agit maintenant de
la construction de fibrés associés, non à un fibré, mais à un système
fini de fibrés ; il semble donc qu'il faille écrire le par.1, n°11 dans
ce cas plus général, ou du moins indiquer qu'on pourrait le faire, et
se borner "pour fixer les idées" au cas d'un seul fibré. On ne peut se
dispenser d'indiquer les propriétés d'associativité et de commutativité
du produit tensoriel de $\mathcal{V}$-fibrés, qui résultent des propriétés analo-
gues pour les produits tensoriels d'espaces vectoriels, en vertu de ce
qui a été dit dans loc. cité à propos des homomorphismes fonctoriels
entre fibrés associés. On a encore des isomorphismes de bifoncteurs
locaux :

\[
\underline{O}_X(E \otimes F) = \underline{O}_X(E) \otimes_{\underline{O}_X} \underline{O}_X(F)
\]

\newpage

%% ===== Page 86 =====

% typescript page 84 — n° 262

qui se généralisent au produit tensoriel d'un nombre fini quelconque
de $\mathcal{V}$-fibrés vectoriels sur l'espace $X$.

On définit de même, pour deux $\mathcal{V}$-fibrés vectoriels $E,F$ sur $X$, le
$\mathcal{V}$-fibré vectoriel $L(E,F)$, qu'il vaut mieux noter $E' \otimes F$. On a
$$ \underline{O}_X(E' \otimes F) = \underline{\mathrm{Hom}}_{\underline{O}_X} (\underline{O}_X(E), \underline{O}_X(F)) $$
(homomorphismes de bifoncteurs locaux), où pour deux $\underline{O}_X$-modules $A$ et $B$,
on désigne par $\underline{\mathrm{Hom}}_{\underline{O}_X}(A,B)$ le faisceau des germes des homomorphismes de
$A$ dans $B$.

On définit de même les puissances extérieures d'un $\mathcal{V}$-fibré vec-
toriel, ses puissances symétriques etc, toutes les formules tensorielles
usuelles s'étendent en des formules anlogues pour les $\mathcal{V}$-fibrés. Proxi-
mus Redactordemerdatur.

Ce serait aussi ici le moment de poser la définition de $L^{\otimes n}$ pour
tout $n \in \mathbf{Z}$, $L$ étant un $\mathcal{V}$-fibré à fibre de dimension $1$ : On pose pour
$n \geqslant 0$, $L^{\otimes n} = L \otimes L \dots \otimes L$ ($n$ facteurs), et pour $n \leqslant 0$, $L^{\otimes n} = (L')^{\otimes (-n)}$.
Avec cette définition on a $L^{\otimes (0)} = \underline{1} , L^{\otimes n} \otimes L^{\otimes m} = L^{\otimes (n+m)}$.

Les opérations précédentes passent aux classes de fibrés vectoriels
sur une base fixée $X \in \mathcal{V}$ . Soit $\underline{E}_X(n)$ ou simplement $\underline{E}(n)$ l'ensemble
$H^1(X, GL(n, \underline{O}_X))$ des classes de $\mathcal{V}$-fibrés vectoriels à fibre $K^n$ sur $X$, et
soit $\underline{E}_X$, ou simplement $\underline{E}$ l'ensemble somme des $\underline{E}(n)$. Alors $\underline{E}$ est muni
de deux lois de composition $\oplus$ et $\otimes$ commutatives et associatives, la
deuxième étant distributive par rapport à la première. Il est alors
commode d'introduire un anneau commutatif $\underline{A}_X$, engendré par le système
de générateurs $\underline{E}$, ses générateurs étant soumis aux relations :
$$ (\xi \oplus \eta) = (\xi) + (\eta) \,,\, (\xi \otimes \eta) = (\xi).(\eta) . $$

PROPOSITION 17 - (Atiyah). Soit $X \in \mathcal{V}$, supposons que pour tout $\mathcal{V}$-fibré

\newpage

%% ===== Page 87 =====

% typescript page 85 - n° 262

vectoriel $E$ sur $X$, l'espace des sections régulières de $E$ sur $X$ soit de
dimension finie sur $k$. Alors tout $\mathcal{V}$-fibré vectoriel $E$ sur $X$ est iso-
morphe à la somme directe d'une suite de fibrés indécomposables, bien
déterminés à isomorphisme et l'ordre des facteurs près.

Cette proposition est un cas particulier d'une variante du théorème
de Krull-Remak-Schmidt, valable dans les catégories additives (Alg.Hom...)
(et qui résulte d'ailleurs immédiatement du th. de KRS lui-même). Il
suffit de noter que l'hypothèse implique que pour tout $\mathcal{V}$-fibré vecto-
riel $E$, le $k$-espace vectoriel des endomorphismes de $E$ est de dimension
finie (puisque c'est l'espace des sections régulières de $E' \otimes E$), donc un
anneau d'Artin.

\underline{Corollaire} - Sous les conditions précédentes, l'anneau $\underline{A}_X$ engendré par
les classes de $\mathcal{V}$-fibrés vectoriels homogènes sur $X$ (cf. ci-dessus), a
pour groupe additif sous-jacent le groupe libre engendré par les classes
des fibrés vectoriels indécomposables homogènes.

La démonstration est immédiate.

L'anneau $\underline{A}_X$ sera appelé par abus de langage \emph{l'anneau des fibrés vec-
toriels sur $X$}. Sauf dans un petit nombre de cas extrêmement spéciaux, on
ne sait à peu près rien sur cet anneau. Notons cependant que les classes
de $\mathcal{V}$-fibrés sur $X$ à fibre de dimension 1 formant un groupe commutatif
$\underline{P}_X$ sous l'opération $\otimes$, qui s'identifie à un sous-groupe du groupe des
éléments inversibles de $\underline{A}_X$. Donc il y a un homomorphisme canonique de
l'algèbre $Z(\underline{P}_X)$ du groupe $\underline{P}_X$ dans $\underline{A}_X$, qui est d'ailleurs injectif quand
les conditions de la prop. 17 (donc de son corollaire) sont remplies.

\underline{Images réciproques de fibrés vectoriels}. La notion d'image réciproque de
$\mathcal{V}$-fibrés vectoriels est claire, étant un cas particulier de par.1 n°13

\newpage

%% ===== Page 88 =====

% typescript page 86 — n° 262

En vertu de la compatibilité de l'opération "image réciproque" et de
l'opération "fibré associé", les opérations tensorielles envisagées
plus haut sur les fibrés vectoriels, sont compatibles avec l'opéra-
tion image réciproque. Il faut relier enfin la notion d'image récipro-
que de fibré vectoriel à une notion d'image réciproque de faisceau de
modules, comme suit.

Soit $f$ un morphisme d'un $\mathcal{V}$-espace $X$ dans un autre $Y$, soit $B$ un
$\underline{O}_Y$-Module. Alors l'image réciproque $f^{-1}(B)$ (Top.Elém. Chap.I) est un
$f^{-1}(\underline{O}_Y)$-Module. D'autre part, on a un homomorphisme canonique de fais-
ceaux d'anneaux
$$f^{-1}(\underline{O}_Y) \to \underline{O}_X$$
(qui est d'ailleurs un monomorphisme), permettant de considérer $\underline{O}_X$
comme une $f^{-1}(\underline{O}_Y)$-Algèbre, ce qui permet de former le $\underline{O}_X$-Module
$$B^f = f^{-1}(B) \otimes_{f^{-1}(\underline{O}_Y)} \underline{O}_X$$

\underline{Définition 6} -. Sous les conditions précédentes, le faisceau $B^f$ est
appelé le $\underline{O}_X$-Module image réciproque du $\underline{O}_Y$-Module $B$ par $f$.

Bien entendu, $B^f$ est un foncteur additif local en $B$. On a
$$(\underline{O}_Y)^f = \underline{O}_X$$
ce qui montre que si $B$ est localement isomorphe à $\underline{O}_Y$, alors $B^f$ est
localement isomorphe à $\underline{O}_X$. Si maintenant $E$ est un $\mathcal{V}$-fibré vectoriel
sur $Y$, on a un isomorphisme naturel
$$\underline{O}_X(f^{-1}(E)) = \underline{O}_Y(E)^f$$
En effet, la notion d'image réciproque d'une section régulière permet
de définir de façon évidente un homomorphisme de faisceaux
$f^{-1}(\underline{O}_Y(E)) \to \underline{O}_X(f^{-1}(E))$, qui est un homomorphisme de $f^{-1}(\underline{O}_Y)$-Modules,
et définit donc (le deuxième terme étant un $\underline{O}_X$-Module) un homomorphisme

\newpage

%% ===== Page 89 =====

% typescript page 87 — n° 262

$f^{-1}(\underline{O}_Y(E)) \otimes_{f^{-1}(\underline{O}_Y)} \underline{O}_X = \underline{O}_Y(E)^f \rightarrow \underline{O}_X(f^{-1}(E))$. Pour montrer que c'est un
isomorphisme, on est ramené aussitôt au cas où $E$ est un produit $X \times K^n$,
puis au cas où $E$ est le produit $X \times K$, auquel cas c'est trivial.

\underline{Exercice} - Soit $X \in \mathcal{V}$, soit $A_X$ l'anneau des $\mathcal{V}$-fibrés vectoriels
sur $X$. Montrer qu'il existe un homomorphisme unique $r$ de $A_X$ sur $\mathbb{Z}$,
tel que $r$ prenne la valeur $n$ sur l'image canonique de l'ensemble
$H^1(X, \mathrm{Gl}(n, \underline{O}_X))$, dans $A_X$. Soit $P_X = H^1(X, \mathrm{Gl}(1, \underline{O}_X))$ l'ensemble des
classes de $\mathcal{V}$-fibrés vectoriels sur $X$, à fibre de dimension $1$, muni
de sa structure de groupe abélien, écrite additivement. Sur le grou-
pe abélien $\mathbb{Z} \times P_X$, mettons la structure d'anneau pour laquelle $P_X$
est un idéal de carré nul, le quotient par cet idéal étant $\mathbb{Z}$, muni
de sa structure d'anneau usuelle. Pour $\xi \in H^1(X, \mathrm{Gl}(n, \underline{O}_X))$, soit
$c^1(\xi)$ son image dans $H^1(X, \mathrm{Gl}(1, \underline{O}_X))$ par l'application qui corres-
pond à l'homomorphisme de faisceaux de groupes $\mathrm{Gl}(n, \underline{O}_X) \rightarrow \mathrm{Gl}(1, \underline{O}_X)$
donné par $u \rightarrow \text{dét.} u$. (Si $\xi$ est la classe d'un fibré vectoriel $E$
à fibre $K^n$, $c^1(\xi)$ est la classe du fibré vectoriel $\bigwedge^n E$ à fibre $K$).
Montrer qu'il existe un homomorphisme d'anneaux et un seul de $A_X$
dans $\mathbb{Z} + P_X$ qui transforme l'image d'un $\xi \in H^1(X, \mathrm{Gl}(n, \underline{O}_X))$ en
$(n, c^1(\xi))$. Montrer que la définition de l'homomorphisme précédent
se généralise quand on se donne seulement un espace topologique $X$
muni d'un faisceau $\underline{O}$ d'anneaux commutatifs avec unité.

9. - \underline{Sous-$\mathcal{V}$-fibrés vectoriels}, \underline{$\mathcal{V}$-fibrés vectoriels quotients}, 
\underline{$\mathcal{V}$-fibrés en drapeaux}.

Soit $E$ un $\mathcal{V}$-fibré vectoriel de base $X$. On appelle sous-$\mathcal{V}$-fibré
vectoriel de $E$ un -fibré vectoriel $F$ de base $X$, qui est en tant que

\newpage

%% ===== Page 90 =====

% typescript page 88 — n° 262

$\mathcal{V}$-espace un sous-$\mathcal{V}$-espace de $E$, dont l'application de projection
est induite par celle de $E$, dont les fibres sont des sous-espace vecto-
riels des fibres de $E$, munis de la structure $\underline{K}$-vectorielle induite, et
qui est tel que pour tout ouvert $U \subset X$ assez petit, il existe un isomor-
phisme de $E|U$ sur $U \times K^n$ appliquant $F \cap (E|U)$ sur $U \times K^m$ (où, pour $m \le n$, on
identifie $K^m$ à un sous-espace vectoriel de $K^n$ de la façon bien connue).
Notons d'abord qu'il résulte de cette définition qu'un \emph{sous-$\mathcal{V}$-fibré
vectoriel $F$ de $E$ est déterminé quand on connait la partie de $E$ qu'elle
définit, et qu'une partie $F$ de $E$ provient d'un sous-$\mathcal{V}$-fibré vectoriel
si et seulement si elle satisfait à la condition de la fin de la défini-
tion précédente.} Ces assertions sont triviales, sauf celle-ci : si $F$ est
une partie de $E$ telle que pour tout ouvert $U$ assez petit dans $X$, existe
un isomorphisme $\varphi_U$ de $E|U$ sur $U \times K^n$ appliquant $F \cap (E|U)$ sur $U \times K^m$, alors
$F$ est un sous-$\mathcal{V}$-fibré vectoriel. On voit aussitôt qu'il suffit de mon-
trer qu'alors $F$ est un sous-$\mathcal{V}$-espace de $E$, et que la $\mathcal{V}$-structure in-
duite sur $F \cap (E|U)$ s'identifie, par $\varphi_U$, à celle de $U \times K^m$.
Grâce au n°8, prop. 8, on est ramené à prouver que $U \times K^m$ est un sous-$\mathcal{V}$-
espace de $U \times K^n$. Or l'application d'injection $U \times K^m \to U \times K^n$ est évidemment
un morphisme, et il existe un morphisme inverse à droite, déduit de la
projection canonique de $K^n$ sur son sous-espace $K^m$, d'où notre assertion.

On appelle \emph{$\mathcal{V}$-fibré vectoriel quotient} de $E$ un $\mathcal{V}$-fibré vectoriel $G$
tel que $G$ soit un $\mathcal{V}$-espace quotient de $E$, l'application canonique $E \to G$
étant compatible avec les projections sur la base $X$, linéaire sur les
fibres, et la condition suivante étant satisfaite : pour tout ouvert $U \subset X$
assez petit, il existe un diagramme commutatif de $\mathcal{V}$-fibrés vectoriels

\begin{tikzcd}
E|U \arrow[r] \arrow[d] & G|U \arrow[d] \\
\phantom{U \times K^n} & \phantom{U \times K^p}
\end{tikzcd}

\newpage

%% ===== Page 91 =====

% typescript page 89 - n° 262

les flèches verticales étant des isomorphismes, et la deuxième flèche
horizontale étant l'application déduite de la projection naturelle
$(x_1, \dots, x_n) \to (x_{m+1}, \dots, x_n)$ de $\underline{K}^n$ sur $\underline{K}^{n-m}$. En tout $x \in X$, la fibre
$G_x$ de $G$ est alors un espace vectoriel quotient de la fibre $E_x$ de $E$ en $x$,
donc de la forme $E_x/F_x$, où $F_x$ est un sous-espace vectoriel bien déter
miné de $E_x$. En vertu de ce qui précède, la réunion $F$ des $F_x$ est un sous
$\mathcal{V}$-fibré vectoriel de $E$. On a alors le résultat : \emph{le $\mathcal{V}$-fibré vectoriel
quotient $G$ de $E$ est bien déterminé par la donnée du sous- $\mathcal{V}$-fibré $F$ de
$E$ qui lui correspond, et tout sous- $\mathcal{V}$-fibré vectoriel $F$ de $E$ provient
d'un $\mathcal{V}$-fibré vectoriel quotient de $E$ (noté $E/F$).} Seul, le fait que tout
sous- $\mathcal{V}$-fibré vectoriel $F$ de $E$ correspond à un $\mathcal{V}$-fibré quotient deman
de une démonstration. Soit donc $G$ l'ensemble quotient de $E$ défini par $F$,
il faut montrer que $G$ est un $\mathcal{V}$-espace quotient de $E$, et que muni de la
projection naturelle $G \to X$ et des structures vectorielles quotient sur
les fibres, c'est un $\mathcal{V}$-fibré vectoriel sur $X$ (il résultera alors de la
définition que $G$ est un $\mathcal{V}$-fibré vectoriel quotient de $E$). Or cette asser
tion est vraie à condition de se restreindre au dessus d'un ouvert
assez petit $U$ de $X$, comme on voit en procédant comme plus haut, en
utilisant l'injection naturelle de $\underline{K}^{n-m}$ dans $\underline{K}^n$, inverse à droite de la
projection $\underline{K}^n \to \underline{K}^{n-m}$. En vertu de par.1, n°7 prop.9, on en conclut que
$G$ est un $\mathcal{V}$-espace quotient de $E$, induisant sur les ouverts $G|U$ les
$\mathcal{V}$-structures qu'on vient d'envisager. Comme ces dernières, quand on
considère les $G|U$ comme des fibrés vectoriels sur $U$, se recollent, il
s'ensuit en vertu de par.1, n°9, prop.13 (qui s'applique, puisque $\mathcal{V}$ est
recollable) que $G$ est même un $\mathcal{V}$-espace, et un $\mathcal{V}$-fibré vectoriel sur $X$.

Ainsi les sous-$\mathcal{V}$-fibrés et les $\mathcal{V}$-fibrés quotients de $E$ se

\newpage

%% ===== Page 92 =====

% typescript page 90 — n° 262

correspondent biunivoquement. On voit de même, si $E'$ désigne le $\mathcal{V}$-fibré
vectoriel dual de $E$, que les sous- $\mathcal{V}$ -fibrés vectoriels $F$ de $E$ et de $E'$
se correspondent biunivoquement par polarité, et qu'on a encore les iso-
morphismes habituels
$$F' = E'/F^\circ \, , \quad (E/F)' = F^\circ$$

Ce serait le moment ici de définir les \emph{$\mathcal{V}$-fibrés de drapeaux}. (De
façon général, le présent paragraphe semble le contexte naturel pour dé-
rouler les sorites géométriques sur des fibrés spéciaux divers dans le
genre de ceux du présent numéro, par exemple sur les fibrés orthogonaux
sympléctiques etc..). Soit $\nu = (n_1, \dots, n_k)$ une suite strictement croissan-
te d'entiers $>0$, avec $n_k = n$. Un \emph{drapeau de type $\nu$} (sur l'espace vecto-
riel $K$) est un $K$-espace vectoriel $V$ de dimension $n$, muni d'une suite
croissante $(V_i)_{1 \le i \le k}$ des sous-espaces vectoriels de dimension $n_i$. On
définit de même la notion de drapeau de type $\nu$ sur $K$ \underline{défini sur $k$}, encore
appelé \emph{$\mathcal{V}$-drapeau de type $\nu$}. Un \emph{morphisme} d'un drapeau $(V_i)$ de type $\nu$
dans un autre $(W_i)$ est une application $K$-linéaire $V_k \xrightarrow{u} W_k$ appliquant
chaque $V_i$ dans $W_i$ ; on définit de même la notion de morphismes de $\mathcal{V}$-
drapeaux. D'où les catégories $\underline{\mathcal{D}}$ et $\underline{\mathcal{D}}_0$ permettant de construire la ca-
tégorie des fibrés de type $(\mathcal{C}, \underline{\mathcal{D}}, \underline{\Phi}, \underline{\mathcal{D}}_0)$ (par. 1, n° 9). Ces fibrés sont
appelés \emph{$\mathcal{V}$-fibrés en drapeaux de type $\nu$}. Considérons alors le $\mathcal{V}$-
drapeau de type $\nu$
$$D(\nu) = (K^{n_1}, K^{n_2}, \dots, K^{n_k})$$
et le sous-groupe $T\ell(\nu)$ de $G\ell(n,K)$ formé des $K$-automorphismes de $K^n$
qui invarient $D(\nu)$. Ce sous-groupe est l'intersection de $G\ell(n,K)$ avec
un sous-espace vectoriel défini sur $k$ de $L(K^n,K^n)$. C'est donc un sous-
$\mathcal{V}$-groupe de $G\ell(n,K)$, et en particulier un $\mathcal{V}$-groupe, et $K^n$ peut être

\newpage

%% ===== Page 93 =====

% typescript page 91 - n° 262
regardé comme un $\underline{\mathcal{V}}$-espace à $\underline{\mathcal{V}}$-groupe d'opérateurs $T\ell(\nu)$. Je dis que
\emph{la structure de $\underline{\mathcal{V}}$-fibré en drapeaux de type $\nu$ est équivalente à la
structure de $\underline{\mathcal{V}}$-fibré à $\underline{\mathcal{V}}$-groupe structural $T\ell(\nu)$ et fibre $\underline{K}^n$} (par .1
n$^\circ$ 10). En effet, $T\ell(\nu)$ est le groupe des $\underline{\Phi}$-automorphismes de $D(\nu)$,
et on peut appliquer par.1, n$^\circ$ 10 prop.14. On peut encore dire que, si
E est un $\underline{\mathcal{V}}$-fibré vectoriel à fibre -type $\underline{K}^n$, \emph{la donnée sur E d'une struc-
ture de $\underline{\mathcal{V}}$-fibré en drapeaux de type $\nu$ , compatible avec sa structure de
$\underline{\mathcal{V}}$-fibré vectoriel, équivaut à la donnée d'une restriction de son $\underline{\mathcal{V}}$-
groupe structural $G\ell(n,\underline{K})$ à $T\ell(\nu)$} (voir fin de par.1 n$^\circ$11 bis). Pour
pouvoir appliquer le critère usuel de restrictabilité du groupe structu-
ral, il nous faut des précisions sur la fibration de $G\ell(n,\underline{K})$ par $T\ell( \ )$
De façon générale, soit V un $\underline{K}$-espace vectoriel de dimension n, on déno-
te par $Dr(\nu,V)$ l'ensemble des drapeaux de type $\nu$ dans V, par $Gr(m,V)$
l'ensemble des sous-espaces vectoriels de dimension m de V, (qui s'iden-
tifie aussi à l'ensemble des drapeaux de type (m,n)) par $P(V)$ l'ensemble
$Gr(1,V)$ des droites de V. On a des applications naturelle
(1) \quad $Dr(\nu,V) \longrightarrow \prod_{i=1}^k Gr(n_i,V) \longrightarrow \prod_{i=1}^k P(\overset{n_i}{\bigwedge} V)$

où la deuxième flèche est déduite des applications naturelle $Gr(m,V) \longrightarrow$
$P(\overset{m}{\bigwedge} V)$, faisant correspondre à un sous-espace vectoriel W de dimension
m de V l'ensemble des multivecteurs $x_1 \wedge \dots \wedge x_m$ ($x_1,\dots,x_m \in W$). D'autre
part, le groupe $G\ell(V)$ opère dans tous les ensembles précédents, et les
applications (1) sont compatibles avec ces opérations. D'ailleurs
$Dr(\nu,V)$ (et en particulier les $Gr(m,V)$) sont des espaces homogènes sous
$G\ell(V)$.

\underline{Théorème 1} - Supposons que V soit un $\underline{\mathcal{V}}$-espace vectoriel de dimension n.
a) Soit D un $\underline{\mathcal{V}}$-drapeau de type $\nu$ dans V, T son stabilisateur dans

\newpage

%% ===== Page 94 =====

% typescript page 92 - n° 262

$G = G\ell(V)$. Alors $T$ est un sous-$\mathcal{V}$-groupe fermé dans $G$, et $G$ fibré par $T$ à
droite est localement trivial (par. 1, n$^\circ$ 11 bis, déf. 7), par suite $G/T$ est un $\mathcal{V}$-
espace quotient de $G$. Ainsi, par transport de structure à partir de la bijection
$s \mapsto s.D$ de $G/T$ sur $Dr(\nu, V)$, ce dernier ensemble se trouve muni d'une struc-
ture de $\mathcal{V}$-espace, d'ailleurs indépendante du choix de $D$, qui en fait au $\mathcal{V}$-espace
à $\mathcal{V}$-groupe d'opérateurs à gauche $G$. Ceci s'applique en particulier aux variétés
grassmaniennes $Gr(m,V)$, et en particulier à $Gr(1,V) = P(V)$. b) la $\mathcal{V}$-structure
ainsi mise sur $P(V)$ induit sur les sous-espaces affines de $P(V)$ définis par les
hyperplans de $V$ définis sur $k$, une $\mathcal{V}$-structure isomorphe à celle de $K^{n-1}$ par
les applications linéaires affines définies sur $k$ de $U$ sur $K^{n-1}$ (ce qui suffit
à caractériser cette $\mathcal{V}$-structure). c) Les applications (1) sont des $\mathcal{V}$-isomor-
phismes sur des sous-$\mathcal{V}$-espaces fermés.

La démonstration de ce théorème ne sera pas donnée ici, elle est bien enten-
du tout à fait élémentaire, le th. 1 est en réalité un énoncé de Géométrie Algé-
brique : il faut construire trois $\mathcal{V}$-morphismes, on construit pour cela des appli-
cations polynomiales définies sur $k$. L'explication a priori du th. 1 (valable pour
toute $\mathcal{V}$) est dans le fait que toutes les variétés algébriques qui y interviennent
sont des variétés rationnelles, et même dont tout point admet un voisinage isomor-
phe à un ouvert d'un $K^m$, et pour cette raison sont munies de $\mathcal{V}$-structures de
façon naturelle, quel que soit $\mathcal{V}$. (Ce qui n'aurait pas été possible si elles
n'avaient pas été rationnelles, puisque les $(k,K)$-ensembles algébriques non sin-
guliers dont les composantes connexes sont des variétés rationnelles sur $k$ forment
une "catégorie de variétés sur $K$", dont le "corps de définition" est $k$).

Nous pouvons maintenant appliquer par. 1, n$^\circ$ 11 bis, prop. 15 bis, cor. 2, et
nous obtenons; compte tenu du th. 1, a) :

Corollaire 1- Soit $E$ un $\mathcal{V}$-fibré vectoriel sur $X$, à fibre-type $K^n$, soit $P$ le fibré

\newpage

%% ===== Page 95 =====

% typescript page 93 - n° 262

principal associé, à groupe structural $Gl(n,K)$. Alors $P/T\ell(\nu)$ est un $\mathcal{V}$-
espace quotient de P et est le fibré associé à E et la représentation na-
turelle de $G\ell(n,K)$ dans la variété de drapeaux $Dr(\nu,K^n)$. Ses $\mathcal{V}$-sections
s'identifient aux restrictions du groupe structural de E à $T\ell(\nu)$, i.e.
aux structures de $\mathcal{V}$-fibré en drapeaux de type $\nu$ sur E, compatibles avec
la structure de $\mathcal{V}$-fibré vectoriel.

En particulier, supposons $\nu = (m,n)$, alors il résulte aussitôt des
définitions que la donnée d'une structure de $\mathcal{V}$-fibré en drapeaux de
type $\nu$ sur E, compatible avec sa structure de $\mathcal{V}$-fibré vectoriel, est
équivalente à la donnée d'un sous-$\mathcal{V}$-fibré vectoriel F à fibre type $K^m$.
Par suite :

Corollaire 2 - La donnée d'un sous-$\mathcal{V}$-fibré vectoriel F de E, à fibre-
type $K^m$, est équivalente à la donné d'une $\mathcal{V}$-section du $\mathcal{V}$-fibré asso-
cié à E, de fibre la variété grassmanienne $Gr(m,K^n)$.

Notons maintenant ($\nu$ étant de nouveau quelconque) qu'en vertu du
th.1, c), une application u de X dans $Dr(\nu,V)$ est un $\mathcal{V}$-morphisme si et
seulement si il en est de même des applications correspondantes de X
dans les $Gr(n_i,V)$. On en conlut, en conjuguant cor.1 et cor.2 :

Corollaire 3 - Une structure de $\mathcal{V}$-fibré en drapeaux de type $\nu$ est équi-
valente à une structure de $\mathcal{V}$-fibré vectoriel E à fibre type $K^n$, muni
d'une suite croissante $(E_1,\dots,E_k)$ de sous-$\mathcal{V}$-fibrés vectoriels à fibres
types $n_1,\dots,n_k$.

10. - \underline{Extensions de fibrés vectoriels.}

PROPOSITION 18 - Soit $u: E \to F$ un homomorphisme de $\mathcal{V}$-fibrés vectoriels
de base X, et pour tout $x \in X$, soit $u_x$ l'application K-linéaire $E_x \to F_x$
définie par u. Les conditions suivantes sont équivalentes :

a) $Rang(u_x)$ est une fonction localement constante de x.

\newpage

%% ===== Page 96 =====

% typescript page 94 — n° 262

b) La réunion des $\mathrm{Ker}(u_x)$ est un sous-$\mathcal{V}$-fibré vectoriel de $E$.

b') La réunion des $\mathrm{Im}(u_x)$ est un sous-$\mathcal{V}$-fibré vectoriel de $F$.

c) On peut factoriser $u$ en le composé d'homomorphismes
$$ E \xrightarrow{\alpha} E_1 \xrightarrow{u'} F_1 \xrightarrow{\beta} F $$
où $\alpha$ est l'homomorphisme canonique de $E$ sur un $\mathcal{V}$-fibré vectoriel quotient
$\beta$ l'homomorphisme canonique d'un sous-$\mathcal{V}$-fibré vedoriel de $F$ dans $F$,
et $u_1$ un isomorphisme de $\mathcal{V}$-fibrés vectoriels.

\underline{Démonstration.} Evidemment c) implique b) et b'), qui chacun impli-
quent a), reste à prouver que a) implique c). On vérifie aussitôt que la
question est locale, et qu'on peut donc supposer $E=X \times K^n$, $F=X \times K^m$. Alors
$u$ est définie par une application régulière $x \to u(x)$ de $X$ dans $L(K^n,K^m)$

Prouvons d'abord le

\underline{Corollaire 1} - Soient $V$ et $W$ deux $\mathcal{V}$-espaces vectoriels, $u$ une applica-
tion régulière de $X$ dans $L(V;W)$, telle que $u(x)$ soit de rang $r$ pour tout
$x \in X$, $u_r$ une application linéaire de rang $r$ de $V$ dans $W$. Il existe alors
pour tout ouvert $U \subset X$ assez petit des applications régulières $v$ de $X$
dans $G\mathcal{L}(V)$ et $w$ de $X$ dans $G\mathcal{L}(W)$ telles que l'on ait $u(x)=w(x)u_r v(x)^{-1}$

Démontrons le corollaire. Il est à peu près trivial si pour tout
$x \in X$, $u(x)$ a même noyau et même image que $u_r$ : prenant un supplémentaire
$W_1$ de $\mathrm{Im} \, u_r$ dans $W$, on peut alors écrire de façon unique $u(x)=w(x)u_r$,
où $w(x)$ est un automorphisme de $W$ induisant l'identité sur $W_1$, et
$x \to w(x)$ sera manifestement une application régulière de $X$ dans $G\mathcal{L}(W)$

Par suite, il suffit de trouver $v(x)$ et $w(x)$ de telle façon que
$w(x)u_r v(x)^{-1}$ ait même noyau et même image que $u_r$, ou encore de touver
des applications régulières $v$ et $w$ telles que pour tout $x$, $\mathrm{Ker} \, u(x)$ soit
transformé de $\mathrm{Ker} \, u_r$ par $v(x)$, et $\mathrm{Im} \, u(x)$ soit transformé de $\mathrm{Im} \, u_r$ par

\newpage

%% ===== Page 97 =====

% typescript page 95 — n° 262

$w(x)$. Le problème de trouver $v$ ou $w$ étant des problèmes duaux, on se
bornera à établir l'existence de $w$. Utilisant th.1, a), on est ramené
à prouver que l'application $x \to \mathrm{Im}\, u(x)$ de $X$ dans la grassmanienne
$Gr(r,W)$ est régulière (puisque $G\ell(W)$ est fibré localement trivial sur
$Gr(r,W)$). En vertu du th.1, c), la $\mathcal{V}$-structure de $Gr(r,W)$ est induite
par celle de $P(\overset{r}{\wedge} W)$, pour l'immersion naturelle de $Gr(r,W)$ dans $P(\overset{r}{\wedge} W)$

Soit $x_o \in X$, soient $a_i$ ($1 \leqslant i \leqslant r$) des vecteurs tels que les $u(x_o).a_i$ sont
linéairement indépendants, on en conclut facilement qu'il existe un
voisinage ouvert $U$ de $x_o$ tel que $x \in U$ implique que les $u(x).a_i$ soient
linéairement indépendants. L'application $U \to P(\overset{r}{\wedge} W)$ déduite de $u$ est
alors celle déduite de l'application \quad $x \to u(x).a_1 \wedge \dots \wedge u(x).a_r$, et
est donc bien régulière en vertu de th.1) b).

Démontrons alors que a)$\Rightarrow$c), en supposant (ce qui est loisible
grâce au corollaire 1) que $u$ est donné par une application de la forme
$x \to v(x)u_r w(x)^{-1}$ de $X$ dans $L(V,W)$. Faisant dans $X \times V$ le changement de
coordonnées défini par l'application, et le changement analogue dans
$X \times W$, on trouve que $u$ est défini par une application linéaire constante
de $V$ dans $W$, ce qui implique évidemment c).

-Soient $E_1$ et $E_2$ deux $\mathcal{V}$-fibrés vectoriels sur $X$. On appelle \emph{$\mathcal{V}$-
fibré vectoriel extension de $E_2$ par $E_1$} un $\mathcal{V}$-fibré vectoriel $E$, muni
de la structure supplémentaire définie par une suite exacte d'homomor-
phismes de $\mathcal{V}$-fibrés vectoriels :
\begin{center}
\begin{tikzcd}
0 \arrow[r] & E_1 \arrow[r, "i"] & E \arrow[r, "j"] & E_2 \arrow[r] & 0
\end{tikzcd}
\end{center}
("exacte" signifiant : exacte sur chaque fibre). Utilisant la proposi-
tion 18, on voit que ceci implique que $E_1$ est isomorphe à un sous-$\mathcal{V}$-
fibré vectoriel de $E$ et $E_2$ est isomorphe au $\mathcal{V}$-fibré vectoriel quotient

\newpage

%% ===== Page 98 =====

% typescript page 96 — n° 262

$E/E_1$. Il en résulte aussitôt, en vertu des définitions, que sur tout
ouvert $U$ assez petit de $X$, $E$ est isomorphe au produit $E_1 \times E_2$, $i$ et $j$
étant transformées par cet isomorphisme en les applications d'injections
et projections naturelles $E_1 \rightarrow E_1 \times E_2 \rightarrow E_2$. Considérons alors les exten-
sions de $E_2$ par $E_1$ sur des ouverts non précisés de $X$ comme une catégorie
topologique sur $X$, en prenant comme morphismes les isomorphismes (le fait
qu'on obtient bien une catégorie topologique et non seulement prétopolo-
gique résulte de par.1, prop.13). Comme les objets de cette catégorie
sont tous localement isomorphes à $E_1 \times E_2$, il s'ensuit (Top. Elem. chap I)
que l'ensemble des classes d'extensions de $E_2$ par $E_1$ sur $X$ s'identifie
à $H^1(X,\underline{G})$ où $\underline{G}$ désigne le faisceau des germes d'automorphismes de $E_1 \times E_2$.
Or un automorphisme de $E_1 \times E_2$ sur un ouvert $U$ est un automorphisme de sa
structure de $\mathcal{V}$-fibré vectoriel, invariant le sous-fibré $E_1$ et induisant
l'identité sur le sous-fibré $E_1$ et dans le fibré quotient $E/E_1 = E_2$, et
est donc donné par une matrice de morphismes
$$
\begin{pmatrix}
1_{E_1} & u \\
0 & 1_{E_2}
\end{pmatrix}
$$
où $u$ est un homomorphisme arbitraire de $E_2|U$ dans $E_1|U$. D'ailleurs,
comme bien connu, la multiplication des matrices de cette forme corres-
pond à l'addition des termes $u$ correspondants. Ainsi, $\underline{G}$ s'identifie donc
au faisceau de groupes abéliens des germes d'homomorphismes de $E_2$ dans $E_1$
i.e. au faisceau $\underline{O}_X(E^\prime_2 \otimes E_1)$ des germes de sections régulières du fibré
vectoriel $E^\prime_2 \otimes E_1$ (où $E^\prime_2$ désigne le fibré dual de $E_2$). On a ainsi prouvé

\underline{Théorème} 2 - Soient $E_1$ et $E_2$ deux $\mathcal{V}$-fibrés vectoriels sur $X$. L'ensembl
des classes de $\mathcal{V}$-fibrés vectoriels extension de $E_2$ par $E_1$ s'identifie à

\newpage

%% ===== Page 99 =====

% typescript page 97 — n° 262

$H^1(X,\underline{O}_X(E'_2 \otimes E_1))$ (où $E'_2$ désigne le $\mathcal{V}$-fibré vectoriel dual de $E_2$).

Cela montre donc que l'ensemble des classes d'extensions considé-
rées est muni de façon naturelle d'une structure d'espace vectoriel sur
$K$, dont le lecteur trouvera l'interprétation géométrique, en s'inspirant
de constructions faites en Algèbre Homologique. - L'élément de
$H^1(X,\underline{O}_X(E'_2 \otimes E_1))$ qui correspond à une extension $E$ de $E_2$ par $E_1$ est appe-
lé \emph{classe fondamentale de l'extension E}. Son annulation est nécessaire
et suffisante pour que $E$ soit isomorphe à l'extension (dite "triviale")
$E_1 \times E_2$.

Il importe aussi de donner la définition de la classe fondamentale
de l'extension $E$ par un cobord : transformant la suite exacte
$0 \rightarrow E_1 \rightarrow E \rightarrow E_2 \rightarrow 0$ par $E'_2 \otimes$ , on trouve la suite exacte
$$0 \rightarrow E'_2 \otimes E_1 \rightarrow E'_2 \otimes E \rightarrow E'_2 \otimes E_2 \rightarrow 0$$
d'où une suite exacte correspondante sur les faisceaux de germes de
sections régulières, d'où un homomorphisme cobord
$$\partial : H^0(X,\underline{O}_X(E'_2 \otimes E_2)) \rightarrow H^1(X,\underline{O}_X(E'_2 \otimes E_1))$$

Le premier terme s'identifie au groupe des endomorphismes de $E_2$, et on
montre facilement, en montant aux définitions, que $\partial(1_{E_2})$ n'est autre
que la classe fondamentale de l'extension $E$, (à moins qu'il n'y ait le
signe -1).

Si $E$ est une extension de $E_2$ par $E_1$, alors $E'$ est une extension de
$E'_1$ par $E'_2$, sa classe fondamentale est alors un élément de $H^1(X,\underline{O}_X(E_1 \otimes E'_2))$.
Identifiant abusivement $E_1 \otimes E'_2$ et $E'_2 \otimes E_1$, on constate alors que $c(E') =$
$- c(E)$ \quad ($c$ désignant la classe fondamentale). Le signe - provient du
fait que la contragrédiente de la matrice
$$
\begin{pmatrix}
1_{E_1} & u \\
0 & 1_{E_2}
\end{pmatrix}
$$

\newpage

%% ===== Page 100 =====

% typescript page 98 - n° 262
\begin{center}
- 98 - \hfill n° 262
\end{center}

s'obtient en remplaçant $u$ par $-u$ et interchangeant $\underline{1}$ et $\underline{2}$.

Si $E$ est une extension de $E_2$ par $E_1$, et si $E_1$ et $E_2$ ont les fibres-
types $K^n$ resp. $K^m$, alors $E$ est aussi muni d'une structure sous-jacente de
$\mathcal{V}$-fibrés en drapeaux de type $(n, n+m)$ (th. 1, cor. 3). Mais on fera atten-
tion que (même si $E_1$ et $E_2$ sont constants et à fibre-type $K$) deux exten-
sions $E$ et $F$ \emph{non isomorphes de $E_2$ par $E_1$ peuvent avoir des structures
sous-jacentes de fibrés en drapeaux isomorphes.} De façon précise, le grou-
pe $Aut(E_1) \times Aut(E_2)$ opère dans $H^1(X, \underline{O}_X(E_2^\prime \otimes E_1))$, et les éléments d'une
orbite sont exactement les classes d'extensions donnant une même classe
de fibrés en drapeaux. (N.B. - Cette remarque permet en même temps de
classifier les fibrés en drapeaux de type $(n, n+m)$ dont les "facteurs" $E_1$
et $E_2$ sont chacun dans une classe donnée de fibrés vectoriels).

\emph{Prime au lecteur} - Si $X$ est paracompact et $\underline{O}_X$ fin, alors toute extension
de $\mathcal{V}$-fibrés vectoriels est triviale. (En effet, pour tout $\mathcal{V}$-fibré vec-
toriel $L$, $\underline{O}_X(L)$ est fin, donc $H^i(X, \underline{O}_X(L)) = 0$ pour $i > 0$).

\emph{Exercices.} a) Soient $E_1$ et $E_2$ deux $\mathcal{V}$-fibrés vectoriels sur $X$. Mon-
trer que la catégorie locale sur $X$ des $\mathcal{V}$-fibrés vectoriels extension de
$E_2$ par $E_1$ est équivalente à la catégorie locale des $\underline{O}_X$-Modules extension
de $\underline{O}_X(E_2)$ par $\underline{O}_X(E_1)$. b) Soit $X$ un espace topologique muni d'un faisceau
d'anneaux $\underline{O}$, soient $M_1$ et $M_2$ deux $\underline{O}$-Modules. Montrer que l'ensemble
des classes de $\underline{O}$-Modules extensions localement triviales de $M_2$ par $M_1$
s'identifie à $H^1(X, \underline{Hom}_{\underline{O}}(M_2, M_1))$. c) déduire de a) et b) une nouvelle
démonstration du th. 2.

\begin{center}
-:-:-:-:-:-:-:-:-:-:-
\end{center}
